\part{Hodge Theory III}\label{hodge-theory-iii}\thispagestyle{fancy}\par{}P. Deligne.
  "Théorie de Hodge III".
  \emph{Pub. Math. de l'IHÉS} \textbf{44} (1974) pp. 5–77.
  publications.ias.edu/node/361\chapter{Introduction}\thispagestyle{fancy}\label{hodge-theory-iii-introduction}\par{}This article follows \hyperref[hodge-theory-i]{Hodge Theory I} and \hyperref[hodge-theory-ii]{Hodge Theory II};
  the numbering of the sections here follows that of \hyperref[hodge-theory-ii]{Volume~\ref*{hodge-theory-ii}} (which contains \cref{hodge-theory-ii-1} to \cref{hodge-theory-ii-4}).\par{}In \hyperref[hodge-theory-ii]{Volume~\ref*{hodge-theory-ii}}, we gave an introduction to the yoga that underlies \hyperref[hodge-theory-ii]{Volume~\ref*{hodge-theory-ii}} and the present article \hyperref[hodge-theory-iii]{Volume~\ref*{hodge-theory-iii}}.
  However, \hyperref[hodge-theory-ii]{Volume~\ref*{hodge-theory-ii}} and \hyperref[hodge-theory-iii]{Volume~\ref*{hodge-theory-iii}} are logically independent of \hyperref[hodge-theory-i]{Volume~\ref*{hodge-theory-i}};
  they also do not contain all the results that were stated in \hyperref[hodge-theory-i]{Volume~\ref*{hodge-theory-i}}.\par{}In \hyperref[hodge-theory-ii]{Volume~\ref*{hodge-theory-ii}}, we introduced the Hodge theory of non-singular (not necessarily complete) algebraic varieties.
  Here we treat the case of arbitrary singularities.
  From \cref{hodge-theory-iii-7} onwards, we will make essential use of results from \cref{hodge-theory-ii-1} to \cref{hodge-theory-ii-3}.\par{}In the case of a \emph{complete} singular algebraic variety \(X\), the fundamental idea is to use the resolution of singularities to replace \(X\) by a simplicial system of smooth projective schemes
  
  \begin {tikzcd}
    X_\bullet  = \Big (\ldots \quad 
    X_2
      \ar [r]
      \ar [r,shift left=4]
      \ar [r,shift right=4]
    & X_1
      \ar [r,shift left=2]
      \ar [r,shift right=2]
      \ar [l,shift right=2]
      \ar [l,shift left=2]
    & X_0\Big )
      \ar [l]
  \end {tikzcd}

  (or, as we will say, by a smooth projective simplicial scheme).
  The results of \cref{hodge-theory-iii-5} and \cref{hodge-theory-iii-6} allow us to find such simplicial schemes \(X_\bullet \) that have, in a suitable sense, the same cohomology as \(X\).
  We can then "express", via a spectral sequence, the cohomology of \(X\) in terms of the cohomology of the \(X_n\).
  Following the general principles of \hyperref[hodge-theory-i]{Volume~\ref*{hodge-theory-i}}, we can then, on each \(X_n\), use classical Hodge theory, and obtain an induced mixed Hodge structure on the cohomology of \(X\).\par{}In the case of an arbitrary algebraic variety \(X\), we "replace" \(X\) by a smooth simplicial scheme \(X_\bullet \) with compactification given by a smooth projective simplicial scheme \(\overline {X}_\bullet \).
  We can further arrange things so that \(\overline {X}_n\setminus  X_n\) is a normal crossing divisor, given by a union of smooth divisors \(D_{n,i}\).
  We thus "express", via a spectral sequence, the cohomology of \(X\) in terms of the cohomology of the \(p\)-fold intersections of the \(D_{n,i}\) (for \(n,p\geqslant 0\)), and we obtain an induced mixed Hodge structure on \(H^\bullet (X)\).\par{}\cref{hodge-theory-iii-5} and \cref{hodge-theory-iii-6} constitute a summary (without proofs) of the theory of "cohomological descent".
  This theory, in the more general framework of a topos, is laid out in \hyperref[SD]{[SD]}.\par{}In \cref{hodge-theory-iii-7.1}, we restate certain results obtained previously, but in the language of filtered derived categories;
  \cref{hodge-theory-iii-7.2} contains a simpler proof of the two-filtrations lemma than that of \cref{hodge-theory-ii-1.3}.\par{}The unappealing \cref{hodge-theory-iii-8.1} is the heart of this work.
  The essential result is \hyperref[hodge-theory-iii-8.1.15]{Theorem~\ref*{hodge-theory-iii-8.1.15}}.
  
  In its proof, to be able to apply the two-filtrations lemma, we use the fact that every morphism of mixed Hodge structures is \emph{strictly} compatible with the filtrations \(W\) and \(W\) (\hyperref[hodge-theory-ii-2.3.5]{Theorem~\ref*{hodge-theory-ii-2.3.5}}).
  In \cref{hodge-theory-iii-8.2}, we define the mixed Hodge structure of \(\operatorname {H}^n(X,\mathbb {Z})\) (for \(X\) a separated scheme).
  We show that the Hodge numbers \(h^{pq}\) of \(\operatorname {H}^n(X,\mathbb {Z})\) can only be non-zero for \((p,q)\in [0,n]\times [0,n]\).
  For \(n>\dim  X\), we have a more precise result (\hyperref[hodge-theory-iii-8.2.4]{Theorem~\ref*{hodge-theory-iii-8.2.4}}).
  If \(X\) is smooth (resp. complete), they can only be non-zero if, further, \(p+q\geqslant  n\) (resp. \(p+q\leqslant  n\)).
  If \(X\) is a complete subscheme of \(Z\), with complement \(U\), with \(Z\) complete and smooth of dimension \(n\), then we can, by N. Katz, interpret this "duality" between the complete case and the smooth case as coming from the "Alexander duality"
  \[     \to  \ldots      \to  \operatorname {H}^i(Z,\mathbb {Q})     \to  \operatorname {H}^i(X,\mathbb {Q})     \to  (\operatorname {H}^{2n-i-1}(U,\mathbb {Q}))^*     \to  \operatorname {H}^{i+1}(Z,\mathbb {Q})     \to \ldots .   \]\par{}From the functorial properties of the theory thus constructed, we obtain the useful corollaries \hyperref[hodge-theory-iii-8.2.5]{Proposition~\ref*{hodge-theory-iii-8.2.5}} to \hyperref[hodge-theory-iii-8.2.8]{Corollary~\ref*{hodge-theory-iii-8.2.8}}, which we state in terms independent of any Hodge theory.\par{}In \cref{hodge-theory-iii-8.3}, we endow the cohomology of any simplicial scheme \(X_\bullet \) with a mixed Hodge structure.
  This generality is not illusory.

  
 \begin{enumerate}
    
 \item[(a)]
      We can interpret relative cohomology spaces as being the cohomology of suitable simplicial schemes.
    

    
 \item[(b)]
      Let \(B_G\) be the classifying space of the underlying Lie group of an algebraic group \(G\).
      We can interpret the cohomology of \(B_G\) as being the cohomology of a suitable simplicial scheme.
    

  \end{enumerate}\par{}In \cref{hodge-theory-iii-9.1}, after having calculated the mixed Hodge structure on the cohomology of \(B_G\) (for \(G\) linear), we thus obtain that of \(G\).
  As a corollary, we find that, if a linear group \(G\) acts on a non-singular complete variety \(X\), then the map
  \[     \operatorname {H}^\bullet (X,\mathbb {Q})     \to  \operatorname {H}^\bullet (X,\mathbb {Q})\otimes \operatorname {H}^\bullet (G,\mathbb {Q})   \]
  factors through \(\operatorname {H}^\bullet (X,\mathbb {Q})\otimes \operatorname {H}^0(G,\mathbb {Q})\subset \operatorname {H}^\bullet (X,\mathbb {Q})\otimes \operatorname {H}^\bullet (G,\mathbb {Q})\).\par{}In \cref{hodge-theory-iii-10}, we interpret in terms of abelian schemes the mixed Hodge structures of pure degree \(\{(-1,-1),(-1,0),(0,-1),(0,0)\}\) (see \cref{hodge-theory-iii-0.5}), and we consider in detail the \(\operatorname {H}^1\) of curves.\par{}The theory developed up until here is an absolute theory (we do not study the functors \(\mathrm {R} f_*\)), and deals only with constant coefficients.
  I conjecture that, if \(\mathscr {H}\) is a variation of polarisable Hodge structures (in the sense of Griffiths \hyperref[G1970a]{[G1970a]}) on a scheme \(X\), then the cohomology of \(X\) with coefficients in the local system \(\mathscr {H}\) is endowed with a natural mixed Hodge structure.
  I can only prove this when \(X\) is complete.\chapter{Terminology and notation}\thispagestyle{fancy}\label{hodge-theory-iii-terminology-and-notation}\label{hodge-theory-iii-0.1}\par{}Let \(u\colon  X\to  Y\) be a continuous map between topological spaces.
    We say that \(u\) is \emph{proper} if it is proper in the sense of Bourbaki (i.e. universally closed) and furthermore separated (i.e. the diagonal of \(X\times _Y X\) is closed).\label{hodge-theory-iii-0.2}\par{}Following Gabriel and Zisman, we say \emph{simplicial} where we would have previously said semi-simplicial.\label{hodge-theory-iii-0.3}\par{}We denote by \(A\) a Noetherian subring of \(\mathbb {R}\) such that \(A\otimes \mathbb {Q}\) is a field.
    The useful cases are \(A=\mathbb {Z}\), \(\mathbb {Q}\), or \(\mathbb {R}\).\label{hodge-theory-iii-0.4}\par{}A \emph{mixed Hodge \(A\)-structure} consists of an \(A\)-module \(H_A\) of finite type, a finite increasing filtration \(W\) on the \(A\otimes \mathbb {Q}\)-vector space \(H_{A\otimes \mathbb {Q}}=H_A\otimes \mathbb {Q}\), and a finite decreasing filtration \(F\) on the \(\mathbb {C}\)-vector space \(H_\mathbb {C}=H_A\otimes _A\mathbb {C}\).
    We demand that the \((\operatorname {Gr}_n^W(H_{A\otimes \mathbb {Q}}),\operatorname {Gr}_n^W(F))\) be Hodge \(A\otimes \mathbb {Q}\)-structures.
    For \(A=\mathbb {Z}\) (resp. \(A=\mathbb {Q}\)), we recover \hyperref[hodge-theory-ii-2.3.1]{Definition~\ref*{hodge-theory-ii-2.3.1}} (resp. \cref{hodge-theory-ii-2.3.8});
    the results of \cref{hodge-theory-ii-2.3} carry over as they are.\label{hodge-theory-iii-0.5}\par{}Let \(\mathscr {E}\) be a subset of \(\mathbb {Z}\times \mathbb {Z}\).
    We say that a mixed Hodge structure \(H\) is of \emph{degree \(\mathscr {E}\)} if the Hodge numbers \(h^{pq}\) are zero for \((p,q)\not \in \mathscr {E}\).\label{hodge-theory-iii-0.6}\par{}From now on, we denote by \(\Omega _X^\bullet (\log  D)\) what we previously denoted by \(\Omega _X^\bullet \langle  D\rangle \) (\cref{hodge-theory-ii-3.1.2}).\label{hodge-theory-iii-0.7}\par{}Unless explicitly stated otherwise, \emph{scheme} means "scheme of finite type over \(\mathbb {C}\)", and a \emph{sheaf} on a scheme \(X\) is a sheaf on the underlying topological space of \(X_\mathrm {an}\).\chapter{Cohomological descent}\thispagestyle{fancy}\label{hodge-theory-iii-5}\section{Simplicial topological spaces}\label{hodge-theory-iii-5.1}\par{}This section starts with some reminders for which we can refer to the homotopical seminar of Strasbourg, 1963/64.\label{hodge-theory-iii-5.1.1}\par{}We will us the following notation, where \(n,k\geqslant -1\) are integers.

    \begin{itemize}\item{}\(\mathscr {A}^\circ \)
        = the opposite category of a category \(\mathscr {A}\)
      \item{}\(\underline {\operatorname {Hom}}(\mathscr {A},\mathscr {B})\)
        = the category of functors from \(\mathscr {A}\) to \(\mathscr {B}\).
      \item{}\(\Delta _n\)
        = the finite totally ordered set \([0,n]\).
      \item{}\(\delta _i\colon \Delta _n\to \Delta _{n+1}\)
        = the increasing injection such that \(i\not \in \delta _i(\Delta _n)\) (for \(0\leqslant  i\leqslant  n+1\)).
      \item{}\(s_i\colon \Delta _{n+1}\to \Delta _n\)
        = the increasing surjection such that \(s_i(i)=s_i(i+1)\) (for \(0\leqslant  i\leqslant  n\)).
      \item{}\(\varepsilon \colon \Delta _{-1}\to \Delta _n\)
        = the unique map from \(\Delta _{-1}\) to \(\Delta _n\).
      \item{}\((\Delta ^+)\)
        = the category whose objects are the \(\Delta _n\) (for \(n\geqslant -1\)) and whose morphisms are the increasing maps between the \(\Delta _n\).
      \item{}\((\Delta )\)
        = the full subcategory of \((\Delta ^+)\) whose objects are the \(\Delta _n\) (for \(n\geqslant 0\)).
      \item{}\((\Delta ^+)_k\)
        = the full subcategory of \((\Delta ^+)\) whose objects are the \(\Delta _n\) (for \(k\geqslant  n\)).
      \item{}\((\Delta )_k\)
        = the full subcategory of \((\Delta ^+)\) whose objects are the \(\Delta _n\) (for \(k\geqslant  n\geqslant 0\)).\end{itemize}\par{}For any category \(\mathscr {C}\), we define a \emph{simplicial object} (resp. \emph{\(k\)-truncated simplicial object}) of \(\mathscr {C}\) to be an object of \(\underline {\operatorname {Hom}}((\Delta )^0,\mathscr {C})\) (resp. of \(\underline {\operatorname {Hom}}((\Delta )^0_k,\mathscr {C})\)).
    Similarly, a \emph{cosimplicial object} (resp. \emph{\(k\)-truncated cosimplicial object}) is an object of \(\underline {\operatorname {Hom}}((\Delta ),\mathscr {C})\) (resp. of \(\underline {\operatorname {Hom}}((\Delta )_k,\mathscr {C})\)).\par{}For \(k\geqslant -1\), the \emph{\(k\)-skeleton} functor is the restriction functor
    \[       \operatorname {sq}_k\colon        \underline {\operatorname {Hom}}((\Delta )^0,\mathscr {C})       \to  \underline {\operatorname {Hom}}((\Delta )_k^0,\mathscr {C})     \]
    and the \emph{\(k\)-coskeleton} functor
    \[       \operatorname {cosq}_k\colon        \underline {\operatorname {Hom}}((\Delta )_k^0,\mathscr {C})       \to  \underline {\operatorname {Hom}}((\Delta )^0,\mathscr {C}).     \]
    is the right adjoint to \(\operatorname {sq}_k\).
    Let \(Y\) be a simplicial object of \(\mathscr {C}\).
    We also call the functor
    \[       \operatorname {sq}_k\colon        \underline {\operatorname {Hom}}((\Delta )^0,\mathscr {C})/Y       \to  \underline {\operatorname {Hom}}((\Delta )_k^0,\mathscr {C})/\operatorname {sq}_k(Y)     \]
    the \emph{skeleton}, and its right adjoint
    \[       \operatorname {cosq}_k^Y\colon        \underline {\operatorname {Hom}}((\Delta )_k^0,\mathscr {C})/\operatorname {sq}_k(Y)       \to  \underline {\operatorname {Hom}}((\Delta )^0,\mathscr {C})/Y     \]
    is called the \emph{coskeleton with respect to \(Y\)}.\par{}The coskeleton functors exist if finite projective limits exist in \(\mathscr {C}\);
    the relative coskeleton functors exist if fibred products do;
    we have
    \[       \operatorname {cosq}_k^Y(X)       = \operatorname {cosq}_k(X) \times _{\operatorname {cosq}_k\operatorname {sq}_k(Y)}Y.     \]\label{hodge-theory-iii-5.1.2}\par{}If \(X_\bullet \colon (\Delta )^0\to \mathscr {C}\) is a simplicial object of \(\mathscr {C}\), then we set

    \begin{itemize}\item{}\(X_n = X_\bullet (\Delta _n)\)
      \item{}\(\delta _i = X_\bullet (\delta _i\colon \Delta _n\to \Delta _{n+1})\colon  X_{n+1}\to  X_n\)
      \item{}\(s_i = X_\bullet (s_i\colon \Delta _{n+1}\to \Delta _n)\colon  X_n\to  X_{n+1}\).\end{itemize}

    
  \begin {tikzcd}
      X_\bullet  = \Big (\ldots \quad 
      X_2
        \ar [r,"\delta _1" description]
        \ar [r,shift left=6,"\delta _0" description]
        \ar [r,shift right=6,"\delta _2" description]
      & X_1
        \ar [r,shift left=3,"\delta _0" description]
        \ar [r,shift right=3,"\delta _1" description]
        \ar [l,shift right=3,"s_0" description]
        \ar [l,shift left=3,"s_1" description]
      & X_0\Big )
        \ar [l,"s_0" description]
    \end {tikzcd}\label{hodge-theory-iii-5.1.3}\par{}Let \(S\in \operatorname {Ob}\mathscr {C}\).
    The \emph{constant simplicial object} \(S_\bullet \) is the simplicial object for which \(S_n=S\) and \(\delta _i=s_i=\mathrm {Id}_S\).
    A simplicial (resp. \(k\)-truncated simplicial) object of \(\mathscr {C}\) \emph{augmented over \(S\)} is a morphism \(a\colon  X_\bullet \to  S_\bullet \) (resp. \(a\colon  X_\bullet \to \operatorname {sq}_k(S_\bullet )\)).
    We identify such a simplicial object (resp. \(k\)-truncated simplicial object) augmented over \(S\) with the object \(X_\bullet ^+\) of \(\underline {\operatorname {Hom}}((\Delta ^+)^0,\mathscr {C})\) (resp. of \(\underline {\operatorname {Hom}}((\Delta ^+)_k^0,\mathscr {C})\)) such that \(X_\bullet ^+(\Delta _{-1})=S\).
    We have that \(a_n=X_\bullet ^+(\varepsilon \colon \Delta _{-1}\to \Delta _n)\).
    These objects will be denoted by notation of the form \(a\colon  X_\bullet \to  S\).
    The relative coskeletons will be mostly used in this setting, and denoted \(\operatorname {cosq}_k^S\) or simply \(\operatorname {cosq}_k\).\label{hodge-theory-iii-5.1.4}\par{}The coskeleton of the augmented \(0\)-truncated simplicial object \(a_0\colon  X\to  S\) is the simplicial object augmented over \(S\) whose components are the cartesian powers of \(X\) in \(\mathscr {C}/S\), i.e.
    \[       \operatorname {cosq}_0(X\to  S)       = \big ( ((X/S)^{\Delta _n})_{n\geqslant 0} \to  S \big ).     \]\label{hodge-theory-iii-5.1.5}\par{}Let \(u\colon  X\to  Y\) be a continuous map between topological spaces, and \(F\) a sheaf on \(X\) and \(G\) a sheaf on \(Y\).
    The set \(\operatorname {Hom}_u(G,F)\) of \emph{\(u\)-morphisms} from \(G\) to \(F\) is the set \(\operatorname {Hom}(u^*G,F)\cong \operatorname {Hom}(G,u_*F)\).\par{}The data of a \(u\)-morphism \(f\) from \(G\) to \(F\) consists of the data of, for each pair of opens \((U\subset  X,V\subset  Y)\) such that \(u(U)\subset  V\), a map \(f_{UV}\colon  G(V)\to  F(U)\);
    these maps must satisfy the condition

    \begin{condition}{*}\setnumber{*}\label{hodge-theory-iii-5.1.5.star}\skipsinglepar\par{}For \(U'\subset  U\subset  X\) and \(V'\subset  V\subset  Y\) such that \(u(U)\subset  V\) and \(u(U')\subset  V'\), the diagram
        
  \begin {tikzcd}
          G(V)
            \ar [r]
            \ar [d]
          & G(V')
            \ar [d]
        \\F(U)
            \ar [r]
          & F(U')
        \end {tikzcd}

        commutes.\end{condition}\label{hodge-theory-iii-5.1.6}\par{}A \emph{simplicial topological space} is a simplicial object of the category whose objects are topological spaces and whose morphisms are continuous maps.\par{}A \emph{sheaf} \(F^\bullet \) on a simplicial topological space \(X_\bullet \) consists of

    
 \begin{enumerate}
      
 \item[(a)]
        a family of sheaves \(F^n\) on the \(X_n\);
      

      
 \item[(b)]
        for all \(f\colon \Delta _n\to \Delta _m\), an \(X_\bullet (f)\)-morphism \(F^\bullet (f)\) from \(F^n\) to \(F^m\).
      

    \end{enumerate}\par{}We further require that \(F^\bullet (f\circ  g)=F^\bullet (f)\circ  F^\bullet (g)\).\par{}A morphism \(u\) from \(F^\bullet \) to \(G^\bullet \) is a family of morphisms \(u^n\colon  F^n\to  G^n\) such that, for all \(f\in \operatorname {Hom}_{(\Delta )}(\Delta _n,\Delta _m)\), we have that \(u^m\circ  F^\bullet (f)=G^\bullet (f)\circ  u^n\).\label{hodge-theory-iii-5.1.7}\par{}For an open \(U\) of \(X_n\), let \(F^\bullet (U)=F^n(U)\).
    For \(f\colon \Delta _n\to \Delta _m\), \(U\subset  X_n\), and \(V\subset  X_m\) such that \(X_\bullet (f)(V)\subset (U)\), let \(F^\bullet (f,V,U)\colon  F^\bullet (U)\to  F^\bullet (V)\) be the map induced by \(F^\bullet (f)\).
    We immediately see that the system of sets \(F^\bullet (U)\) (indexed by \(n\geqslant 0\) and \(U\subset  X_n\)) and of maps \(F^\bullet (f,V,U)\) uniquely determines \(F^\bullet \) (cf. \cref{hodge-theory-iii-5.1.5}).
    For such a system to come from a sheaf on \(X_\bullet \), it is necessary and sufficient that:

    
 \begin{enumerate}
      
 \item[(a)]
        \(F^\bullet (fg,U,W)=F^\bullet (f,U,V)F^\bullet (g,V,W)\) whenever this makes sense;
      

      
 \item[(b)]
        for any \(n\), the \(F^\bullet (U)\) for \(U\subset  X_n\) form a sheaf on \(X_n\).
      

    \end{enumerate}\label{hodge-theory-iii-5.1.8}\par{}This shows that sheaves on \(X_\bullet \) can be interpreted as sheaves on a suitable site, and, in particular, form a topos \((X_\bullet )^\sim \).
    We freely use, for sheaves on \(X_\bullet \), the modern terminology for sheaves on a site.
    With this in mind, a sheaf \(F^\bullet \) of abelian groups (resp. of rings, ...) is a system of sheaves \(F^n\) of abelian groups (resp. of rings, ...) on the \(X_n\) along with the morphisms \(F^\bullet (f)\).\begin{examples}{5.1.9}\setnumber{5.1.9}\label{hodge-theory-iii-5.1.9}\skipsinglepar\begin{enumerate}
    
 \item[(I)]
      Let \(X_\bullet \) be a simplicial analytic space.
      The structure sheaves \(\mathcal {O}_{X_n}\) form a sheaf of rings on \(X_\bullet \).
    


    
 \item[(II)]
      Let \(X_\bullet \) be a simplicial analytic space augmented over \(S\).
      The sheaves \(\Omega _{X_n/S}^1\) form a sheaf of \(\mathcal {O}\)-modules on \(X_\bullet \).
      Its \(i\)-th exterior power is the sheaf of \(\mathcal {O}\)-modules \((\Omega _{X_n/S}^i)_{n\geqslant 0}\) denoted \(\Omega _{X_\bullet {/}S}^i\).
      
      The de Rham complexes \((\Omega _{X_n/S}^\bullet )_{n\geqslant 0}\) form a complex of sheaves on \(X_\bullet \).
    


    
 \item[(III)]
      Let \(F^\bullet \) be a sheaf of abelian groups on \(X_\bullet \).
      The canonical flasque Godement resolutions \(\mathscr {C}^\bullet (F^n)\) form a complex of sheaves on \(X_\bullet \) that is a resolution of \(F^\bullet \).
    


    
 \item[(IV)]
      Let \(S\) be a topological space.
      Sheaves \(F^\bullet \) on the constant simplicial space \(S_\bullet \) can be identified with cosimplicial sheaves on \(S^\bullet \).
      In particular, an abelian sheaf \(F^\bullet \) on \(S_\bullet \) defines a chain complex \((F^n,d=\sum _i(-1)^i\delta _i)\).
      A complex \(K\) of abelian sheaves on \(S_\bullet \) defines a double complex \((K^{n,m})\) that we again denote by \(K\) (here \(m\) is the cosimplicial degree) and whose associated simple complex \(\mathrm {s}{K}\) is

      \begin{eqn}{5.1.9.1}\setnumber{5.1.9.1}\label{hodge-theory-iii-5.1.9.1}\skipsinglepar\[           (\mathrm {s}{K})^n           = \bigoplus _{p+q=n} K^{pq}         \tag{5.1.9.1}         \]\end{eqn}

      with differential given by

      \begin{eqn}{5.1.9.2}\setnumber{5.1.9.2}\label{hodge-theory-iii-5.1.9.2}\skipsinglepar\[           \mathrm {d}(x^{pq})           = \mathrm {d}_K(x^{pq}) + (-1)^p\sum _i(-1)^i\delta _i x^{pq}.         \tag{5.1.9.2}         \]\end{eqn}

      We denote by \(L\) the second filtration of \(\mathrm {s}{K}\), i.e.

      \begin{eqn}{5.1.9.3}\setnumber{5.1.9.3}\label{hodge-theory-iii-5.1.9.3}\skipsinglepar\[           L^r(\mathrm {s}{K})           = \bigoplus _{q\geqslant  r} K^{pq}.         \tag{5.1.9.3}         \]\end{eqn}
    

  \end{enumerate}\end{examples}\label{hodge-theory-iii-5.1.10}\par{}Let \(u\colon  X_\bullet \to  Y_\bullet \) be a morphism of simplicial topological spaces, with components \(u_n\colon  X_n\to  Y_n\).
    If \(G\) (resp. \(F\)) is a sheaf on \(Y_\bullet \) (resp. on \(X_\bullet \)), then \((u_n^*G^n)_{n\geqslant 0}\) (resp. \(({u_n}_*F^n)_{n\geqslant 0}\)) is a sheaf on \(X_\bullet \) (resp. on \(Y_\bullet \));
    we denote it by \(u^*G\) (resp. by \(u_*F\)).
    The functors \(u^*\) and \(u_*\) are adjoint to one another;
    they are the inverse image and direct image morphisms of topos morphism \(u\colon (X_\bullet )^\sim \to (Y_\bullet )^\sim \).\label{hodge-theory-iii-5.1.11}\par{}Let \(a\colon  X_\bullet \to  S\) be an augmented simplicial topological space.
    If \(F\) is a sheaf on \(S\), then \(a^*F=(a_n^*F)_{n\geqslant 0}\) "is" a sheaf on \(X_\bullet \).
    The functor \(a^*\) has a right adjoint

    \begin{eqn}{5.1.11.1}\setnumber{5.1.11.1}\label{hodge-theory-iii-5.1.11.1}\skipsinglepar\[         a_*\colon  F^\bullet          \mapsto  \operatorname {Ker}\big (           {a_0}_* F^0           \,\overset {\delta _0^*}{\underset {\delta _1^*}{\rightrightarrows }}\,           {a_1}_* F^1         \big ).       \tag{5.1.11.1}       \]\end{eqn}

    The functors \(a^*\) and \(a_*\) are the inverse image and direct image morphisms of topos morphism \(a\colon (X_\bullet )^\sim \to (S)^\sim \).\label{hodge-theory-iii-5.1.12}\par{}Let \(S_\bullet \) be the constant simplicial space associated to \(S\), and \(a_\bullet \colon  X_\bullet \to  S_\bullet \) the morphism defined by \(a\).
    For an abelian sheaf \(F^\bullet \) on \(X_\bullet \), we can identify \({a_\bullet }_*F^\bullet \) with a cosimplicial sheaf on \(S\) ((IV) of \hyperref[hodge-theory-iii-5.1.9]{Examples~\ref*{hodge-theory-iii-5.1.9}}).
    For a complex \(K\) of abelian sheaves on \(X_\bullet \), if \(K^{pq}\) is the degree \(p\) component of the restriction of \(K\) to \(X_q\), then the components of the simple complex associated to the double complex defined by \({a_\bullet }_*K\) are

    \begin{eqn}{5.1.12.1}\setnumber{5.1.12.1}\label{hodge-theory-iii-5.1.12.1}\skipsinglepar\[         (\mathrm {s}{a_\bullet }_*K)^n         = \bigoplus _{p+q=n} {a_q}_* K^{pq}.       \tag{5.1.12.1}       \]\end{eqn}

    
    The spectral sequence defined by the filtration \(L\) (\hyperref[hodge-theory-iii-5.1.9.3]{Equation~\ref*{hodge-theory-iii-5.1.9.3}}) of \(\mathrm {s}{a_\bullet }_*K\) can be written as

    \begin{eqn}{5.1.12.2}\setnumber{5.1.12.2}\label{hodge-theory-iii-5.1.12.2}\skipsinglepar\[         E_1^{pq}         = \operatorname {H}^q({a_p}_*(K|X^p))         \Rightarrow  \operatorname {H}^{p+q}(\mathrm {s}{a_\bullet }_*K).       \tag{5.1.12.2}       \]\end{eqn}\label{hodge-theory-iii-5.1.13}\par{}Let \(S=P^t\) (the topological space consisting of a single point).
    We see that, for a sheaf \(F^\bullet \) on \(X_\bullet \), the \(\Gamma (X_n,F^n)\) are the components of a cosimplicial set \(\Gamma ^\bullet (X_\bullet ,F^\bullet )\).
    The functor \(\Gamma \) (of global sections) is the functor

    \begin{eqn}{5.1.13.1}\setnumber{5.1.13.1}\label{hodge-theory-iii-5.1.13.1}\skipsinglepar\[         \Gamma \colon  F^\bullet          \mapsto  \operatorname {Ker}\big (           \Gamma (X_0,F^0)           \rightrightarrows  \Gamma (X_1,F^1)         \big ).       \tag{5.1.13.1}       \]\end{eqn}

    If \(K\) is a complex of abelian sheaves, we denote by \(\Gamma ^\bullet (X_\bullet ,K)\) the chain complex whose components are the cosimplicial abelian groups \(\Gamma ^\bullet (X_\bullet ,K^n)\), and by \(\mathrm {s} \Gamma ^\bullet (X_\bullet ,K)\) the associated simple complex.\section{Cohomology of simplicial topological spaces}\label{hodge-theory-iii-5.2}\label{hodge-theory-iii-5.2.1}\par{}To each simplicial topological space \(X_\bullet \) is associated a usual topological space \(|X_\bullet |\), called its geometric realisation (see below).
    In the particular case where the \(X_n\) are discrete, we recover the usual notion of geometric realisation of a simplicial set.\par{}In [\hyperref[S1968]{S1968}], G. Segal defined the cohomology of \(X_\bullet \) with values in an abelian group \(A\) as \(\operatorname {H}^\bullet (|X_\bullet |,A)\).
    Under suitable hypotheses, the filtration of \(|X_\bullet |\) by successive skeletons gives a spectral sequence

    \begin{eqn}{5.2.1.1}\setnumber{5.2.1.1}\label{hodge-theory-iii-5.2.1.1}\skipsinglepar\[         E_1^{pq}         = \operatorname {H}^q(X_p,A)         \Rightarrow  \operatorname {H}^{p+q}(X_\bullet ,A).       \tag{5.2.1.1}       \]\end{eqn}\par{}One of the interests in using this definition is that it also applies, for example, to the definition of the groups \(K(X_\bullet )\).
    We will adopt another definition, which is better adapted to sheaf-theoretic techniques.Geometric realisation\label{hodge-theory-iii-5.2.1-geometric-realisation}\par{}For an integer \(n\geqslant 0\), we denote by \(|\Delta _n|\) the simple in \(\mathbb {R}^{\Delta _n}\) whose vertices are the set of basis vectors.
      We identify \(\Delta _n\) with the set of vertices of \(|\Delta _n|\).
      Every function \(f\colon \Delta _n\to \Delta _m\) extends by linearity to \(|f|\colon |\Delta _n|\to |\Delta _m|\).\par{}Let \(X_\bullet \) be a simplicial topological space.
      Let
      \[         Y=\coprod _{n\geqslant 0}(X_n\times |\Delta _n|).       \]
      Let \(R\) be the finest equivalence relation on \(Y\) for which, for any \(f\colon \Delta _n\to \Delta _m\) in \((\Delta )\), any \(x_m\in  X_m\), and any \(a\in |\Delta _n|\), we have
      \[         (x_m,|f|(a))         \equiv  (f(x_m),a)         \operatorname {mod} R.       \]
      The \emph{geometric realisation} of \(X_\bullet \) is by definition
      \[         |X_\bullet |         = Y/R.       \]\begin{definition}{5.2.2}\setnumber{5.2.2}\label{hodge-theory-iii-5.2.2}\skipsinglepar\par{}Let \(X_\bullet \) be a simplicial topological space.
    The functors \(\operatorname {H}^i(X_\bullet ,F^\bullet )\) of \emph{cohomology with coefficients in the abelian sheaf \(F^\bullet \) on \(X_\bullet \)} are the derived functors of the functor \(\Gamma \) (\hyperref[hodge-theory-iii-5.1.13.1]{Equation~\ref*{hodge-theory-iii-5.1.13.1}}).\end{definition}\par{}This definition is equivalent to the following, sometimes more convenient.\label{hodge-theory-iii-5.2.3}\par{}Let \(F^\bullet \) be an abelian sheaf on \(X_\bullet \).
    We can show (by example, using (III) of \hyperref[hodge-theory-iii-5.1.9]{Examples~\ref*{hodge-theory-iii-5.1.9}}) that \(F^\bullet \) always admits resolutions \(K\) on the right such that
    \[       \operatorname {H}^r(X_q,K^{p,q})       =0       \quad \text {for } p,q\geqslant 0 \text { and } r>0.     \]
    If \(K\) is such a resolution, then we canonically have that

    \begin{eqn}{5.2.3.1}\setnumber{5.2.3.1}\label{hodge-theory-iii-5.2.3.1}\skipsinglepar\[         \operatorname {H}^n(X_\bullet ,F^\bullet )         \simeq  \operatorname {H}^n(\mathrm {s}\Gamma ^\bullet (X_\bullet ,K)).       \tag{5.2.3.1}       \]\end{eqn}

    It is easy to show directly that the right-hand side is independent (up to unique isomorphism) of the choice of \(K\);
    for what follows, we are free to \emph{define} \(\operatorname {H}^\bullet (X_\bullet ,F^\bullet )\) by \hyperref[hodge-theory-iii-5.2.3.1]{Equation~\ref*{hodge-theory-iii-5.2.3.1}}.
    We can further show that the spectral sequence of the complex \(\mathrm {s}\Gamma ^\bullet (X_\bullet ,K)\), filtered by \(L\) (\hyperref[hodge-theory-iii-5.1.9.3]{Equation~\ref*{hodge-theory-iii-5.1.9.3}} and \hyperref[hodge-theory-iii-5.1.12.2]{Equation~\ref*{hodge-theory-iii-5.1.12.2}} for \(S=P^t\))

    \begin{eqn}{5.2.3.2}\setnumber{5.2.3.2}\label{hodge-theory-iii-5.2.3.2}\skipsinglepar\[         E_1^{pq}         = \operatorname {H}^q(X_p,F^p)         \Rightarrow  \operatorname {H}^{p+q}(X_\bullet ,F^\bullet )       \tag{5.2.3.2}       \]\end{eqn}

    is independent (up to unique isomorphism) of the choice of \(K\) (cf. \hyperref[hodge-theory-iii-5.2.1.1]{Equation~\ref*{hodge-theory-iii-5.2.1.1}}).\label{hodge-theory-iii-5.2.4}\par{}We will need to make precise the above construction when passing to derived categories, and to give a relative variant of it.\par{}Let \(a\colon  X_\bullet \to  S\), let \(S_\bullet \) be the constant simplicial space defined by \(S\) with augmentation map \(\varepsilon \colon  S_\bullet \to  S\), and let \(a_\bullet \colon  X_\bullet \to  S_\bullet \) be the map induced by \(a\).
    Then

    \begin{eqn}{5.2.4.1}\setnumber{5.2.4.1}\label{hodge-theory-iii-5.2.4.1}\skipsinglepar\[         a_*         = \varepsilon _* {a_\bullet }_*\colon          (X_\bullet )^\sim          \to  (S)^\sim .       \tag{5.2.4.1}       \]\end{eqn}\par{}The derived version of this equation is

    \begin{eqn}{5.2.4.2}\setnumber{5.2.4.2}\label{hodge-theory-iii-5.2.4.2}\skipsinglepar\[         \mathrm {R} a_*         = \mathrm {R}\varepsilon _* \mathrm {R}{a_\bullet }_*\colon          \mathrm {D}^+(X_\bullet )         \to  \mathrm {D}^+(S)         \to  (S)^\sim .       \tag{5.2.4.2}       \]\end{eqn}

    where \(\mathrm {D}^+(X_\bullet )\) is the bounded-below derived category of the category of abelian sheaves on \(X_\bullet \).
    We will calculate \(\mathrm {R}{a_\bullet }_*\) and \(\mathrm {R}\varepsilon _*\).\label{hodge-theory-iii-5.2.5}\par{}Let \(u\colon  X_\bullet \to  Y_\bullet \) be a morphism of simplicial topological spaces.
    The functor \(\mathrm {R} u_*\colon \mathrm {D}^+(X_\bullet )\to \mathrm {D}^+(Y_\bullet )\) can be calculated "component by component": if \(K\) is a complex of abelian sheaves on \(X_\bullet \), then to calculate \(\mathrm {R} u_*K\) we take a resolution \(K\xrightarrow {\sim } K'\) such that the components \(F\) of \(K'\) satisfy \(\mathrm {R}^i {u_p}_*(F^p)=0\) for \(i>0\), \(p\geqslant 0\) (cf. (III) of \hyperref[hodge-theory-iii-5.1.9]{Examples~\ref*{hodge-theory-iii-5.1.9}});
    then \(\mathrm {R} u_*K\approx  u_*K'\).\label{hodge-theory-iii-5.2.6}\par{}The functor \(\mathrm {s}\colon \operatorname {C}^+(S_\bullet )\to \operatorname {C}^+(S)\) sends acyclic complexes to acyclic complexes.
    It thus trivially derives to
    \[       \mathrm {s}\colon  \mathrm {D}^+(S_\bullet )       \to  \mathrm {D}^+(S).     \]\par{}If \(F\) is an injective sheaf on \(S_\bullet \), then the chain complex \(\mathrm {s} F\) is a resolution of \(\varepsilon _*F\).
    We thus obtain an isomorphism \(\mathrm {R}\varepsilon _*\xrightarrow {\sim }\mathrm {s}\), and \hyperref[hodge-theory-iii-5.2.4.2]{Equation~\ref*{hodge-theory-iii-5.2.4.2}} becomes

    \begin{eqn}{5.2.6.1}\setnumber{5.2.6.1}\label{hodge-theory-iii-5.2.6.1}\skipsinglepar\[         \mathrm {R} a_*         = \mathrm {s}\mathrm {R}{a_\bullet }_*       \tag{5.2.6.1}       \]\end{eqn}\label{hodge-theory-iii-5.2.7}\par{}Combined with \cref{hodge-theory-iii-5.2.5} and specialised to the case where \(S=P^t\), this equation proves \hyperref[hodge-theory-iii-5.2.3.1]{Equation~\ref*{hodge-theory-iii-5.2.3.1}}.
    In concrete terms, this implies that, to calculate \(\mathrm {R} a_*K\), we can proceed in two steps:

    
 \begin{enumerate}
      
 \item[(a)]
        We take a resolution \(K\xrightarrow {\sim } K'\) such that the components \(F\) of \(K'\) satisfy \(\mathrm {R}^i {a_p}_*(F^p)=0\) for \(i>0\).
        The complex \({a_\bullet }_*K'\in \mathrm {D}^+(S_\bullet )\) (the derived category of the category of cosimplicial abelian sheaves on \(S\)) can be identified with \(\mathrm {R}{a_\bullet }_*K\).
      


      
 \item[(b)]
        \(\mathrm {R} a_*K\) is the simple complex \(\mathrm {s}{a_\bullet }_*K'\) associated to the double complex \({a_\bullet }_*K'\).
      

    \end{enumerate}\par{}The spectral sequence \hyperref[hodge-theory-iii-5.2.3.2]{Equation~\ref*{hodge-theory-iii-5.2.3.2}} generalises to a spectral sequence

    \begin{eqn}{5.2.7.1}\setnumber{5.2.7.1}\label{hodge-theory-iii-5.2.7.1}\skipsinglepar\[         E_1^{pq}         = \mathrm {R}^q{a_p}_*(K|X_p)         \Rightarrow  \mathrm {R}^{p+q}a_*K       \tag{5.2.7.1}       \]\end{eqn}

    induced by \hyperref[hodge-theory-iii-5.1.12.2]{Equation~\ref*{hodge-theory-iii-5.1.12.2}}.\section{Cohomological descent}\label{hodge-theory-iii-5.3}\label{hodge-theory-iii-5.3.1}\par{}Let \(a\colon  X_\bullet \to  S\) be an augmented simplicial topological space.
    For every sheaf \(F\) on \(S\), we have a morphism
    \[       \varphi \colon  F       \to  a_*a^*F     \]
    from the adjunction \(a^*\dashv  a_*\).
    This morphism derives to a morphism of functors from \(\mathrm {D}^+(S)\) to \(\mathrm {D}^+(S)\), namely

    \begin{eqn}{5.3.1.1}\setnumber{5.3.1.1}\label{hodge-theory-iii-5.3.1.1}\skipsinglepar\[         \varphi \colon \mathrm {Id}         \to  \mathrm {R} a_*a^*.       \tag{5.3.1.1}       \]\end{eqn}\begin{definition}{5.3.2}\setnumber{5.3.2}\label{hodge-theory-iii-5.3.2}\skipsinglepar\par{}We say that \(a\) is of \emph{cohomological descent} if, for every abelian sheaf \(F\) on \(S\), we have
    \[       F\xrightarrow {\sim }\operatorname {Ker}({a_0}_*{a_0}^*F\to {a_1}_*{a_1}^*F)     \]
    and
    \[       \mathrm {R}^i a_*a^*F       = 0       \quad \text {for }i>0.     \]\par{}This is equivalent to asking that \hyperref[hodge-theory-iii-5.3.1.1]{Equation~\ref*{hodge-theory-iii-5.3.1.1}} be an isomorphism.\end{definition}\label{hodge-theory-iii-5.3.3}\par{}If \(a\) is of cohomological descent, then, for \(K\in \operatorname {Ob}\mathrm {D}^+(S)\), the canonical map

    \begin{eqn}{5.3.3.1}\setnumber{5.3.3.1}\label{hodge-theory-iii-5.3.3.1}\skipsinglepar\[         \mathrm {R}\Gamma (S,K)         \to  \mathrm {R}\Gamma (S,\mathrm {R} a_*a^*K)         \simeq  \mathrm {R}\Gamma (X_\bullet ,a^*K)       \tag{5.3.3.1}       \]\end{eqn}

    is an isomorphism.
    In particular, for \(F\) an abelian sheaf on \(S\), we have a spectral sequence (\hyperref[hodge-theory-iii-5.2.3.2]{Equation~\ref*{hodge-theory-iii-5.2.3.2}})

    \begin{eqn}{5.3.3.2}\setnumber{5.3.3.2}\label{hodge-theory-iii-5.3.3.2}\skipsinglepar\[         E_1^{pq}         = \operatorname {H}^q(X_p,a_p^*F)         \Rightarrow  \operatorname {H}^{p+q}(S,F).       \tag{5.3.3.2}       \]\end{eqn}\par{}For a complex, we again have, in hypercohomology, a spectral sequence

    \begin{eqn}{5.3.3.3}\setnumber{5.3.3.3}\label{hodge-theory-iii-5.3.3.3}\skipsinglepar\[         E_1^{pq}         = \operatorname {\mathbb {H}}^q(X_p,a_p^*K)         \Rightarrow  \operatorname {\mathbb {H}}^{p+q}(S,K).       \tag{5.3.3.3}       \]\end{eqn}\par{}In both cases, the \(E_1^{\bullet  q}\) (for fixed \(q\)) form a simplicial group, and
    \[       d_1       = \sum _i (-1)^i\delta _i       \colon  E_1^{p,q}       \to  E_1^{p+1,q}.     \]\begin{definition}{5.3.4}\setnumber{5.3.4}\label{hodge-theory-iii-5.3.4}\skipsinglepar\par{}A continuous map \(a\colon  X\to  S\) is of \emph{cohomological descent} if the augmentation morphism of \(\operatorname {cosq}(X\to  S)\), namely
    \[       \big ((X/S)^{\Delta _n}\big )_{n\geqslant 0}       \to  S,     \]
    is of cohomological descent.
    We say that \(a\) is of \emph{universal cohomological descent} if, for every \(u\colon  S'\to  S\), the continuous map \(a'\colon  X\times _S S'\to  S'\) is of cohomological descent.\end{definition}\label{hodge-theory-iii-5.3.5}\par{}The fundamental results, proven in [\hyperref[SD]{SD}], are the following.

    
 \begin{enumerate}
      
 \item[(I)]
        The continuous maps of universal cohomological descent form a Grothendieck topology on the category of topological spaces, which we call the \emph{universal cohomological descent topology}.
      


      
 \item[(II)]
        A proper (\cref{hodge-theory-iii-0.1}) surjective map is of universal cohomological descent.
      


      
 \item[(III)]
        A map \(a\colon  X\to  S\) that admits sections locally on \(S\) is of universal cohomological descent.
      


      
 \item[(IV)]
        Let \(a\colon  X_\bullet \to  S\) be a \(k\)-truncated augmented simplicial space (with \(-1\leqslant  k\leqslant \infty \)).
        For \(k\geqslant  n\geqslant -1\), let \(\varphi _n\colon \operatorname {cosq} X_\bullet \to \operatorname {cosq}\operatorname {sq}_n X_\bullet \) be the evident map.
        We say that \(X_\bullet \) is a \emph{\(k\)-truncated hypercover of \(S\), for the universal cohomological descent topology}, if the maps

        \begin{eqn}{5.3.5.1}\setnumber{5.3.5.1}\label{hodge-theory-iii-5.3.5.1}\skipsinglepar\[             (\varphi _n)_{n+1}             \colon  X_{n+1}             \to  (\operatorname {cosq}\operatorname {sq}_n X_\bullet )_{n+1}           \tag{5.3.5.1}           \]\end{eqn}

        (for \(-1\leqslant  n\leqslant  k-1\)) are of universal cohomological descent.
        If \(X_\bullet \) is such a hypercover, then the simplicial space \(\operatorname {cosq}(X_\bullet )\) augmented over \(S\) is of cohomological descent.
      


      
 \item[(V)]
        Let \(a\) be a morphism of simplicial topological spaces augmented over \(S\).
        \[           \begin {CD}             X_\bullet  @>a>> Y_\bullet            \\@VxVV @VVyV           \\S @= S           \end {CD}         \]
        We say that \(a\) is a \emph{hypercover} for the universal cohomological descent topology if the evident maps \(X_n\to (\operatorname {cosq}_{n-1}^{Y_\bullet }\operatorname {sq}_{n-1}X_\bullet )_n\) are of universal cohomological descent.
        If \(a\) is such a hypercover then, for every \(K\in \operatorname {Ob}\mathrm {D}^+(S)\),
        \[           a^*\colon            \mathrm {R} y_*y^*K           \xrightarrow {\sim } \mathrm {R} x_*x^*K.         \]
      

    \end{enumerate}\label{hodge-theory-iii-5.3.6}\par{}For \(k=\infty \), (IV) of \cref{hodge-theory-iii-5.3.5} implies that the \(a\colon  X_\bullet \to  S\) are of cohomological descent if the \((\varphi _n)_{n+1}\) are of universal cohomological descent.
    For \(n=-1,0\), these maps are
    \[       (\varphi _n)_{n+1}       =       \begin {cases}         X_0\xrightarrow {a}S         &\text {if }n=-1       \\X_1\xrightarrow {(\delta _0,\delta _1)}X_0\times _S X_0         &\text {if }n=0.       \end {cases}     \]
    
    For \(n=1\), \((\operatorname {cosq}\operatorname {sq}_1(X_\bullet ))_1\) is the subspace of \(X_1\times _S X_1\times _S X_1\) consisting of the triples \((x,y,z)\) such that \(\delta _0x=\delta _0y\), \(\delta _1x=\delta _0z\), and \(\delta _1y=\delta _1z\).
    The map \((\varphi _1)_2\) is \(x\mapsto (\delta _0x,\delta _1x,\delta _2x)\).\par{}For \(k=0\), (IV) of \cref{hodge-theory-iii-5.3.5} is \hyperref[hodge-theory-iii-5.3.4]{Definition~\ref*{hodge-theory-iii-5.3.4}}.\begin{example}{5.3.7}\setnumber{5.3.7}\label{hodge-theory-iii-5.3.7}\skipsinglepar\par{}Let \(\mathscr {U}=(U_i)_{i\in  I}\) be an open cover, or a finite locally closed cover, of \(S\).
    Let \(X=\coprod _{i\in  I}U_i\).
    Then \(a\colon  X\to  S\) is of cohomological descent.
    The spectral sequence in \hyperref[hodge-theory-iii-5.3.3.2]{Equation~\ref*{hodge-theory-iii-5.3.3.2}} for \(X_\bullet =\operatorname {cosq}(X\to  S)\) is then exactly the Leray spectral sequence of the cover \(\mathscr {U}\).\end{example}\label{hodge-theory-iii-5.3.8}\par{}Let \(a\colon  X_\bullet \to  S\) be as in (IV) of \cref{hodge-theory-iii-5.3.5}.
    We say that \(X_\bullet \) is a \emph{proper \(k\)-truncated hypercover} of \(S\) if the arrows in \hyperref[hodge-theory-iii-5.3.5.1]{Equation~\ref*{hodge-theory-iii-5.3.5.1}} are proper and surjective.
    For \(k=\infty \), we simply say "proper hypercover".\chapter{Examples of simplicial topological spaces}\thispagestyle{fancy}\label{hodge-theory-iii-6}\section{Classifying spaces}\label{hodge-theory-iii-6.1}\label{hodge-theory-iii-6.1.1}\par{}Let \(u\colon  X\to  S\) be a continuous map.
    For every sheaf \(F\) on \(S\), the sheaf \(u^*F\) is endowed with a "descent data" with respect to \(u\), i.e. we have an isomorphism between the two inverse images of \(u^*F\) on \(X\times _S X\), and this isomorphism satisfies a cocycle condition.
    If \(u\) admits a section locally on \(S\), then this construction defines an equivalence between the category of sheaves on \(S\) and that of sheaves on \(X\) endowed with a descent data.\par{}Take \(X\) to be a (left) principal homogeneous space for the group \(G\) on a space \(S\).
    Then the \(G\)-equivariant sheaves on \(X\) are exactly the sheaves endowed with a descent data: every equivariant sheaf on \(X\) is, in a unique way (as an equivariant sheaf), the inverse image of a sheaf on \(S=X/G\).\label{hodge-theory-iii-6.1.2}\par{}If a topological group \(G\) acts on a space \(X\), then \(G\) acts on \(G^{\Delta _n}\times  X\) by
    \[       g\cdot (g_0,\ldots ,g_n,x)       = (g_0g^{-1},\ldots ,g_ng^{-1},gx).     \]
    We denote by \([X/G]_\bullet \) the simplicial space

    \begin{eqn}{6.1.2.1}\setnumber{6.1.2.1}\label{hodge-theory-iii-6.1.2.1}\skipsinglepar\[         [X/G]_\bullet          = \big ((G^{\Delta _n}\times  X)/G\big )_{n\geqslant 0}.       \tag{6.1.2.1}       \]\end{eqn}\begin{enumerate}
    
 \item[(a)]
      If \(X\) is a principal homogeneous space for the group \(G\) on \(S=X/G\), then the map
      \[         \begin {aligned}           G^{\Delta _n}\times  X           &\to  X^{\Delta _n}         \\(g_0,\ldots ,g_n,x)           &\mapsto  (g_0x,\ldots ,g_nx)         \end {aligned}       \]
      identifies \([X/G]_n\) with the iterated fibre product \((X/S)^{\Delta _n}\):
      \[         \operatorname {cosq}(X\to  S)         = ([X/G]_\bullet \to  S).       \]
      
      In particular, we have (by (III) of \cref{hodge-theory-iii-5.3.5})

      \begin{eqn}{6.1.2.2}\setnumber{6.1.2.2}\label{hodge-theory-iii-6.1.2.2}\skipsinglepar\[           \operatorname {H}^\bullet ([X/G]_\bullet )           \simeq  \operatorname {H}^\bullet (X/G)         \tag{6.1.2.2}         \]\end{eqn}

      (for a principal homogeneous space).
    


    
 \item[(b)]
      For all \(n\), \(G^{\Delta _n}\times  X\) is a principal homogeneous space for the group \(G\) on \([X/G]_n\).
      For every equivariant sheaf \(F\) on \(X\), \(\operatorname {pr}_2^*F\) is an equivariant sheaf on \(G^{\Delta _n}\times  X\);
      by \cref{hodge-theory-iii-6.1.1}, the latter is the inverse image of \(F^n\) on \([X/G]_n\).
      Every equivariant sheaf \(F\) on \(X\) thus defines a sheaf on \([X/G]_\bullet \).
      It is easy to show that we thus obtain an equivalence between the category of equivariant sheaves on \(X\) and the category of sheaves \(F^\bullet \) on \([X/G]_\bullet \) that satisfy

      \begin{property}{*}\setnumber{*}\label{hodge-theory-iii-6.1.star}\skipsinglepar\par{}\emph{\((*)\) For every \(f\colon \Delta _n\to \Delta _m\), the structure morphism \(f^*F^n\to  F^m\) is an isomorphism.}\end{property}
    


    
 \item[(c)]
      Construction (b) above is natural in \((G,X,F)\).
      We set

      \begin{eqn}{6.1.2.3}\setnumber{6.1.2.3}\label{hodge-theory-iii-6.1.2.3}\skipsinglepar\[           \operatorname {H}^\bullet (X,G;F)           = \operatorname {H}^\bullet ([X/G]_\bullet ,F^\bullet )         \tag{6.1.2.3}         \]\end{eqn}

      (mixed cohomology of \(X,G\) with coefficients in \(F\)).
      Under the hypotheses of (a), if \(F\) is the inverse image of \(F^{-1}\) on \(S=X/G\), then

      \begin{eqn}{6.1.2.4}\setnumber{6.1.2.4}\label{hodge-theory-iii-6.1.2.4}\skipsinglepar\[           \operatorname {H}^\bullet (X,G;F)           = \operatorname {H}^\bullet (X/G,F^{-1})         \tag{6.1.2.4}         \]\end{eqn}

      (for a principal homogeneous space).
      This generalises \hyperref[hodge-theory-iii-6.1.2.2]{Equation~\ref*{hodge-theory-iii-6.1.2.2}} (which is the case \(F=\underline {\mathbb {Z}}\)).
    

  \end{enumerate}\label{hodge-theory-iii-6.1.3}\par{}Let \(P^t\) be the topological space consisting of a single point.
    We define the \emph{simplicial classifying space} of \(G\), denoted \(B_{\bullet  G}\), to be the simplicial space
    \[       B_{\bullet  G}       = [P^t/G]_\bullet .     \]\par{}Let \(P\) be a principal homogeneous space for the group \(G\) on \(S\).
    The evident morphism
    \[       \operatorname {cosq}(P\to  S)       = [P/G]_\bullet        \to  [P^t/G]_\bullet        = B_{\bullet  G}     \]
    defines a composite morphism

    \begin{eqn}{6.1.3.1}\setnumber{6.1.3.1}\label{hodge-theory-iii-6.1.3.1}\skipsinglepar\[         [P]         \colon  \operatorname {H}^\bullet (B_{\bullet  G})         \to  \operatorname {H}^\bullet ([P/G]_\bullet )         \xleftarrow [(6.1.2.2)]{\sim } \operatorname {H}^\bullet (S).       \tag{6.1.3.1}       \]\end{eqn}\par{}We will see below that, in good cases, \(\operatorname {H}^\bullet (B_{\bullet  G})=\operatorname {H}^\bullet (B_G)\), and that the image of \([P]\) consists of the \emph{characteristic classes} of \(P\).\label{hodge-theory-iii-6.1.4}\par{}Let \(G\) be a Lie group, \(B_G\) a classifying space for \(G\), and \(a\colon  U_G\to  B_G\) the universal principal homogeneous \(G\)-space.
    Let \(X\) be a \(G\)-space;
    then \(X\times  U_G\) is a principal homogeneous \(G\)-space over \(X\times  U_G/G\) such that, for every equivariant sheaf \(F\) on \(X\), \(\operatorname {pr}_1^*F\) is the inverse image of a sheaf \(F^G\) on \(X\times  U_G/G\).
    Since \(U_G\) is contractible, and by \hyperref[hodge-theory-iii-6.1.2.2]{Equation~\ref*{hodge-theory-iii-6.1.2.2}}, we have

    \begin{eqn}{6.1.4.1}\setnumber{6.1.4.1}\label{hodge-theory-iii-6.1.4.1}\skipsinglepar\[         \operatorname {H}^\bullet (X,G;F)         \xrightarrow {\sim } \operatorname {H}^\bullet (X\times  U_G,G;\operatorname {pr}_1^*F)         \xleftarrow {\sim } \operatorname {H}^\bullet (X\times  U_G/G,F^G).       \tag{6.1.4.1}       \]\end{eqn}\par{}In particular, for \(X=P^t\),

    \begin{eqn}{6.1.4.2}\setnumber{6.1.4.2}\label{hodge-theory-iii-6.1.4.2}\skipsinglepar\[         \operatorname {H}^\bullet (B_{\bullet  G})         = \operatorname {H}^\bullet (B_G).       \tag{6.1.4.2}       \]\end{eqn}\par{}We can see that the isomorphism in \hyperref[hodge-theory-iii-6.1.4.2]{Equation~\ref*{hodge-theory-iii-6.1.4.2}} is a particular case of \hyperref[hodge-theory-iii-6.1.3.1]{Equation~\ref*{hodge-theory-iii-6.1.3.1}} where \(S=B_G\) and \(P=U_G\).\label{hodge-theory-iii-6.1.5}\par{}The spectral sequence in \hyperref[hodge-theory-iii-5.2.3.2]{Equation~\ref*{hodge-theory-iii-5.2.3.2}} for \(B_{\bullet  G}\)
    \[       E_1^{pq}       = \operatorname {H}^q(G^{\Delta _p}/G)       \Rightarrow  \operatorname {H}^{p+q}(B_{\bullet  G})       =\operatorname {H}^{p+q}(B_G)     \]
    is essentially the \emph{Eilenberg–Moore spectral sequence}.
    We briefly recall how it allows us to relate the rational cohomologies of \(G\) and of \(B_G\), for connected \(G\).\begin{enumerate}
    
 \item[(a)]
      The algebra \(\operatorname {H}^\bullet (G,\mathbb {Q})\) is a connected graded Hopf algebra of finite dimension over \(\mathbb {Q}\).
      If \(P^\bullet (G)\) is the graded module of its primitive elements, then we have
      \[         H^\bullet (G,\mathbb {Q})         = \bigwedge  P^\bullet (G)       \]
      and the generators of \(P^\bullet (G)\) are of odd degree.
    


    
 \item[(b)]
      The simplicial algebra \((E_1^{p\bullet })_{p\geqslant 0}\) is
      \[         E_1^{p\bullet }         = \bigwedge \big (P^\bullet (G)^{\Delta _p}/P^\bullet (G)\big )       \]
      which is the exterior algebra of the suspension of the constant cosimplicial module \(P^\bullet (G)\);
      we thus have (by Quillen [\hyperref[Q1968]{Q1968}]) that \(E_2^{p\bullet }=\operatorname {Sym}^p(P^\bullet (G))\).
      The pages \(E_2^{pq}\) are only zero for \(p+q\) even;
      we thus have that \(E_2^{pq}=E_\infty ^{pq}\), and, for a suitable filtration, we canonically have that
      \[         \operatorname {Gr}\operatorname {H}^\bullet (B_G,\mathbb {Q})         \simeq  \operatorname {Sym}^\bullet (P^\bullet (G)[-1])       \]
      and non-canonically that
      \[         \operatorname {H}^\bullet (B_G,\mathbb {Q})         \simeq  \operatorname {Sym}^\bullet (P^\bullet (G)[-1]).       \]
    

  \end{enumerate}\label{hodge-theory-iii-6.1.6}\par{}Let \(G\) be a (complex) linear algebraic group.\begin{enumerate}
    
 \item[(a)]
      If \(T\) is a maximal torus of \(G\), with Weyl group \(W\), then

      \begin{eqn}{6.1.6.1}\setnumber{6.1.6.1}\label{hodge-theory-iii-6.1.6.1}\skipsinglepar\[           \operatorname {H}^\bullet (B_G,\mathbb {Q})           \xrightarrow {\sim } \operatorname {H}^\bullet (B_T,\mathbb {Q})^W.         \tag{6.1.6.1}         \]\end{eqn}
    


    
 \item[(b)]
      If \(T\) is a torus with character group \(X(T)\), then

      \begin{eqn}{6.1.6.2}\setnumber{6.1.6.2}\label{hodge-theory-iii-6.1.6.2}\skipsinglepar\[           \operatorname {H}^\bullet (T,\mathbb {Z})           \simeq  \bigwedge ^\bullet  X(T)         \tag{6.1.6.2}         \]\end{eqn}

      (an isomorphism of graded Hopf algebras).
    

  \end{enumerate}\par{}We will only use (a) in the following weaker form:

    
 \begin{enumerate}
      
 \item[(a')]
        \emph{(The splitting principle).}
        The map \(\operatorname {H}^\bullet (B_G,\mathbb {Q})\to \operatorname {H}^\bullet (B_T,\mathbb {Q})\) is injective.
      

    \end{enumerate}\par{}For completion, we recall a proof of (a').\par{}If \(B\) is a Borel subgroup of \(G\), then the bundle \(U_G/B\) on \(B_G\) is a fibre in flag spaces.
    By [2.1 and 2.6.3, \hyperref[D1968]{D1968}], which is better explained in [Proposition 3.1, \hyperref[G1970]{G1970}], or by [\hyperref[B1956]{B1956}], the Leray spectral sequence of \(U_G/B\to  B_G\) degenerates to rational cohomology.
    We thus have that \(\operatorname {H}^\bullet (B_G,\mathbb {Q})\hookrightarrow \operatorname {H}^\bullet (U_G/B,\mathbb {Q})\).
    We conclude by noting that \(U_G/B\sim  B_B\sim  B_T\).\section{Construction of hypercovers}\label{hodge-theory-iii-6.2}\label{hodge-theory-iii-6.2.1}\par{}Let \(X_\bullet \) be a simplicial set.
    Denote by \(D(\Delta _n,\Delta _M)\) the set of increasing surjective maps from \(\Delta _n\) to \(\Delta _m\) (degeneracy operators), and set

    \begin{eqn}{6.2.1.1}\setnumber{6.2.1.1}\label{hodge-theory-iii-6.2.1.1}\skipsinglepar\[         N(X_n)         = X_n\setminus \bigcup _{{s\in  D(\Delta _n,\Delta _{n-1})}}s(X_{n-1}).       \tag{6.2.1.1}       \]\end{eqn}\par{}Recall that, for all \(n\), the map

    \begin{eqn}{6.2.1.2}\setnumber{6.2.1.2}\label{hodge-theory-iii-6.2.1.2}\skipsinglepar\[         \coprod  s         \colon  \coprod _{\mathclap {\substack {m\leqslant  n,\\s\in  D(\Delta _n,\Delta _m)}}} N(X_m)         \to  X_n       \tag{6.2.1.2}       \]\end{eqn}

    is bijective.\begin{definition}{6.2.2}\setnumber{6.2.2}\label{hodge-theory-iii-6.2.2}\skipsinglepar\par{}We say that a simplicial topological space is \emph{\(s\)-split} if the maps in \hyperref[hodge-theory-iii-6.2.1.2]{Equation~\ref*{hodge-theory-iii-6.2.1.2}} are homeomorphisms.\end{definition}\par{}Let \(X_\bullet \) be a \(k\)-truncated simplicial set.
  For \(n\leqslant  k\), we again define \(N(X_n)\) by \hyperref[hodge-theory-iii-6.2.1.1]{Equation~\ref*{hodge-theory-iii-6.2.1.1}}, and then \hyperref[hodge-theory-iii-6.2.1.2]{Equation~\ref*{hodge-theory-iii-6.2.1.2}} is a bijection.
  We say that a \(k\)-truncated simplicial topological space is \emph{\(s\)-split} if \hyperref[hodge-theory-iii-6.2.1.2]{Equation~\ref*{hodge-theory-iii-6.2.1.2}} is a homeomorphism for \(n\leqslant  k\).\label{hodge-theory-iii-6.2.3}\par{}For an \(s\)-split \((n+1)\)-truncated topological simplicial space \(X\) augmented over \(S\), let \(\alpha (X)\) be the triple consisting of \(\operatorname {sq}_n(X)\), \(NX_{n+1}\), and the evident map from \(NX_{n+1}\) to \((\operatorname {cosq}\operatorname {sq}_n X)_{n+1}\).\par{}This triple \(\alpha (X)=(Y,N,\beta )\) satisfies the following:

    \begin{property}{*}\setnumber{*}\label{hodge-theory-iii-6.2.star}\skipsinglepar\par{}\emph{\((*)\) \(Y\) is an \(s\)-split \(n\)-truncated simplicial topological space augmented over \(S\), and \(\beta \) is a continuous map from \(N\) to \((\operatorname {cosq} Y)_{n+1}\).}\end{property}\begin{proposition}{6.2.4}\setnumber{6.2.4}\label{hodge-theory-iii-6.2.4}\skipsinglepar\par{}Let \((Y,N,\beta )\) satisfy \hyperref[hodge-theory-iii-6.2.star]{Property~\ref*{hodge-theory-iii-6.2.star}} above.

    
 \begin{enumerate}
      
 \item[(i)]
        Up to unique isomorphism, there exists exactly one \(s\)-split \((n+1)\)-truncated topological space \(X\) augmented over \(S\) such that \(\alpha (X)\simeq (Y,N,\beta )\).
      


      
 \item[(ii)]
        It is equivalent to give either \(f\colon  X\to  Z\) or:

        
 \begin{enumerate}
          
 \item[(a)]
            a morphism \(f'\colon  Y\to \operatorname {sq}_n(Z)\); and
          

          
 \item[(b)]
            a morphism \(f''\colon  N\to  Z_{n+1}\), such that the diagram
            \[               \begin {CD}                 N @>\beta >> (\operatorname {cosq} Y)_{n+1}               \\@V{f''}VV @VV{f'}V               \\Z_{n+1} @>>> (\operatorname {cosq}\operatorname {sq}_n Z)_{n+1}               \end {CD}             \]
            commutes.
          

        \end{enumerate}

      

    \end{enumerate}\begin{proof}\par{}[5.1.3, \hyperref[SD]{SD}].\end{proof}\par{}This proposition also applies to simplicial objects in other categories \(\mathscr {C}\) apart from that of topological spaces; it suffices that \(\mathscr {C}\) satisfy the following:

    \label{hodge-theory-iii-6.2.4.1}\par{}Finite projective limits exist in \(\mathbb {C}\).
        Finite projective sums exist in \(\mathbb {C}\), and they are disjoint and universal.\end{proposition}\label{hodge-theory-iii-6.2.5}\par{}\hyperref[hodge-theory-iii-6.2.4]{Proposition~\ref*{hodge-theory-iii-6.2.4}} allows us to construct, by induction, proper hypercovers of \(S\).\begin{enumerate}
    
 \item[(0)]
      We take \(f_0\colon  X_0\to  S\), proper and surjective.
      Then \(\{X_0\}\) is a \(0\)-truncated proper hypercover of \(S\) (\cref{hodge-theory-iii-5.3.8}), and it is \(s\)-split.
    


    
 \item[(1)]
      We take \(f_1\colon  N_1\to \operatorname {cosq}(\{X_0\})_1\), i.e. \(f_1\colon  N_1\to  X_0\times _S X_0\).
      Applying \hyperref[hodge-theory-iii-6.2.4]{Proposition~\ref*{hodge-theory-iii-6.2.4}}, we associate to \(f_1\) the \(s\)-split 1-truncated augmented simplicial topological space
      \[         {}_1X_\bullet          = \left (         N_1{\textstyle \coprod } X_0         \xleftarrow [\to ]{\to } X_0         \to  S         \right ).       \]
      Suppose \(f_1\) to be chosen such that
      \[         f'_1         \colon  N_1{\textstyle \coprod } X_0\to  X_0\times _S X_0       \]
      is proper and surjective (for example, if \(f_1\) is proper and surjective).
      Then \({}_1X_\bullet \) is an \(s\)-split 1-truncated proper hypercover of \(S\).
    


    
 \item[(k+1)]
      Assume that we have already constructed an \(s\)-split \(k\)-truncated proper hypercover \({}_kX_\bullet \to  S\).
      We take \(f_{k+1}\colon  N_{k+1}\to (\operatorname {cosq}({}_kX_\bullet ))_{k+1}\) and, applying \hyperref[hodge-theory-iii-6.2.4]{Proposition~\ref*{hodge-theory-iii-6.2.4}}, we associate to \(f_{k+1}\) an \(s\)-split \((k+1)\)-truncated augmented semi-simplicial space \({}_{k+1}X_\bullet \).
      Suppose that \(f_{k+1}\) is such that
      \[         f'_{k+1}         \colon  {}_{k+1}X_{k+1}         \to  \operatorname {cosq}({}_kX_\bullet )_{k+1}       \]
      is proper and surjective (for example, if \(f_{k+1}\) is proper and surjective).
      Then \({}_{k+1}X_\bullet \) is an \(s\)-split \((k+1)\)-truncated proper hypercover of \(S\).
    


    
 \item[(∞)]
      The \({}_kX_\bullet \) thus constructed are the successive skeletons of an \(s\)-split proper hypercover of \(S\).
    

  \end{enumerate}\label{hodge-theory-iii-6.2.6}\par{}We say that a simplicial scheme \(X_\bullet \) over \(\mathbb {C}\) is \emph{smooth} if the \(X_n\) are smooth;
    it is said to be \emph{proper} if the \(X_n\) are compact.
    A \emph{normal crossing divisor} \(D_\bullet \) of \(X_\bullet \), assumed to be smooth, is a family \(D_n\subset  X_n\) of normal crossing divisors (\cref{hodge-theory-ii-3.1.2}) such that the \(U_n=X_n\setminus  D_n\) form a simplicial subscheme \(U_\bullet \) of \(X_\bullet \).
    This definition is justified by the following lemma.\begin{lemma}{6.2.7}\setnumber{6.2.7}\label{hodge-theory-iii-6.2.7}\skipsinglepar\par{}If \(D_\bullet \) is a normal crossing divisor of \(X_\bullet \), then the logarithmic de Rham complexes \((\Omega _{X_n}^\bullet (\log  D_n))_{n\geqslant 0}\), endowed with the weight filtration (\cref{hodge-theory-ii-3.1.5}), form a filtered complex on \(X_\bullet \).\begin{proof}\par{}This follows from (ii) of \hyperref[hodge-theory-ii-3.1.3]{Proposition~\ref*{hodge-theory-ii-3.1.3}}.\end{proof}\par{}We denote the complex \((\Omega _{X_n}^\bullet (\log  D_n))_{n\geqslant 0}\) by \(\Omega _{X_\bullet }^\bullet (\log  D_\bullet )\).\end{lemma}\label{hodge-theory-iii-6.2.8}\par{}Using \cref{hodge-theory-iii-6.2.5}, we can show that, for every separated scheme \(S\) over \(\mathbb {C}\), there exists:

    
 \begin{enumerate}
      
 \item[(𝛼)]
        
        a simplicial scheme \(X_\bullet \) over \(\mathbb {C}\), smooth and proper, that we can take to be \(s\)-split;
      

      
 \item[(𝛽)]
        a normal crossing divisor \(D_\bullet \) of \(X_\bullet \); we set \(U_\bullet =X_\bullet \setminus  D_\bullet \);
      

      
 \item[(𝛾)]
        an augmentation \(a\colon  U_\bullet \to  S\) that realises \(U_\bullet ^\mathrm {an}\) as a proper hypercover of \(S^\mathrm {an}\).
      

    \end{enumerate}\par{}Furthermore, any two such systems are covered by a third, and a morphism \(u\colon  S\to  T\) can be covered by a morphism
    
  \begin {tikzcd}
      X_\bullet 
        \ar [r]
      & Y_\bullet 
    \\U_\bullet 
        \ar [r]
        \ar [u,hook]
        \ar [d,swap,"a"]
      & V_\bullet 
        \ar [u,hook]
        \ar [d,"b"]
    \\S \ar [r,swap,"u"]
      & T
    \end {tikzcd}

    of such systems (see [\hyperref[SD]{SD}]).\section{Relative cohomology}\label{hodge-theory-iii-6.3}\label{hodge-theory-iii-6.3.1}\par{}The \emph{mapping cone} construction for morphisms of simplicial sets works in the same way for simplicial objects in any category \(\mathscr {C}\) that has a final object \(e\) and finite sums.
    For \(u\colon  Y_\bullet \to  X_\bullet \), the cone \(C(u)\) satisfies
    \[       C(u)_n       = X_n{\,\textstyle \coprod \,}\coprod _{i<n} Y_i{\,\textstyle \coprod \,}e.     \]\par{}We take \(\mathscr {C}\) to be:

    
 \begin{enumerate}
      
 \item[(a)]
        the category of topological spaces and continuous maps, with final object \(e=P^t\);
      

      
 \item[(b)]
        the category of pairs \((X,F)\) where \(X\) is a topological space and \(F\) is an abelian sheaf on \(X\), with an arrow \((u,f)\colon (Y,F)\to (X,G)\) consisting of a continuous map \(u\colon  Y\to  X\) and a \(u\)-morphism (\cref{hodge-theory-iii-5.1.5}) \(f\colon  G\to  F\), with final object \(e=(P^t,0)\).
      

    \end{enumerate}\label{hodge-theory-iii-6.3.2}\par{}Let \(u\colon  Y_\bullet \to  X_\bullet \) be a morphism of topological simplicial spaces, with cone \(C(u)\).
    Let \(F\) be an abelian sheaf on \(X_\bullet \) and \(G\) an abelian sheaf on \(Y_\bullet \), and let \(f\colon  G\to  F\) be a \(u\)-morphism.
    The cone \(C(f)\) of \(f\) is an abelian sheaf on \(C(u)\), and we set

    \begin{eqn}{6.3.2.1}\setnumber{6.3.2.1}\label{hodge-theory-iii-6.3.2.1}\skipsinglepar\[         \operatorname {H}^n(C(u),C(f))         = \operatorname {H}^n(X_\bullet \operatorname {mod} Y_\bullet , F\operatorname {mod} G).       \tag{6.3.2.1}       \]\end{eqn}

    These are the \emph{relative cohomology groups}.
    We can easily show that they fit into a long exact sequence

    \begin{eqn}{6.3.2.2}\setnumber{6.3.2.2}\label{hodge-theory-iii-6.3.2.2}\skipsinglepar\[         \ldots          \to  \operatorname {H}^i(X_\bullet \operatorname {mod} Y_\bullet , F\operatorname {mod} G)         \to  \operatorname {H}^i(X_\bullet ,F)         \to  \operatorname {H}^i(Y_\bullet ,G)         \to  \ldots .       \tag{6.3.2.2}       \]\end{eqn}\label{hodge-theory-iii-6.3.3}\par{}More generally, let \(L\) and \(K\) be bounded-below complexes of abelian sheaves on \(Y_\bullet \) and \(X_\bullet \) (respectively), and let \(f\colon  K\to  L\) be a \(u\)-morphism.
    We thus obtain a complex \(C(f)\) on \(C(u)\).
    We again define the hypercohomology
    \[       \operatorname {\mathbb {H}}^n(C(u),C(f))       = \operatorname {\mathbb {H}}^n(X_\bullet \operatorname {mod} Y_\bullet , K\operatorname {mod} L).     \]
    
    These groups appear in an exact sequence analogous to that of \hyperref[hodge-theory-iii-6.3.2.2]{Equation~\ref*{hodge-theory-iii-6.3.2.2}}, coming, in the suitable derived category, from a distinguished triangle

    \begin{eqn}{6.3.3.1}\setnumber{6.3.3.1}\label{hodge-theory-iii-6.3.3.1}\skipsinglepar\end{eqn}\label{hodge-theory-iii-6.3.4}\par{}The construction presented above is not the only one possible.
    It has the inconvenience that, even if we start with true topological spaces \(X\) and \(Y\) (i.e. constant simplicial spaces), we are led to consider non-constant simplicial spaces.\section{Multisimplicial spaces}\label{hodge-theory-iii-6.4}\label{hodge-theory-iii-6.4.1}\par{}Let \(r\geqslant 0\) be an integer.
    An \emph{\(r\)-simplicial object} \(Z_\bullet \) in a category \(\mathscr {C}\) is a contravariant functor from the \(r\)-fold product \((\Delta )^r\) to \(\mathscr {C}\).
    The \emph{diagonal} simplicial object \(\delta  Z_\bullet \) is the composite functor \((\Delta )\to (\Delta )^r\to \mathscr {C}\).\label{hodge-theory-iii-6.4.2}\par{}As in \cref{hodge-theory-iii-5.1.6}, we define the topos of \emph{sheaves} on an \(r\)-simplicial topological space.
    Let \(\Gamma ^\bullet \) be the functor
    \[       F       \mapsto  (\Gamma (X_{n_1\ldots  n_r},F^{n_1\ldots  n_r}))     \]
    from sheaves on \(X_\bullet \) to \(r\)-cosimplicial sets.
    For small \(r\), we often prefer to write \(\Gamma ^{\bullet \ldots \bullet }\) (with \(r\)-many copies of \(\bullet \)).
    We have a co-augmentation \(\Gamma (X_\bullet ,F^\bullet )\to \Gamma ^\bullet (X_\bullet ,F^\bullet )\).
    The functors of \emph{cohomology with values in an abelian sheaf \(F\)} are the derived functors of the "global sections" functor \(\Gamma \).
    They can be calculated by a procedure parallel to that of \cref{hodge-theory-iii-5.2.3}, i.e.

    \begin{eqn}{6.4.2.1}\setnumber{6.4.2.1}\label{hodge-theory-iii-6.4.2.1}\skipsinglepar\[         \mathrm {R}\Gamma          =\mathrm {s}\mathrm {R}\Gamma ^\bullet          \colon  \mathrm {D}^+(X_\bullet )         \to \mathrm {D}^+((\mathsf {Ab})).       \tag{6.4.2.1}       \]\end{eqn}\par{}A sheaf \(F\) on an \(r\)-simplicial topological space \(X_\bullet \) induces a sheaf \(\delta  F\) on the diagonal simplicial space \(\delta  X_\bullet \).
    It follows from the Cartier–Eilenberg–Zilber theorem that

    \begin{eqn}{6.4.2.2}\setnumber{6.4.2.2}\label{hodge-theory-iii-6.4.2.2}\skipsinglepar\[         \mathrm {R}\Gamma (X_\bullet ,K)         \xrightarrow {\sim } \mathrm {R}\Gamma (\delta  X_\bullet ,\delta  K).       \tag{6.4.2.2}       \]\end{eqn}\label{hodge-theory-iii-6.4.3}\par{}We restrict to the case \(r=2\).
    A \emph{bisimplicial object 1-augmented over a simplicial object \(S_\bullet \)} is a contravariant functor from \((\Delta ^+)\times (\Delta )\) to \(\mathscr {C}\) such that \(S_\bullet \) is the composite functor
    \[       (\Delta )       \xrightarrow {(\Delta _{-1},\bullet )} (\Delta ^+)\times (\Delta )       \to  \mathscr {C}.     \]
    To denote a bisimplicial object 1-augmented over \(S_\bullet \), with underlying bisimplicial object \(X_{\bullet \bullet }\), we will use notation of the form
    \[       a\colon  X_{\bullet \bullet }       \to  S_\bullet .     \]
    
    For \(n\geqslant 0\), \(a\colon (X_{\bullet  n})\to  S_n\) is a simplicial object augmented over \(S_n\).
    If \(F\) is a sheaf on \(X_{\bullet \bullet }\), then the \((a_*(F|X_{\bullet  n})\text { on }S_n)_{n\geqslant 0}\) form a sheaf on \(S_\bullet \);
    we thus define a topos morphism
    \[       a       \colon  X_{\bullet \bullet }^\sim        \to  S_\bullet ^\sim .     \]\par{}We explain in [\hyperref[SD]{SD}] that \(\mathrm {R} a_*\) can be calculated "component by component", i.e.
    \[       \mathrm {R} a_*K|S_n       = \mathrm {R} a_*(K|X_{\bullet  n}).     \]
    It thus follows that if, for each \(n\), the morphism \(a\colon (X_{\bullet  n})\to  S_n\) is of cohomological descent, then \(a\colon  X_{\bullet \bullet }\to  S_\bullet \) is of cohomological descent: for every complex \(K\in \mathrm {D}^+(S_\bullet )\) of abelian sheaves, we have
    \[       K       \xrightarrow {\approx } \mathrm {R} a_*a^*K.     \]\label{hodge-theory-iii-6.4.4}\par{}In [\hyperref[SD]{SD}], we show that, for every separated simplicial scheme \(S_\bullet \), there exists a bisimplicial scheme \(X_{\bullet \bullet }\) that is 1-augmented over \(S_\bullet \), along with \(i\colon  X_{\bullet \bullet }\hookrightarrow \overline {X}_{\bullet \bullet }\) such that:

    
 \begin{enumerate}
      
 \item[(a)]
        The \(\overline {X}_{nm}\) are smooth and projective;
        \(X_{nm}\) is the complement of a normal crossing divisor \(D_{nm}\) in \(\overline {X}_{nm}\), and we can suppose the \(D_{nm}\) to be a union of smooth divisors.
      


      
 \item[(b)]
        For \(n\geqslant 0\), \(X_{\bullet  n}\) is a proper hypercover of \(S_n\), and we can take it to be \(s\)-split.
      

    \end{enumerate}\par{}The construction proceeds as in \cref{hodge-theory-iii-6.2}, but the induction is more complicated.
    The claims of uniqueness (\cref{hodge-theory-iii-6.2.8}) remain true, mutatis mutandis.\chapter{Supplements to §1}\thispagestyle{fancy}\label{hodge-theory-iii-7}\section{Filtered derived category}\label{hodge-theory-iii-7.1}\par{}This section completes \cref{hodge-theory-ii-1.4}.\label{hodge-theory-iii-7.1.1}\par{}Let \(\mathscr {A}\) be an abelian category.
    We set:

    \begin{itemize}\item{}\(\mathrm {F}\mathscr {A}\) (resp. \(\mathrm {F}_2\mathscr {A}\))
        = the category of filtered (resp. bifiltered) objects with finite filtration(s) of \(\mathscr {A}\).
      \item{}\(\mathrm {K}^+\mathrm {F}\mathscr {A}\) (resp. \(\mathrm {K}^+\mathrm {F}_2\mathscr {A}\))
        = the category of filtered (resp. bifiltered) bounded-below complexes of objects of \(\mathscr {A}\), up to homotopy that respects the filtration(s).
      \item{}\(\mathrm {D}^+\mathrm {F}\mathscr {A}\) (resp. \(\mathrm {D}^+\mathrm {F}_2\mathscr {A}\))
        = the triangulated category induced from \(\mathrm {K}^+\mathrm {F}\mathscr {A}\) (resp. from \(\mathrm {K}^+\mathrm {F}_2\mathscr {A}\)) by inverting the filtered (resp. bifiltered) quasi-isomorphisms (\hyperref[hodge-theory-ii-1.3.6]{Definition~\ref*{hodge-theory-ii-1.3.6}});
        this is the \emph{derived filtered category} (resp. \emph{derived bifiltered category}).\end{itemize}\label{hodge-theory-iii-7.1.2}\par{}A filtered quasi-isomorphism \(u\colon (K,F)\to (K',F')\) induces an isomorphism of spectral sequences \(u\colon  E_\bullet ^{\bullet \bullet }(K,F)\to  E_\bullet ^{\bullet \bullet }(K',F')\).
    An object \(K\) of \(\mathrm {D}^+\mathrm {F}\mathscr {A}\) thus defines a spectral sequence \(E_\bullet ^{\bullet \bullet }(K)\).
    
    Similarly, an object \(L\) of \(\mathrm {D}^+\mathrm {F}_2\mathscr {A}\) defines an accumulation of spectral sequences of the type considered in \cref{hodge-theory-ii-1.4.9}.\label{hodge-theory-iii-7.1.3}\par{}Let \(T\) be a left exact functor from \(\mathscr {A}\) to an abelian category \(\mathscr {B}\).
    Suppose that every object of \(\mathscr {A}\) injects into an injective object.
    The functor \(T\) can then be "derived" to give the functors

    \begin{eqn}{7.1.3.1}\setnumber{7.1.3.1}\label{hodge-theory-iii-7.1.3.1}\skipsinglepar\[         \mathrm {R}\colon          \mathrm {D}^+(\mathscr {A})         \to  \mathrm {D}^+(\mathscr {B})       \tag{7.1.3.1}       \]\end{eqn}

    \begin{eqn}{7.1.3.2}\setnumber{7.1.3.2}\label{hodge-theory-iii-7.1.3.2}\skipsinglepar\[         \mathrm {R}\colon          \mathrm {D}^+\mathrm {F}(\mathscr {A})         \to  \mathrm {D}^+\mathrm {F}(\mathscr {B})       \tag{7.1.3.2}       \]\end{eqn}

    \begin{eqn}{7.1.3.3}\setnumber{7.1.3.3}\label{hodge-theory-iii-7.1.3.3}\skipsinglepar\[         \mathrm {R}\colon          \mathrm {D}^+\mathrm {F}_2(\mathscr {A})         \to  \mathrm {D}^+\mathrm {F}_2(\mathscr {B}).       \tag{7.1.3.3}       \]\end{eqn}

    They can be calculated as follows: if \(K'\) is a \(T\)-acyclic resolution (resp. a filtered resolution, resp. a bifiltered resolution) of \(K\) (\cref{hodge-theory-ii-1.4.5} and \cref{hodge-theory-ii-1.4.11}), then \(\mathrm {R} T(K)=T(K')\).\par{}The hypercohomology spectral sequence (for \(T\)) of \(K\in \operatorname {Ob}\mathrm {D}^+\mathrm {F}(\mathscr {A})\) is the spectral sequence of \(\mathrm {R} TK\in \operatorname {Ob}\mathrm {D}^+\mathrm {F}(\mathscr {B})\) (cf. \cref{hodge-theory-ii-1.4.5}).\label{hodge-theory-iii-7.1.4}\par{}We will need more precise results for the functors \(\mathrm {R} a_*\), where \(a\colon  X_\bullet \to  S\) is an augmentation of a simplicial topological space.
    The case \(S=P^t\), where \(\mathrm {R} a_*=\mathrm {R}\Gamma \) will suffice.\par{}We reuse the notation of \cref{hodge-theory-iii-5.1.12}, and recall \hyperref[hodge-theory-iii-5.2.6.1]{Equation~\ref*{hodge-theory-iii-5.2.6.1}}:
    \[       \mathrm {R} a_*       = \mathrm {s}\mathrm {R} {a_\bullet }_*.     \]
    For every complex \(K\in \operatorname {Ob}\mathrm {C}^+(S_\bullet )\), the simple complex \(\mathrm {s} K\) is endowed with a natural filtration \(L\) (\hyperref[hodge-theory-iii-5.1.9.3]{Equation~\ref*{hodge-theory-iii-5.1.9.3}}).
    A quasi-isomorphism \(u\colon  K'\xrightarrow {\sim } K''\) induces a filtered quasi-isomorphism \(u\colon (\mathrm {s} K',L)\xrightarrow {\sim }(\mathrm {s} K'',L)\).
    Then \(\mathrm {s}\) factors as
    \[       \mathrm {s}\colon  \mathrm {D}^+(S_\bullet )       \to  \mathrm {D}^+\mathrm {F}(S)     \]
    and \(\mathrm {R} a_*\) factors as
    \[       \mathrm {R} a_*\colon  \mathrm {D}^+(X_\bullet )       \to  \mathrm {D}^+\mathrm {F}(S).     \]
    The spectral sequence of the filtered complex \((\mathrm {R} a_* K,L)\in \mathrm {D}^+\mathrm {F}(S)\) is exactly \hyperref[hodge-theory-iii-5.2.7.1]{Equation~\ref*{hodge-theory-iii-5.2.7.1}}.\label{hodge-theory-iii-7.1.5}\par{}If \(K\) is filtered (resp. bifiltered), then \(\mathrm {R}{a_\bullet }_*K\) is filtered (resp. bifiltered): we have

    \begin{eqn}{7.1.5.1}\setnumber{7.1.5.1}\label{hodge-theory-iii-7.1.5.1}\skipsinglepar\[         \mathrm {R}{a_\bullet }_*         \colon  \mathrm {D}^+\mathrm {F}(X_\bullet )         \to  \mathrm {D}^+\mathrm {F}(S_\bullet )       \tag{7.1.5.1}       \]\end{eqn}

    \begin{eqn}{7.1.5.2}\setnumber{7.1.5.2}\label{hodge-theory-iii-7.1.5.2}\skipsinglepar\[         \mathrm {R}{a_\bullet }_*         \colon  \mathrm {D}^+\mathrm {F}_2(X_\bullet )         \to  \mathrm {D}^+\mathrm {F}_2(S_\bullet ).       \tag{7.1.5.2}       \]\end{eqn}\label{hodge-theory-iii-7.1.6}\par{}Let \(K\) be a complex of cosimplicial sheaves on \(S\), endowed with an increasing filtration \(W\).
    We define the \emph{diagonal filtration} \(\delta (W,L)\) of \(W\) and \(L\) to be the increasing filtration of \(\mathrm {s} K\) given by

    \begin{eqn}{7.1.6.1}\setnumber{7.1.6.1}\label{hodge-theory-iii-7.1.6.1}\skipsinglepar\[         \begin {aligned}           \delta (W,L)_n(\mathrm {s} K)           &= \bigoplus _{p,q} W_{n+p}(K^{p,q})         \\&= \sum _p \mathrm {s}(W_{n+p}(K)) \cap  L^p(\mathrm {s} K).         \end {aligned}       \tag{7.1.6.1}       \]\end{eqn}\par{}We have

    \begin{eqn}{7.1.6.2}\setnumber{7.1.6.2}\label{hodge-theory-iii-7.1.6.2}\skipsinglepar\[         \operatorname {Gr}_n^{\delta (W,L)}(\mathrm {s} K)         \simeq  \bigoplus _p \operatorname {Gr}_{n+p}^W(K^{\bullet  p})[-p].       \tag{7.1.6.2}       \]\end{eqn}\par{}The functor \((K,W)\mapsto (\mathrm {s} K,\delta (W,L))\) sends filtered quasi-isomorphisms to filtered quasi-isomorphisms, and defines

    \begin{eqn}{7.1.6.3}\setnumber{7.1.6.3}\label{hodge-theory-iii-7.1.6.3}\skipsinglepar\[         \begin {aligned}           (s,\delta )           \colon  \mathrm {D}^+\mathrm {F}(S_\bullet )           &\to  \mathrm {D}^+\mathrm {F}(S)         \\(K,W)           &\mapsto  (\mathrm {s} K,\delta (W,L))         \end {aligned}       \tag{7.1.6.3}       \]\end{eqn}

    whence, by composition with \hyperref[hodge-theory-iii-7.1.5.1]{Equation~\ref*{hodge-theory-iii-7.1.5.1}},

    \begin{eqn}{7.1.6.4}\setnumber{7.1.6.4}\label{hodge-theory-iii-7.1.6.4}\skipsinglepar\[         \begin {aligned}           (\mathrm {R}\Gamma ,\delta (-,L))           \colon  \mathrm {D}^+\mathrm {F}(X_\bullet )           &\to  \mathrm {D}^+\mathrm {F}(S)         \\(K,W)           &\mapsto  (\mathrm {R}\Gamma  K,\delta (W,L)).         \end {aligned}       \tag{7.1.6.4}       \]\end{eqn}

    From \hyperref[hodge-theory-iii-7.1.6.2]{Equation~\ref*{hodge-theory-iii-7.1.6.2}}, we see that

    \begin{eqn}{7.1.6.5}\setnumber{7.1.6.5}\label{hodge-theory-iii-7.1.6.5}\skipsinglepar\[         \operatorname {Gr}_n^{\delta (W,L)}         = \bigoplus _p \mathrm {R}{a_p}_*(\operatorname {Gr}_{n+p}^W K)[-p].       \tag{7.1.6.5}       \]\end{eqn}\label{hodge-theory-iii-7.1.7}\par{}If \((K,W,F)\) is a bifiltered complex of cosimplicial sheaves, then \(\mathrm {s} K\) is endowed with the three filtrations \(W\), \(F\), and \(L\), and defines different bifiltered complexes.\par{}For example, for increasing \(W\), the functor \((K,W,F)\mapsto (K,\delta (W,L),F)\) sends bifiltered quasi-isomorphisms to bifiltered quasi-isomorphisms, and thus defines

    \begin{eqn}{7.1.7.1}\setnumber{7.1.7.1}\label{hodge-theory-iii-7.1.7.1}\skipsinglepar\[         \begin {aligned}           \mathrm {D}^+\mathrm {F}_2(S_\bullet )           &\to  \mathrm {D}^+\mathrm {F}_2(S)         \\(K,W,F)           &\mapsto  (K,\delta (W,L),F).         \end {aligned}       \tag{7.1.7.1}       \]\end{eqn}

    By composition with \hyperref[hodge-theory-iii-7.1.5.2]{Equation~\ref*{hodge-theory-iii-7.1.5.2}}, we thus obtain

    \begin{eqn}{7.1.7.2}\setnumber{7.1.7.2}\label{hodge-theory-iii-7.1.7.2}\skipsinglepar\[         \begin {aligned}           \mathrm {D}^+\mathrm {F}_2(X_\bullet )           &\to  \mathrm {D}^+\mathrm {F}_2(S)         \\(K,W,F)           &\mapsto  (\mathrm {R}\Gamma  K,\delta (W,L),F).         \end {aligned}       \tag{7.1.7.2}       \]\end{eqn}

    and we have that

    \begin{eqn}{7.1.7.3}\setnumber{7.1.7.3}\label{hodge-theory-iii-7.1.7.3}\skipsinglepar\[         \operatorname {Gr}_n^{\delta (W,L)}(\mathrm {R}\Gamma  K,F)         = \bigoplus _p \mathrm {R}{a_p}_*(\operatorname {Gr}_{n+p}^W K,F)[-p]       \tag{7.1.7.3}       \]\end{eqn}

    in \(\mathrm {D}^+\mathrm {F}(S)\).\section{Supplements to the two filtrations lemma}\label{hodge-theory-iii-7.2}\par{}In this section, we give a new proof of the two filtrations lemma (\hyperref[hodge-theory-ii-1.3.16]{Theorem~\ref*{hodge-theory-ii-1.3.16}}) and some supplements.\label{hodge-theory-iii-7.2.1}\par{}Let \((K,F,W)\) be a bounded-below bifiltered complex with objects in an abelian category \(\mathscr {A}\), with \(F\) assumed to be biregular.\par{}We say that \((K,F,W)\) is \emph{\(F\)-splitable} if the filtered complex \((K,W)\) can be written as a sum of filtered complexes
    \[       (K,W)       = \bigoplus _{n\in \mathbb {Z}} (K_n,W_n)     \]
    with
    \[       F^nK       = \bigoplus _{n'\geqslant  n} K_{n'}.     \]\par{}Let \(r_0\geqslant 0\) be an integer or \(+\infty \).
  The following condition was considered in \hyperref[hodge-theory-ii-1.3.16]{Theorem~\ref*{hodge-theory-ii-1.3.16}} and \hyperref[hodge-theory-ii-1.3.17]{Corollary~\ref*{hodge-theory-ii-1.3.17}}:\begin{condition}{7.2.2}\setnumber{7.2.2}\label{hodge-theory-iii-7.2.2}\skipsinglepar\par{}For every non-negative integer \(r<r_0\), the differentials \(d_r\) of the graded complex \(E_r(K,W)\) are strictly compatible with the recurrent filtration defined by \(F\).\end{condition}\label{hodge-theory-iii-7.2.3}\par{}It is clear that, if \((K,F,W)\) is \(F\)-splitable, then \hyperref[hodge-theory-iii-7.2.2]{Condition~\ref*{hodge-theory-iii-7.2.2}} is satisfied for \(r_0=\infty \).
    Conversely, it seems that if \hyperref[hodge-theory-iii-7.2.2]{Condition~\ref*{hodge-theory-iii-7.2.2}} is satisfied for \(r_0=\infty \), then everything is as if the functor \(\operatorname {Gr}_F\) were exact.
    For example, we will show that the \(\operatorname {Gr}_F\) of the spectral sequence \(E(K,W)\) can then be identified with the spectral sequence \(E(\operatorname {Gr}_FK,W)\), and that the spectral sequence \(E(K,F)\) degenerates (\(E_1=E_\infty \)).\label{hodge-theory-iii-7.2.4}\par{}We immediately deduce from the definition (\cref{hodge-theory-ii-1.3.8}) that the first direct filtration of \(E_r(K,W)\) is the filtration \(F_d\) of \(E_r(K,W)\) by the images
    \[       F_d^p(E_r(K,W))       = \operatorname {Im} \big (E_r(F^pK,W) \to  E_r(K,W)\big ).     \]
    Dually, the second direct filtration (\cref{hodge-theory-ii-1.3.9}) of \(E_r(K,W)\) is the filtration \(F_{d^*}\) of \(E_r(K,W)\) by the kernels
    \[       F_{d^*}^p(E_r(K,W))       = \operatorname {Ker} \big (E_r(K,W) \to  E_r(K/F^pK,W)\big ).     \]
    The recurrent filtration \(F_{\mathrm {rec}}\) of \(E_r(K,W)\) is intermediary between these two filtrations ((iii) of \hyperref[hodge-theory-ii-1.3.13]{Proposition~\ref*{hodge-theory-ii-1.3.13}}).\begin{proposition}{7.2.5}\setnumber{7.2.5}\label{hodge-theory-iii-7.2.5}\skipsinglepar\par{}Suppose that \((K,F,W)\) satisfies \hyperref[hodge-theory-iii-7.2.2]{Condition~\ref*{hodge-theory-iii-7.2.2}} for some \(r_0\geqslant 0\).
    Then

    
 \begin{enumerate}
      
 \item[(i)]
        \((F^aK/F^bK,F,W)\) also satisfies \hyperref[hodge-theory-iii-7.2.2]{Condition~\ref*{hodge-theory-iii-7.2.2}} for \(r_0\).
      

      
 \item[(ii)]
        For \(r\leqslant  r_0\), the sequence
        \[           0           \to  E_r(F^pK,W)           \to  E_r(K,W)           \to  E_r(K/F^pK,W)           \to  0         \]
        is exact;
        for \(r=r_0+1\), the sequence
        \[           E_r(F^pK,W)           \to  E_r(K,W)           \to  E_r(K/F^pK,W)         \]
        is exact.
        In particular, for \(r\leqslant  r_0+1\), the two direct filtrations and the recurrent filtration of \(E_r(K,W)\) agree.
      

    \end{enumerate}\begin{proof}\par{}Fix an integer \(p\).
      We will prove the claim

      \label{hodge-theory-iii.7.2.5-inductive-claim}\par{}\emph{\((*)_r\) If \(r\leqslant  r_0\), then \(E_r(F^pK,W)\) injects into \(E_r(K,W)\);
            if \(r\leqslant  r_0+1\), then its image is \(F_\mathrm {rec}(E_r(K,W))\).}

      by induction on \(r\).\par{}The claim \((*)_0\) is always true.
      Assume \((*)_r\), and let us prove \((*)_{r+1}\).
      We can suppose that \(0\leqslant  r\leqslant  r_0\).
      We have a diagram
      
      
  \begin {tikzcd}
        E_r(K,W)
          \ar [r,"d_r"]
        & E_r(K,W)
          \ar [r,"d_r"]
        & E_r(K,W)
      \\E_r(F^pK,W)
          \ar [r]
          \ar [u,hook]
        & E_r(F^pK,W)
          \ar [r]
          \ar [u,hook]
        & E_r(F^pK,W)
          \ar [u,hook,swap,"i"]
      \end {tikzcd}


      
      and the image of the vertical inclusions is \(F_\mathrm {rec}^p(E_r(K,W))\).
      Since \(i\) is injective,

      \[         \begin {aligned}           F_\mathrm {rec}^p(E_{r+1}(K,W))           &= \operatorname {Im}\Big ( \operatorname {Ker} \big (F_\mathrm {rec}^pE_r(K,W)\xrightarrow {d_r}E_r(K,W) \big ) \to  E_{r+1}(K,W) \Big )         \\&= \operatorname {Im}\Big ( \operatorname {Ker}\big ( E_r(F^pK,W)\xrightarrow {d_r}E_r(F^pK,W) \big ) \to  E_{r+1}(K,W) \Big )         \\&= \operatorname {Im}\big ( E_{r+1}(F^pK,W) \to  E_{r+1}(K,W) \big ).         \end {aligned}       \]

      If \(r<r_0\), then \(d_r\) is strictly compatible with \(F_\mathrm {rec}\), whence
      \[         d_rE_r(K,W)\cap  E_r(F^pK,W)         = d_rE_r(F^pK,W)       \]
      and \(E_{r+1}(F^pK,W)\) thus injects into \(E_{r+1}(K,W)\).
      This proves \((*)_{r+1}\).\par{}The statements \((*)_r\), combined with the dual statements, implies (ii).\par{}It follows from (ii) that, for \(r<r_0\) and \(a\leqslant  b\), \(E_r(F^bK,W)\) injects into \(E_r(F^aK,W)\), and that the differential \(d_r\) of \(E_r(F^aK,W)\) is strictly compatible with the first direct filtration of \(E_r(F^aK,W)\).
      We thus deduce, by induction on \(r<r_0\), that the first direct filtration of \(E_r(F^aK,W)\) coincides with the recurrent filtration, and (i) then follows.\end{proof}\end{proposition}\label{hodge-theory-iii-7.2.6}\par{}If \hyperref[hodge-theory-iii-7.2.2]{Condition~\ref*{hodge-theory-iii-7.2.2}} is satisfied for \(r_0\), and \(r\leqslant  r_0+1\), then we denote by \(F\) the filtration \(F_d=F_{d^*}=F_\mathrm {rec}\) of \(E_r(W,K)\).
    The exact sequences from (ii) of \hyperref[hodge-theory-iii-7.2.5]{Proposition~\ref*{hodge-theory-iii-7.2.5}}

    
  \begin {tikzcd}
      0
      & E_r(K/F^pK,W)
        \ar [l]
      && E_r(K/F^{p+1}K,W)
        \ar [ll]
      && 0
    \\&& E_r(K,W)
        \ar [ul]
        \ar [ur]
      && E_r(\operatorname {Gr}_F^pK,W)
        \ar [ul]
        \ar [ur]
    \\0
        \ar [r]
      & E_r(F^{p+1}K,W)
        \ar [rr]
        \ar [ur]
      && E_r(F^pK,W)
        \ar [ul]
        \ar [ur]
      && 0
        \ar [ul]
    \end {tikzcd}


    define, for \(r\leqslant  r_0\), an (autodual) isomorphism that is compatible with the differentials \(d_r\), namely

    \begin{eqn}{7.2.6.1}\setnumber{7.2.6.1}\label{hodge-theory-iii-7.2.6.1}\skipsinglepar\[         \operatorname {Gr}_F^p(E_r(K,W))         \simeq  E_r(\operatorname {Gr}_F^pK,W)         \qquad \text {for }r\leqslant  r_0.       \tag{7.2.6.1}       \]\end{eqn}\label{hodge-theory-iii-7.2.7}\par{}If \hyperref[hodge-theory-iii-7.2.2]{Condition~\ref*{hodge-theory-iii-7.2.2}} is satisfied for \(r_0=\infty \), then the sequences
    \[       0       \to  E_r(F^pK,W)       \to  E_r(K,W)       \to  E_r(K/F^pK,W)       \to  0     \]
    are exact for all \(r\).
    If the filtration \(W\) is biregular, then the sequence
    \[       0       \to  E_\infty (F^pK,W)       \to  E_\infty (K,W)       \to  E_\infty (K/F^pK,W)       \to  0     \]
    is thus exact.
    We can rewrite this sequence as
    \[       0       \to  \operatorname {Gr}_W(\operatorname {H}(F^pK))       \to  \operatorname {Gr}_W(\operatorname {H}(K))       \to  \operatorname {Gr}_W(\operatorname {H}(K/F^pK))       \to  0.     \]\begin{lemma}{7.2.7.1}\setnumber{7.2.7.1}\label{hodge-theory-iii-7.2.7.1}\skipsinglepar\par{}Let \(L\) be a complex endowed with a biregular filtration \(G\).
      For \(\operatorname {Gr}_G(L)\) to be acyclic, it is necessary and sufficient that \(L\) be acyclic and that the differentials of \(L\) be strictly compatible with the filtration \(G\).\begin{proof}\par{}This is a particular case of \hyperref[hodge-theory-ii-1.3.2]{Proposition~\ref*{hodge-theory-ii-1.3.2}}.\end{proof}\end{lemma}\par{}Applying \hyperref[hodge-theory-iii-7.2.7.1]{Lemma~\ref*{hodge-theory-iii-7.2.7.1}}, taking \(L\) to be the complex
    \[       0       \to  \operatorname {H}(F^pK)       \to  \operatorname {H}(K)       \to  \operatorname {H}(K/F^pK)       \to  0,     \]
    we see that:

    
 \begin{enumerate}
      
 \item[(a)]
        This sequence is exact, i.e. the differentials of \(K\) are strictly compatible with the filtration \(F\);
      


      
 \item[(b)]
        the filtration of \(\operatorname {H}(F^pK)\) induced by the filtration \(W\) of \(F^pK\) coincides with the filtration induced by the filtration \(W\) of \(\operatorname {H}(K)\);
      


      
 \item[(c)]
        an analogous statement holds for \(K/F^pK\).
      

    \end{enumerate}\par{}In conclusion:\begin{proposition}{7.2.8}\setnumber{7.2.8}\label{hodge-theory-iii-7.2.8}\skipsinglepar\par{}If \((K,F,W)\) satisfies \hyperref[hodge-theory-iii-7.2.2]{Condition~\ref*{hodge-theory-iii-7.2.2}} for \(r_0=\infty \), then the spectral sequence \(E(K,F)\) degenerates (\(E_1=E_\infty \)).
    Furthermore, we have an isomorphism of spectral sequences
    \[       \operatorname {Gr}_F^p(E_\bullet (K,W))       \simeq  E_\bullet (\operatorname {Gr}_F^pK,W).     \]\end{proposition}\chapter{Hodge theory of algebraic spaces}\thispagestyle{fancy}\label{hodge-theory-iii-8}\section{Hodge complexes}\label{hodge-theory-iii-8.1}\label{hodge-theory-iii-8.1.1}\par{}Let \(A\) be as in \cref{hodge-theory-iii-0.3}.
    Then a \emph{Hodge \(A\)-complex} \(K\) of weight \(n\) consists of

    
 \begin{enumerate}
      
 \item[(𝛼)]
        a complex \(K_A\in \operatorname {Ob}\mathrm {D}^+(A)\) such that \(\operatorname {H}^k(K_A)\) is an \(A\)-module of finite type for all \(k\);
      


      
 \item[(𝛽)]
        a filtration \(F\) on \(K_A\otimes \mathbb {C}\), i.e. a filtered complex \((K_\mathbb {C},F)\in \operatorname {Ob}\mathrm {D}^+F(\mathbb {C})\), and an isomorphism \(\alpha \colon  K_\mathbb {C}\simeq  K_A\otimes \mathbb {C}\) in \(\mathrm {D}^+(\mathbb {C})\).
      

    \end{enumerate}\par{}The following axioms have to be satisfied:

    \begin{axiom}{CH1}\setnumber{CH1}\label{hodge-theory-iii-8.1.1-ch1}\skipsinglepar\par{}The differential \(d\) of \(K_\mathbb {C}\) is strictly compatible with the filtration \(F\);
        in other words, the spectral sequence defined by \((K_\mathbb {C},F)\) degenerates at \(E_1\) (i.e. \(E_1=E_\infty \)) \hyperref[hodge-theory-ii-1.3.2]{Proposition~\ref*{hodge-theory-ii-1.3.2}}.\end{axiom}

    \begin{axiom}{CH2}\setnumber{CH2}\label{hodge-theory-iii-8.1.1-ch2}\skipsinglepar\par{}For all \(k\), the filtration \(F\) on \(\operatorname {H}^k(K_\mathbb {C})\simeq \operatorname {H}^k(K_A)\otimes \mathbb {C}\) defines a Hodge \(A\)-structure of weight \(n+k\) on \(\operatorname {H}^k(K_A)\), i.e. the filtration \(F\) is \((n+k)\)-opposite to its complex conjugate (which makes sense, since \(A\subset \mathbb {R}\)).\end{axiom}\label{hodge-theory-iii-8.1.2}\par{}Let \(X\) be a topological space.
    Then a \emph{cohomological Hodge \(A\)-complex \(K\) of weight \(n\)} on \(X\) consists of:

    
 \begin{enumerate}
      
 \item[(𝛼)]
        a complex \(K_A\in \operatorname {Ob}\mathrm {D}^+(X,A)\),
      

      
 \item[(𝛽)]
        a filtered complex \((K_\mathbb {C},F)\in \operatorname {Ob}\mathrm {D}^+F(X,\mathbb {C})\),
      

      
 \item[(𝛾)]
        an isomorphism \(\alpha \colon  K_\mathbb {C}\simeq  K_A\otimes \mathbb {C}\) in \(\mathrm {D}^+(X,\mathbb {C})\).
      

    \end{enumerate}\par{}The following axiom has to be satisfied:

    \begin{axiom}{CHC}\setnumber{CHC}\label{hodge-theory-iii-8.1.2-chc}\skipsinglepar\par{}The triple \((\mathrm {R}\Gamma (K_A),\mathrm {R}\Gamma ((K_\mathbb {C},F)),\mathrm {R}\Gamma (\alpha ))\) is a Hodge complex of weight \(n\).\end{axiom}\par{}When \(A=\mathbb {Z}\), we simply speak of Hodge complexes and cohomological Hodge complexes.\par{}The following statement reformulates a part of classical Hodge theory (cf. \cref{hodge-theory-ii-2.2.1}).\begin{scholium}{8.1.3}\setnumber{8.1.3}\label{hodge-theory-iii-8.1.3}\skipsinglepar\par{}Let \(X\) be a complex Kähler variety.
    Let \(K_\mathbb {Z}\) be the complex \(\mathbb {Z}[0]\) consisting of the constant sheaf \(\mathbb {Z}\) in degree \(0\), and let \(K_\mathbb {C}\) be the holomorphic de Rham complex \(\Omega _X^\bullet \), with \(F\) its stupid filtration.
    We have \(\alpha \colon  K_\mathbb {Z}\otimes \mathbb {C}\simeq \mathbb {C}[0]\xrightarrow {\approx }\Omega _X^\bullet \) (the Poincaré lemma), and \((K_\mathbb {Z},(K_\mathbb {C},F),\alpha )\) is a cohomological Hodge complex of weight \(0\).\end{scholium}\begin{remark}{8.1.4}\setnumber{8.1.4}\label{hodge-theory-iii-8.1.4}\skipsinglepar\par{}If \(K\) is a Hodge complex (resp. cohomological Hodge complex) of weight \(n\), then \(K[m]\) is a Hodge complex (resp. cohomological Hodge complex) of weight \(n+m\).\end{remark}\label{hodge-theory-iii-8.1.5}\par{}A \emph{mixed Hodge \(A\)-complex} \(K\) consists of:

    
 \begin{enumerate}
      
 \item[(𝛼)]
        a complex \(K_A\in \operatorname {Ob}\mathrm {D}^+(A)\) such that \(\operatorname {H}^k(K_A)\) is an \(A\)-module of finite type for all \(k\);
      


      
 \item[(𝛽)]
        a filtered complex \((K_{A\otimes \mathbb {Q}},W)\in \operatorname {Ob}\mathrm {D}^+F(A\otimes \mathbb {Q})\) (for an \emph{increasing} filtration \(W\)) and an isomorphism \(K_{A\otimes \mathbb {Q}}\simeq  K_A\otimes \mathbb {Q}\) in \(\mathrm {D}^+(A\otimes \mathbb {Q})\);
      


      
 \item[(𝛾)]
        a bifiltered complex \((K_\mathbb {C},W,F)\in \operatorname {Ob}\mathrm {D}^+F_w(\mathbb {C})\) (for an \emph{increasing} filtration \(W\) and a \emph{decreasing} filtration \(F\)) and an isomorphism \(\alpha \colon (K_\mathbb {C},W)\simeq (K_{A\otimes \mathbb {Q}},W)\otimes \mathbb {C}\) in \(\mathrm {D}^+F(\mathbb {C})\).
      

    \end{enumerate}\par{}The following axiom must be satisfied:

    \begin{axiom}{CHM}\setnumber{CHM}\label{hodge-theory-iii-8.1.5-chm}\skipsinglepar\par{}For all \(n\), the system consisting of the complex \(\operatorname {Gr}_n^W(K_{A\otimes \mathbb {Q}})\in \operatorname {Ob}\mathrm {D}^+(A\otimes \mathbb {Q})\), the complex \(\operatorname {Gr}_n^W(K_\mathbb {C})\in \operatorname {Ob}\mathrm {D}^+F(\mathbb {C})\) filtered by \(F\), and the isomorphism
        \[           \operatorname {Gr}_n^W(\alpha )\colon            \operatorname {Gr}_n^W(K_{A\otimes \mathbb {Q}})\otimes \mathbb {C}           \simeq  \operatorname {Gr}_n^W(K_\mathbb {C}),         \]
        is a Hodge \((A\otimes \mathbb {Q})\)-complex of weight \(n\).\end{axiom}\label{hodge-theory-iii-8.1.6}\par{}A \emph{cohomological mixed Hodge \(A\)-complex} \(K\) on a topological space \(X\) consists of:

    
 \begin{enumerate}
      
 \item[(𝛼)]
        a complex \(K_A\in \operatorname {Ob}\mathrm {D}^+(X,A)\) such that \(\operatorname {\mathbb {H}}^k(X,K_A)\) is an \(A\)-module of finite type for all \(k\), and such that \(\operatorname {\mathbb {H}}(X,K_A)\otimes \mathbb {Q}\xrightarrow {\sim }\operatorname {\mathbb {H}}(X,K_A\otimes \mathbb {Q})\);
      


      
 \item[(𝛽)]
        a filtered complex \((K_{A\otimes \mathbb {Q}},W)\in \operatorname {Ob}\mathrm {D}^+F(X,A\otimes \mathbb {Q})\) (for an \emph{increasing} filtration \(W\)) and an isomorphism \(K_{A\otimes \mathbb {Q}}\simeq  K_A\otimes \mathbb {Q}\) in \(\mathrm {D}^+(X,A\otimes \mathbb {Q})\);
      


      
 \item[(𝛾)]
        a bifiltered complex \((K_\mathbb {C},W,F)\in \operatorname {Ob}\mathrm {D}^+F_w(X,\mathbb {C})\) (for an \emph{increasing} filtration \(W\) and a \emph{decreasing} filtration \(F\)) and an isomorphism \(\alpha \colon (K_\mathbb {C},W)\simeq (K_{A\otimes \mathbb {Q}},W)\otimes \mathbb {C}\) in \(\mathrm {D}^+F(X,\mathbb {C})\).
      

    \end{enumerate}\par{}The following axiom must be satisfied:

    \begin{axiom}{CHMC}\setnumber{CHMC}\label{hodge-theory-iii-8.1.6-chmc}\skipsinglepar\par{}For all \(n\), the system consisting of the complex \(\operatorname {Gr}_n^W(K_{A\otimes \mathbb {Q}})\in \operatorname {Ob}\mathrm {D}^+(X,A\otimes \mathbb {Q})\), the complex \(\operatorname {Gr}_n^W(K_\mathbb {C})\in \operatorname {Ob}\mathrm {D}^+F(X,\mathbb {C})\) filtered by \(F\), and the isomorphism
        \[           \operatorname {Gr}_n^W(\alpha )\colon            \operatorname {Gr}_n^W(K_{A\otimes \mathbb {Q}})\otimes \mathbb {C}           \simeq  \operatorname {Gr}_n^W(K_\mathbb {C}),         \]
        is a cohomological Hodge \((A\otimes \mathbb {Q})\)-complex of weight \(n\).\end{axiom}\par{}We can trivially show:\begin{proposition}{8.1.7}\setnumber{8.1.7}\label{hodge-theory-iii-8.1.7}\skipsinglepar\par{}If \(K=(K_A,(K_{A\otimes \mathbb {Q}},W),(K_\mathbb {C},W,F))\) is a cohomological mixed Hodge \(A\)-complex, then \(\mathrm {R}\Gamma  K=(\mathrm {R}\Gamma  K_A,\mathrm {R}\Gamma (K_{A\otimes \mathbb {Q}},W),\mathrm {R}\Gamma (K_\mathbb {C},W,F))\) is a mixed Hodge \(A\)-complex.\end{proposition}\begin{scholium}{8.1.8}\setnumber{8.1.8}\label{hodge-theory-iii-8.1.8}\skipsinglepar\par{}Let \(U\) be the complement in a smooth proper scheme \(X\), of a normal crossing divisor \(D\), with \(j\colon  U\hookrightarrow  X\) the inclusion morphism.
    Let \(W\) be the canonical filtration on \(\mathrm {R} j_*\mathbb {Q}=\mathrm {R} j_*\mathbb {Z}\otimes \mathbb {Q}\colon  W_n(\mathrm {R} j_*\mathbb {Q})=\tau _{\leqslant  n}(\mathrm {R} j_*\mathbb {Q})\).
    We have (\hyperref[hodge-theory-ii-3.1.8]{Proposition~\ref*{hodge-theory-ii-3.1.8}})
    \[       (\mathrm {R} j_*\mathbb {Q})\otimes \mathbb {C}       \xrightarrow {\approx } \mathrm {R} j_*\mathbb {C}       \xrightarrow {\approx } j_*\Omega _U^\bullet        \xleftarrow {\approx } \Omega _X^\bullet (\log  D).     \]
    Let \(W\) be the weight filtration of \(\Omega _+X^\bullet (\log  D)\), and \(F\) its Hodge filtration (by the \(\sigma _{\geqslant  p}\)).
    Then \hyperref[hodge-theory-ii-3.1.8]{Proposition~\ref*{hodge-theory-ii-3.1.8}} gives an isomorphism \((\mathrm {R} j_*\mathbb {Q},W)\otimes \mathbb {C}\simeq (\Omega _X^\bullet (\log  D),W)\) in \(\mathrm {D}^+F(X,\mathbb {C})\), and, by \cref{hodge-theory-ii-3.1.5} and \hyperref[hodge-theory-ii-3.1.9]{Proposition~\ref*{hodge-theory-ii-3.1.9}},
    \[       (\mathrm {R} j_*\mathbb {Z},(\mathrm {R} j_*\mathbb {Q},W),(\Omega _X^\bullet (\log  D), W, F))     \]
    is a cohomological mixed Hodge complex on \(X\).\end{scholium}\par{}The following result is an "abstract" version of \cref{hodge-theory-ii-3.2}.\begin{scholium}{8.1.9}\setnumber{8.1.9}\label{hodge-theory-iii-8.1.9}\skipsinglepar\par{}Let \(K\) be a mixed Hodge \(A\)-complex.\begin{enumerate}
    
 \item[(i)]
      On the pages \({}_WE_r^{pq}\) of the spectral sequence of \((K_\mathbb {C},W)\), the recurrent filtration and the two direct filtrations defined by \(F\) all coincide.
    


    
 \item[(ii)]
      The filtration \(W[n]\) (\hyperref[hodge-theory-ii-1.1.2]{Definition~\ref*{hodge-theory-ii-1.1.2}}) on \(\operatorname {H}^n(K_{A\otimes \mathbb {Q}})\simeq \operatorname {H}^n(K_A)\otimes \mathbb {Q}\) and the filtration \(F\) on \(\operatorname {H}^n(K_\mathbb {C})\simeq \operatorname {H}^n(K_A)\otimes _A\mathbb {C}\) define a mixed Hodge \(A\)-structure on \(\operatorname {H}^n(K_A)\) (\hyperref[hodge-theory-ii-2.3.1]{Definition~\ref*{hodge-theory-ii-2.3.1}} along with \cref{hodge-theory-iii-0.4}).
    


    
 \item[(iii)]
      The morphisms \({}_Wd_1\colon {}_WE_1^{pq}\to {}_WE_1^{p+1,q}\) are compatible with the Hodge bigrading;
      in particular, they are strictly compatible with the Hodge filtration.
    


    
 \item[(iv)]
      The spectral sequence of \((K_{A\otimes \mathbb {Q}},W)\) degenerates at \(E_2\) (i.e. \({}_WE_2={}_WE_\infty \)).
    


    
 \item[(v)]
      The spectral sequence of \((K_\mathbb {C},F)\) degenerates at \(E_1\) (i.e. \({}_FE_2={}_FE_\infty \)).
    


    
 \item[(vi)]
      The spectral sequence of the complex \(\operatorname {Gr}_F^p(K_\mathbb {C})\), endowed with the filtration induced by \(W\), degenerates at \(E_2\).
    

  \end{enumerate}\begin{proof}\par{}The proof of (i), (ii), (iii), (iv) is parallel to that of \hyperref[hodge-theory-ii-3.2.5]{Theorem~\ref*{hodge-theory-ii-3.2.5}}.
      The analogy of \hyperref[hodge-theory-ii-3.2.6]{Lemma~\ref*{hodge-theory-ii-3.2.6}} and \hyperref[hodge-theory-ii-3.2.7]{Lemma~\ref*{hodge-theory-ii-3.2.7}} has been taken as an axiom (\hyperref[hodge-theory-iii-8.1.1-ch1]{Axiom~\ref*{hodge-theory-iii-8.1.1-ch1}} and \hyperref[hodge-theory-iii-8.1.1-ch2]{Axiom~\ref*{hodge-theory-iii-8.1.1-ch2}}).
      The proofs of \hyperref[hodge-theory-ii-3.2.8]{Lemma~\ref*{hodge-theory-ii-3.2.8}}, \hyperref[hodge-theory-ii-3.2.9]{Lemma~\ref*{hodge-theory-ii-3.2.9}}, and \hyperref[hodge-theory-ii-3.2.10]{Lemma~\ref*{hodge-theory-ii-3.2.10}} can be carried over as they are, and, as in \hyperref[hodge-theory-ii-3.2.10]{Lemma~\ref*{hodge-theory-ii-3.2.10}}, we thus deduce (i), (ii), (iii), and (iv).
      Claims (v) and (vi) then follow from \hyperref[hodge-theory-iii-7.2.8]{Proposition~\ref*{hodge-theory-iii-7.2.8}}.\end{proof}\end{scholium}\label{hodge-theory-iii-8.1.10}\par{}We abbreviate "differential graded" to DG, and "bounded-below differential graded" to DG+.
    A DG complex of \(A\)-modules can be seen as a double complex of \(A\)-modules;
    the \emph{first} degree is that of the complex, and the \emph{second} degree is that defined by the grading of DG modules.
    
    We denote by \(\mathrm {D}^+(A\mathsf {-mod}_{\mathsf {DG}^+})\) the derived category of the category of bounded-below complexes of DG \(A\)-modules, with degrees uniformly bounded below.\par{}A \emph{DG mixed Hodge \(A\)-complex} \(K\) consists of:

    
 \begin{enumerate}
      
 \item[(a)]
        a complex \(K_A\in \operatorname {Ob}\mathrm {D}^+(A\mathsf {-mod}_{\mathsf {DG}^+})\);
      

      
 \item[(b)]
        a filtered complex \((K_{A\otimes \mathbb {Q}},W)\in \operatorname {Ob}\mathrm {D}^+F((A\otimes \mathbb {Q})\mathsf {-mod}_{\mathsf {DG}^+})\), and is isomorphism \(K_A\otimes \mathbb {Q}\simeq  K_{A\otimes \mathbb {Q}}\) in \(\mathrm {D}^+\);
      

      
 \item[(c)]
        a bifiltered complex \((K_\mathbb {C},F,W)\in \operatorname {Ob}\mathrm {D}^+F_2(\mathbb {C}\mathsf {-mod}_{\mathsf {DG}^+})\), and an isomorphism \(K_{A\otimes \mathbb {Q}}\otimes \mathbb {C}\simeq  K_\mathbb {C}\) in \(\mathrm {D}^+F\).
      

    \end{enumerate}\par{}We require that, for each \(n\), the component \((K^{\bullet  n},F,W)\) of \(K\) whose second degree is \(n\) be a mixed Hodge \(A\)-complex.\par{}We denote by \(L\) the filtration by the second degree of \(\mathrm {s} K\).
    We denote by \(\operatorname {Dec}_1W\) the filtration of \(K_{A\otimes \mathbb {Q}}\) induced by shifting (\cref{hodge-theory-ii-1.3.3}) \(W\).
    We will not confuse the filtration that this induces on \(\mathrm {s} K_{A\otimes \mathbb {Q}}\) (also denoted by \(\operatorname {Dec}_1W\)) with the shifted filtration of the filtration induced on \(\mathrm {s} K_{A\otimes \mathbb {Q}}\) by \(W\).\begin{variant}{8.1.10.1}\setnumber{8.1.10.1}\label{hodge-theory-iii-8.1.10.1}\skipsinglepar\par{}We similarly define \emph{cosimplicial} (resp. \emph{\(r\)-cosimplicial}) \emph{mixed Hodge \(A\)-complexes} by replacing "DG" with "cosimplicial" (resp. \(r\)-cosimplicial) everywhere in the definition.
      The functor from cosimplicial \(A\)-modules to DG \(A\)-modules sends cosimplicial mixed Hodge \(A\)-complexes to DG mixed Hodge \(A\)-complexes (and similarly for \(r\)-cosimplicial).\end{variant}\label{hodge-theory-iii-8.1.11}\par{}Let \(X_\bullet \) be a simplicial topological space.
    A \emph{cohomological mixed Hodge \(A\)-complex} \(K\) on \(X_\bullet \) consists of:

    
 \begin{enumerate}
      
 \item[(𝛼)]
        a complex \(K_A\in \operatorname {Ob}\mathrm {D}^+(X_\bullet ,A)\);
      


      
 \item[(𝛽)]
        a filtered complex \((K_{A\otimes \mathbb {Q}},W)\in \operatorname {Ob}\mathrm {D}^+F(X_\bullet ,A\otimes \mathbb {Q})\), and an isomorphism \(K_{A\otimes \mathbb {Q}}\simeq  K_A\otimes \mathbb {Q}\) in \(\mathrm {D}^+(X_\bullet ,A\otimes \mathbb {Q})\);
      


      
 \item[(𝛾)]
        a bifiltered complex \((K_\mathbb {C},W,F)\in \operatorname {Ob}\mathrm {D}^+F_2(X_\bullet ,\mathbb {C})\), and an isomorphism \((K_\mathbb {C},W)\simeq (K_{A\otimes \mathbb {Q}},W)\otimes \mathbb {C}\) in \(\mathrm {D}^+F(X_\bullet ,\mathbb {C})\).
      

    \end{enumerate}\par{}The following axiom has to be satisfied:

    \begin{axiom}{CHM●}\setnumber{CHM●}\label{hodge-theory-iii-8.1.11-chmbullet}\skipsinglepar\par{}The restriction of \(K\) to each of the \(X_n\) is a cohomological mixed Hodge \(A\)-complex.\end{axiom}\par{}For \(A=\mathbb {Z}\), we simply speak of a cohomological mixed Hodge complex.\begin{example}{8.1.12}\setnumber{8.1.12}\label{hodge-theory-iii-8.1.12}\skipsinglepar\par{}Let \(X_\bullet \) be a smooth proper scheme (over \(\mathbb {C}\)), and let \(Y_\bullet \) be a normal crossing divisor in \(X_\bullet \).
    We set \(U_\bullet =X_\bullet \setminus  Y_\bullet \), and we denote by \(j\) the inclusion of \(U_\bullet \) into \(X_\bullet \).
    The constructions in \hyperref[hodge-theory-iii-8.1.8]{Scholium~\ref*{hodge-theory-iii-8.1.8}} can be performed "uniformly in \(X_n\)", as \cref{hodge-theory-iii-5.2.5} and \hyperref[hodge-theory-iii-6.2.7]{Lemma~\ref*{hodge-theory-iii-6.2.7}} show.
    
    They thus give a cohomological mixed Hodge complex
    \[       \big (\mathrm {R} j_*\mathbb {Z}, (\mathrm {R} j_*\mathbb {Q},\tau _{\leqslant }), (\Omega _{X_\bullet }^\bullet (\log  Y_\bullet ),W,F)\big )     \]
    on \(X_\bullet \).\end{example}\label{hodge-theory-iii-8.1.13}\par{}Let \(K\) be a cohomological mixed Hodge \(A\)-complex on \(X_\bullet \).
    Write \(A\mathsf {-mod}_\mathsf {cosimp}\) to mean the category of cosimplicial \(A\)-modules.
    If we apply the functor \(\mathrm {R}\Gamma ^\bullet \) to \(K\), then we obtain

    
 \begin{enumerate}
      
 \item[(a)]
        a cosimplicial complex \(\mathrm {R}\Gamma ^\bullet  K_A\in \operatorname {Ob}\mathrm {D}^+(A\mathsf {-mod}_\mathsf {cosimp})\);
      


      
 \item[(b)]
        a cosimplicial filtered complex \(\mathrm {R}\Gamma ^\bullet (K_{A\otimes \mathbb {Q}},W)\in \operatorname {Ob}\mathrm {D}^+F(A\otimes \mathbb {Q}\mathsf {-mod}_\mathsf {cosimp})\);
      


      
 \item[(c)]
        a cosimplicial bifiltered complex \(\mathrm {R}\Gamma ^\bullet (K_\mathbb {C},W,F)\in \operatorname {Ob}\mathrm {D}^+F_2(\mathbb {C}\mathsf {-mod}_\mathsf {cosimp})\);
      


      
 \item[(d)]
        isomorphisms as in (b) and (c) of \cref{hodge-theory-iii-8.1.10} between these objects.
      

    \end{enumerate}\par{}It is clear that the \(\mathrm {R}\Gamma ^\bullet (K,W,F)\) described above is a cosimplicial mixed Hodge \(A\)-module.
    It defines a DG mixed Hodge \(A\)-complex (\hyperref[hodge-theory-iii-8.1.10.1]{Variant~\ref*{hodge-theory-iii-8.1.10.1}}).\label{hodge-theory-iii-8.1.14}\par{}Let \(K\) be a DG mixed Hodge \(A\)-complex, and consider the spectral sequence abutting to the cohomology of \(\mathrm {s} K\), which is defined by the filtration \(L\).
    By hypothesis, the \(E_1\) pages of this spectral sequence are endowed with mixed Hodge \(A\)-structures.
    The content of \hyperref[hodge-theory-iii-8.1.15]{Theorem~\ref*{hodge-theory-iii-8.1.15}} is that every spectral sequence, including its abutment, is endowed with a mixed Hodge \(A\)-structure.
    In other words, \(c'\) is a spectral sequence in the abelian category of mixed Hodge \(A\)-structures.
    The structure on the abutment corresponds to a natural mixed Hodge \(A\)-complex "structure" on \(\mathrm {s} K\).\begin{theorem}{8.1.15}\setnumber{8.1.15}\label{hodge-theory-iii-8.1.15}\skipsinglepar\par{}Let \(K\) be a DG mixed Hodge \(A\)-complex.

    
 \begin{enumerate}
      
 \item[(i)]
        With the notation of \cref{hodge-theory-iii-8.1.10}, \hyperref[hodge-theory-iii-8.1.10.1]{Variant~\ref*{hodge-theory-iii-8.1.10.1}}, and \cref{hodge-theory-iii-7.1.6}, \((\mathrm {s} K,\delta (W,L),F)\) is a mixed Hodge \(A\)-complex.
      


      
 \item[(ii)]
        We have

        \begin{eqn}{8.1.15.1}\setnumber{8.1.15.1}\label{hodge-theory-iii-8.1.15.1}\skipsinglepar\[             {}_{\delta (W,L)}E_1^{ab}(\mathrm {s} K\otimes \mathbb {Q})             = \sum _{\substack {\alpha =\gamma -\beta \\b=\alpha +\beta }} \operatorname {H}^\alpha (\operatorname {Gr}_\beta ^W K^{\bullet \gamma })           \tag{8.1.15.1}           \]\end{eqn}

        where the complex \(({}_{\delta (W,L)}E_1^{\bullet  b},d_1)\) is the simple complex associated to the double complex of Hodge \(A\otimes \mathbb {Q}\)-structures of weight \(b\) given by

        \begin{eqn}{8.1.15.2}\setnumber{8.1.15.2}\label{hodge-theory-iii-8.1.15.2}\skipsinglepar\[             \small              \begin {CD}               \operatorname {H}^{b-(\beta +1)}(\operatorname {Gr}_{\beta +1}^W K^{\bullet ,\gamma +1}) @>{\partial }>> \operatorname {H}^{b-\beta }(\operatorname {Gr}_\beta ^W K^{\bullet ,\gamma +1}) @>{\partial }>> \operatorname {H}^{b-(\beta -1)}(\operatorname {Gr}_{\beta -1}^W K^{\bullet ,\gamma +1})             \\@. @A{d''}AA @AA{d''}A             \\@. \operatorname {H}^{b-\beta }(\operatorname {Gr}_\beta ^W K^{\bullet ,\gamma }) @>>{\partial }> \operatorname {H}^{b-(\beta -1)}(\operatorname {Gr}_{\beta -1}^W K^{\bullet ,\gamma })             \end {CD}           \tag{8.1.15.2}           \]\end{eqn}
      


      
 \item[(iii)]
        
        On the \({}_LE_r\) pages of the spectral sequence defined by \((\mathrm {s} K_{A\otimes \mathbb {Q}},L)\), the recurrent filtration and the two direct filtrations defined by \(\operatorname {Dec}_1(W)\) agree;
        similarly, on the \({}_LE_r(\mathrm {s} K_\mathbb {C})\), the recurrent filtration and the two direct filtrations induced by \(F\) agree.
        We denoted these filtrations again by \(W\) and \(F\) (respectively).
      


      
 \item[(iv)]
        For \(r\geqslant 1\), \(({}_LE_r,W,F)\) is a mixed Hodge \(A\)-structure (\cref{hodge-theory-iii-0.4}).
        The differentials \(d_r\) are morphisms of mixed Hodge \(A\)-structures.
      


      
 \item[(v)]
        The filtration \(L\) on \(\operatorname {H}^\bullet (\mathrm {s} K)\) is a filtration in the category of mixed Hodge structures, and

        \begin{eqn}{8.1.15.3}\setnumber{8.1.15.3}\label{hodge-theory-iii-8.1.15.3}\skipsinglepar\[             \operatorname {Gr}_L^p\big ( \operatorname {H}^{p+q}(\mathrm {s} K), \delta (W,L)[p+q], F \big )             = ({}_LE_\infty ^{pq},W,F).           \tag{8.1.15.3}           \]\end{eqn}
      

    \end{enumerate}\end{theorem}\begin{remark}{8.1.16}\setnumber{8.1.16}\label{hodge-theory-iii-8.1.16}\skipsinglepar\par{}In \hyperref[hodge-theory-iii-8.1.15]{Theorem~\ref*{hodge-theory-iii-8.1.15}}, the filtration \(\operatorname {Dec}_1(W)\) (\cref{hodge-theory-iii-8.1.10}) plays a more important role than \(W\).
    We have

    \begin{eqn}{8.1.16.1}\setnumber{8.1.16.1}\label{hodge-theory-iii-8.1.16.1}\skipsinglepar\[         \operatorname {Dec}(\delta (W,L))         = \operatorname {Dec}_1(W)         \qquad \text {on }\mathrm {s} K       \tag{8.1.16.1}       \]\end{eqn}

    and \hyperref[hodge-theory-iii-8.1.15.3]{Equation~\ref*{hodge-theory-iii-8.1.15.3}} can be rewritten as

    \begin{eqn}{8.1.16.2}\setnumber{8.1.16.2}\label{hodge-theory-iii-8.1.16.2}\skipsinglepar\[         \operatorname {Gr}_L(\operatorname {H}(K),\operatorname {Dec}_1(W,F))         = ({}_LE_\infty ,W,F).       \tag{8.1.16.2}       \]\end{eqn}

    In the same spirit, we note that, with the notation of \hyperref[hodge-theory-iii-8.1.8]{Scholium~\ref*{hodge-theory-iii-8.1.8}},
    \[       \mathrm {R}\Gamma \big ( \mathrm {R} j_*\mathbb {Z}, (\mathrm {R} j_*\mathbb {Q},\operatorname {Dec} W), (\Omega _X^\bullet (\log  D),\operatorname {Dec} W,F) \big )     \]
    depends only on \(U\), and not on the choice of compactification \(X\).\end{remark}\label{hodge-theory-iii-8.1.17}\par{}We now prove \hyperref[hodge-theory-iii-8.1.15]{Theorem~\ref*{hodge-theory-iii-8.1.15}}.
    Claims (i) and (ii) follow from the fact that
    \[       \begin {aligned}         \operatorname {Gr}_n^{\delta (W,L)}(\mathrm {s} K\otimes \mathbb {Q})         &= \sum _{p-q=n} \operatorname {Gr}_p^W\operatorname {Gr}_L^q(\mathrm {s} K\otimes \mathbb {Q})       \\&= \sum _{p-q=n} (\operatorname {Gr}_p^W K^{\bullet  q})[-q]       \end {aligned}     \]
    (this isomorphism is compatible with \(F\), and we apply \hyperref[hodge-theory-iii-8.1.4]{Remark~\ref*{hodge-theory-iii-8.1.4}}).
    The proof of \hyperref[hodge-theory-iii-8.1.15.2]{Equation~\ref*{hodge-theory-iii-8.1.15.2}} is left to the reader.\par{}Consider the claims:

    
 \begin{enumerate}
      
 \item[(aᵣ)]
        The morphisms \(d_r\) of \({}_LE_r(\mathrm {s} K_{A\otimes \mathbb {Q}})\) are strictly compatible with the recurrent filtration defined by \(\operatorname {Dec}_1(W)\).
      


      
 \item[(bᵣ)]
        The morphisms \(d_r\) of \({}_LE_r(\mathrm {s} K_\mathbb {C})\) are strictly compatible with the recurrent filtration defined by \(F\).
      


      
 \item[(cᵣ)]
        \({}_LE_r\), endowed with the filtrations considered in (a) and (b) is a mixed Hodge \(A\)-structure.
      


      
 \item[(dᵣ)]
        The \(d_r\) are morphisms of mixed Hodge \(A\)-structures.
      

    \end{enumerate}\par{}We will prove (a\({}_r\)) and (b\({}_r\)) for \(r\geqslant 0\), and (c\({}_r\)) and (d\({}_r\)) for \(r\geqslant 1\), by simultaneous induction.

    
 \begin{enumerate}
      
 \item[(A)]
        \emph{Proof of (a\({}_0\)).}
        We have
        \[           \big ( {}_LE_0^{p\bullet }(\mathrm {s} K_{A\otimes \mathbb {Q}}), \operatorname {Dec}_1 W \big )           = (K_{A\otimes \mathbb {Q}}^{\bullet  p}, \operatorname {Dec} W).         \]
        By (iv) of \hyperref[hodge-theory-iii-8.1.9]{Scholium~\ref*{hodge-theory-iii-8.1.9}}, the spectral sequence of \((K_{A\otimes \mathbb {Q}}^{\bullet  p},W)\) degenerates at \(E_2\).
        By \hyperref[hodge-theory-ii-1.3.4]{Proposition~\ref*{hodge-theory-ii-1.3.4}}, the spectral sequence of \((K_{A\otimes \mathbb {Q}}^{\bullet  p},\operatorname {Dec} W)\) degenerates at \(E_1\), and this proves (a\({}_0\)) (\hyperref[hodge-theory-ii-1.3.2]{Proposition~\ref*{hodge-theory-ii-1.3.2}}).
      


      
 \item[(B)]
        
        \emph{Proof of (b\({}_0\)).}
        We have
        \[           \big ( {}_LE_0^{p\bullet }(\mathrm {s} K_\mathbb {C}),F \big )           = (K_\mathbb {C}^{\bullet  p}, F).         \]
        We apply (v) of \hyperref[hodge-theory-iii-8.1.9]{Scholium~\ref*{hodge-theory-iii-8.1.9}} and \hyperref[hodge-theory-ii-1.3.2]{Proposition~\ref*{hodge-theory-ii-1.3.2}}.
      


      
 \item[(C)]
        \emph{Proof of (c\({}_1\)).}
        We have to show that \(\operatorname {H}^n(K_A^{\bullet  p})\), endowed with the filtrations induced by \(\operatorname {Dec} W\) on \(\operatorname {H}^n(K_{A\otimes \mathbb {Q}}^{\bullet  p})=\operatorname {H}^n(K_A^{\bullet  p})\otimes \mathbb {Q}\) and by \(F\) on \(\operatorname {H}^n(K_\mathbb {C}^{\bullet  p})=\operatorname {H}^n(K_A^{\bullet  p})\otimes _A\mathbb {C}\), is a mixed Hodge \(A\)-structure.
        But \(\operatorname {Dec} W\) induces on \(\operatorname {H}^n(K_{A\otimes \mathbb {Q}}^{\bullet  p})\) the same filtration as \(W[n]\), and so we conclude by (ii) of \hyperref[hodge-theory-iii-8.1.9]{Scholium~\ref*{hodge-theory-iii-8.1.9}}.
      


      
 \item[(D)]
        \emph{Proof of (c\({}_r\))+(d\({}_r\))\(\implies \)(a\({}_r\))+(b\({}_r\))+(c\({}_{r+1}\)).}
        This is the key point of the proof, and is a simple application of \hyperref[hodge-theory-ii-2.3.5]{Theorem~\ref*{hodge-theory-ii-2.3.5}} (along with \cref{hodge-theory-iii-0.4}).
      


      
 \item[(E)]
        \emph{Proof of (a\({}_r\))\({}_{r\leqslant  r_0-2}\)+(b\({}_r\))\({}_{r\leqslant  r_0-2}\)+(c\({}_{r_0}\))\(\implies \)(d\({}_{r_0}\)).}
        By \hyperref[hodge-theory-ii-1.3.16]{Theorem~\ref*{hodge-theory-ii-1.3.16}}, the recurrent filtration and the direct filtrations induced by \(\operatorname {Dec}_1W\) (resp. \(F\)) on \({}_LE_r\) agree for \(r\leqslant  r_0\), and we apply (i) of \hyperref[hodge-theory-ii-1.3.13]{Proposition~\ref*{hodge-theory-ii-1.3.13}}.
      

    \end{enumerate}\par{}This finishes the proof by induction of (a), (b), (c), and (d).\par{}Claim (iii) was proven in part (E) of the proof by induction, and the claims of (iv) are (c\({}_r\)) and (d\({}_r\)).
    We will prove \hyperref[hodge-theory-iii-8.1.15.3]{Equation~\ref*{hodge-theory-iii-8.1.15.3}} in the form of \hyperref[hodge-theory-iii-8.1.16.2]{Equation~\ref*{hodge-theory-iii-8.1.16.2}}.\par{}The fact that
    \[       \operatorname {Gr}_L(\operatorname {H}(\mathrm {s} K_{A\otimes \mathbb {Q}}),\operatorname {Dec}_1W)       = ({}_LE_\infty (\mathrm {s} K_{A\otimes \mathbb {Q}}),W)     \]
    follows from (a\({}_r\)) and \hyperref[hodge-theory-ii-1.3.17]{Corollary~\ref*{hodge-theory-ii-1.3.17}}.
    The fact that
    \[       \operatorname {Gr}_L(\operatorname {H}(\mathrm {s} K_\mathbb {C}),F)       = ({}_LE_\infty (\mathrm {s} K_\mathbb {C}),F)     \]
    follows from (b\({}_r\)) and \hyperref[hodge-theory-ii-1.3.17]{Corollary~\ref*{hodge-theory-ii-1.3.17}}.\par{}We conclude by the following lemma, whose proof is left to the reader:\begin{lemma}{8.1.18}\setnumber{8.1.18}\label{hodge-theory-iii-8.1.18}\skipsinglepar\par{}Let \(H=(H_A,W,F)\) be a mixed Hodge \(A\)-structure.
    For a filtration \(L\) of \(H_A\) to come from a filtration of \(H\), in the abelian category of mixed Hodge \(A\)-structures, it is necessary and sufficient that, for all \(n\), \((\operatorname {Gr}_L^n(H_A),\operatorname {Gr}_L^n(W),\operatorname {Gr}_L^n(F))\) be a mixed Hodge \(A\)-structure.\end{lemma}\label{hodge-theory-iii-8.1.19}\par{}We now apply \cref{hodge-theory-iii-8.1.13} and \hyperref[hodge-theory-iii-8.1.15]{Theorem~\ref*{hodge-theory-iii-8.1.15}} to \hyperref[hodge-theory-iii-8.1.12]{Example~\ref*{hodge-theory-iii-8.1.12}}.\par{}We obtain a mixed Hodge complex consisting of (a), (b), and (c) below.\begin{enumerate}
    
 \item[(a)]
      The complex \(\mathrm {R}\Gamma \mathrm {R} j_*\mathbb {Z}\simeq \mathrm {R}\Gamma (U_\bullet ,\mathbb {Z})\), whose cohomology groups are the \(\operatorname {H}^n(U_\bullet ,\mathbb {Z})\).
    


    
 \item[(b)]
      The filtration \(\delta (W,L)\) of \(\mathrm {R}\Gamma \mathrm {R} j_*\mathbb {Q}\simeq \mathrm {R}\Gamma (U_\bullet ,\mathbb {Q})\).
      With the notation of \cref{hodge-theory-ii-3.1.4}, we have
      \[         \begin {aligned}           \operatorname {Gr}_n^{\delta (W,L)}\mathrm {R}\Gamma (U_\bullet ,\mathbb {Q})           &\underset {(7.1.6.5)}{\simeq } \bigoplus _m \operatorname {Gr}_{n+m}^W \mathrm {R}\Gamma (U_m,\mathbb {Q})[-m]         \\&\simeq  \bigoplus _m \mathrm {R}\Gamma (\widetilde {Y}_m^{n+m},\varepsilon _\mathbb {Q}^{n+m}[-n-m])[-m]         \\&= \bigoplus _m \mathrm {R}\Gamma (\widetilde {Y}_m^{n+m},\varepsilon _\mathbb {Q}^{n+m})[-n-2m].         \end {aligned}       \]
      
      The \(E_1\) page of the corresponding spectral sequence is the sum of the groups of form \(\operatorname {H}^p(\widetilde {Y}_q^r,\varepsilon _\mathbb {Q}^r)\);
      this contributes to the \(\operatorname {H}^{p+q+r}(U_\bullet )\);
      it corresponds to \(\operatorname {Gr}_a^{\delta (W,F)}\) for \(p+2r=(p+q+r)+a\):

      \begin{eqn}{8.1.19.1}\setnumber{8.1.19.1}\label{hodge-theory-iii-8.1.19.1}\skipsinglepar\[           {}_{\delta (W,L)}E_1^{-a,b}           = \bigoplus _{\substack {p+2r=b\\q-r=-a}} \operatorname {H}^p(\widetilde {Y}_q^r,\varepsilon _\mathbb {Q}^r)           \Rightarrow  \operatorname {H}^{-a+b}(U_\bullet ,\mathbb {Q}).         \tag{8.1.19.1}         \]\end{eqn}
    


    
 \item[(c)]
      The bifiltration \((\delta (W,L),F)\) of \(\mathrm {R}\Gamma \mathrm {R} j_*\mathbb {C}\simeq \mathrm {s}\mathrm {R}\Gamma ^\bullet \Omega _X^\bullet (\log  D)\).
    

  \end{enumerate}\par{}The filtration \(F\) endows \({}_{\delta (W,L)}E_1^{-a,b}\) with a Hodge structure of weight \(b\), and the differentials \(d_1\) of \({}_{\delta (W,L)}E\) are the morphisms of the Hodge \(\mathbb {Q}\)-structures.
    The complex \({}_{\delta (W,L)}E_q^{\bullet  b}\) is the simple complex associated the double complex of Hodge structures of weight \(b\) given by:
    \[       \small        \begin {CD}         \operatorname {H}^{b-2(r+2)}(\widetilde {Y}_{q+1}^{r+1},\varepsilon _\mathbb {Q}^{r+1})[-r-1] @>{\mathrm {Gysin}}>> \operatorname {H}^{b-2r}(\widetilde {Y}_{q+1}^{r},\varepsilon _\mathbb {Q}^{r})[-r] @>{\mathrm {Gysin}}>> \operatorname {H}^{b-2(r-2)}(\widetilde {Y}_{q+1}^{r-1},\varepsilon _\mathbb {Q}^{r-1})[-r+1]       \\ @. @A{\sum _i(-1)^i\delta _i}AA @AA{\sum _i(-1)^i\delta _i}A       \\ @. \operatorname {H}^{b-2r}(\widetilde {Y}_q^r,\varepsilon _\mathbb {Q}^r)[-r] @>>{\mathrm {Gysin}}> \operatorname {H}^{b-2(r-2)}(\widetilde {Y}_q^{r-1},\varepsilon _\mathbb {Q}^{r-1})[-r+1]       \end {CD}     \]
    The indices in the terms in \(E_1^{ab}\) satisfy \(-a=q-r\).\par{}In particular, the \(\operatorname {H}^n(U_\bullet ,\mathbb {Z})\) are endowed with mixed Hodge structures (\hyperref[hodge-theory-iii-8.1.9]{Scholium~\ref*{hodge-theory-iii-8.1.9}}).\begin{proposition}{8.1.20}\setnumber{8.1.20}\label{hodge-theory-iii-8.1.20}\skipsinglepar\par{}Under the above hypotheses:

    
 \begin{enumerate}
      
 \item[(i)]
        The mixed Hodge structure of \(\operatorname {H}^n(U_\bullet ,\mathbb {Z})\) is functorial in the pair \((U_\bullet ,X_\bullet )\).
      


      
 \item[(ii)]
        In rational cohomology, the spectral sequence \hyperref[hodge-theory-iii-8.1.19.1]{Equation~\ref*{hodge-theory-iii-8.1.19.1}} degenerates at \(E_2\) (i.e. \(E_2=E_\infty \)) and abuts to the weight filtration of \(\operatorname {H}^n(U_\bullet ,\mathbb {Q})\).
        The Hodge structure of \(E_2\) induced by that of \(E_1\) is also that of \(\operatorname {Gr}_W\operatorname {H}^n(U_\bullet ,\mathbb {Q})\).
      


      
 \item[(iii)]
        The Hodge numbers of \(\operatorname {H}^n(U_\bullet ,\mathbb {Z})\) satisfy
        \[           h^{pq}\neq 0           \implies            0\leqslant  p,q\leqslant  n.         \]
      


      
 \item[(iv)]
        For \(Y_\bullet =\varnothing \), the Hodge numbers of \(\operatorname {H}^n(X_\bullet ,\mathbb {Z})\) satisfy
        \[           h^{pq}\neq 0           \implies            [0\leqslant  p,q\leqslant  n           \quad \text {and}\quad  p+q\leqslant  n].         \]
      

    \end{enumerate}


    
  
    
    \begin{proof}\par{}Both (i) and (ii) follow from the general theory (\hyperref[hodge-theory-iii-8.1.9]{Scholium~\ref*{hodge-theory-iii-8.1.9}}).
        To prove (iii), we contemplate \cref{hodge-theory-iii-8.1.19} and note that the Hodge structures \(\operatorname {H}^p(\widetilde {Y}_q^r,\varepsilon _\mathbb {Q}^r)[-r]\) (for \(p+q+r=n\)) satisfy (iii).
        To prove (iv), we note that the Hodge structures \(\operatorname {H}^p(X_q)\) (for \(p+q=n\)) satisfy (iv).\end{proof}\end{proposition}\label{hodge-theory-iii-8.1.21}\par{}We define a \emph{cohomological mixed Hodge \(A\)-complex on an \(r\)-simplicial topological space} by transposing \cref{hodge-theory-iii-8.1.11}.
    If \(K\) is such a complex, then we see, as in \cref{hodge-theory-iii-8.1.13}, that \(\mathrm {R}\Gamma ^\bullet  K\) is an \(r\)-cosimplicial mixed Hodge \(A\)-complex (\hyperref[hodge-theory-iii-8.1.10.1]{Variant~\ref*{hodge-theory-iii-8.1.10.1}}).\par{}Let \(X_\bullet \) be a smooth and proper \(r\)-simplicial scheme.
    A normal crossing divisor \(D_\bullet \) on \(X_\bullet \) is a system of normal crossing divisors on the \(X_n\) such that \(X_\bullet \setminus  D_\bullet =U_\bullet \) is an \(r\)-simplicial subscheme of \(X_\bullet \).
    Let \(j\colon  U_\bullet \hookrightarrow  X_\bullet \).
    We define, as in \hyperref[hodge-theory-iii-8.1.12]{Example~\ref*{hodge-theory-iii-8.1.12}}, a cohomological mixed Hodge complex \(\mathrm {R} j_*\mathbb {Z}\) on \(X_\bullet \).\par{}We can, as in \hyperref[hodge-theory-iii-8.1.20]{Proposition~\ref*{hodge-theory-iii-8.1.20}}, thus deduce a mixed Hodge structure on \(\operatorname {H}^\bullet (U_\bullet ,\mathbb {Z})\), but, since
    \[       \operatorname {H}^\bullet (U_\bullet ,\mathbb {Z})       = \operatorname {H}^\bullet (\delta  U_\bullet ,\mathbb {Z}),     \]
    this does not teach us anything new (the above equality would be an identity between mixed Hodge structures).\label{hodge-theory-iii-8.1.22}\par{}Set \(r=2\).
    For every bounded-below bicosimplicial complex \(K=(K^{pn_1n_2})\) (where \(p\) is the differential degree, and \(n_i\) is the \(i\)-th simplicial degree), we denote by \(\mathrm {s}_1K\) the cosimplicial complex of simplicial degree \(n_2\) given by
    \[       (\mathrm {s}_1K)^{\bullet  n_2}       = \mathrm {s} K^{\bullet \bullet  n_2}.     \]
    We denote by \(L_1\) the filtration on \(\mathrm {s}_1K\) that, on each \(\mathrm {s} K^{\bullet \bullet  n_2}\), induces the filtration by the simplicial degree \(n_1\).\par{}Let \(K\) be a bicosimplicial mixed Hodge \(A\)-complex.
    It follows from \hyperref[hodge-theory-iii-8.1.15]{Theorem~\ref*{hodge-theory-iii-8.1.15}} that \((\mathrm {s}_1K,\delta (W,L_1),F)\) is a cosimplicial mixed Hodge complex.
    Let \(L_2\) be the filtration of \(\mathrm {s} K=\mathrm {s}\mathrm {s}_1K\) by the simplicial degree \(n_2\).
    By \cref{hodge-theory-iii-8.1.14} and \hyperref[hodge-theory-iii-8.1.15]{Theorem~\ref*{hodge-theory-iii-8.1.15}}, the spectral sequence defined by \(L_2\) comes from a spectral sequence in the abelian category of mixed Hodge \(A\)-structures

    \begin{eqn}{8.1.22.1}\setnumber{8.1.22.1}\label{hodge-theory-iii-8.1.22.1}\skipsinglepar\[         {}_{L_2}E_1^{pq}         = \operatorname {H}^q(\mathrm {s} K^{\bullet \bullet  p})         \Rightarrow  \operatorname {H}^{p+q}(\mathrm {s} K).       \tag{8.1.22.1}       \]\end{eqn}\par{}In this spectral sequence, the mixed Hodge \(A\)-structure on \(E_1^{pq}\) is the mixed Hodge \(A\)-structure ((ii) of \hyperref[hodge-theory-iii-8.1.9]{Scholium~\ref*{hodge-theory-iii-8.1.9}}) of the cohomology of the mixed Hodge \(A\)-complex \((\mathrm {s} K^{\bullet \bullet  p},\delta (W,L_1),F)\).
    The mixed Hodge \(A\)-structure of the abutment is the mixed Hodge \(A\)-structure of the cohomology of the mixed Hodge \(A\)-complex \((\mathrm {s} K,\delta (\delta (W,L_1),L_2),F)\).
    The filtration \(\delta (\delta (W,L_1),L_2)\eqqcolon \delta (W,L_1,L_2)\) is given by
    \[       \delta (W,L_1,L_2)_k(K^{pn_1n_2})       = W_{k+n_1+n_2}K^{pn_1n_2}.     \]\label{hodge-theory-iii-8.1.23}\par{}In particular, a cohomological mixed Hodge \(A\)-complex on a bisimplicial topological space \(X_{\bullet \bullet }\) defines a spectral sequence of mixed Hodge \(A\)-structures

    \begin{eqn}{8.1.23.1}\setnumber{8.1.23.1}\label{hodge-theory-iii-8.1.23.1}\skipsinglepar\[         E_1^{pq}         = \operatorname {H}^q(X_{\bullet  p},K)         \Rightarrow  \operatorname {H}^{p+q}(X_{\bullet \bullet },K).       \tag{8.1.23.1}       \]\end{eqn}\par{}In the particular case considered in \cref{hodge-theory-iii-8.1.21}, for \(r=2\), we have \(A=\mathbb {Z}\), and \hyperref[hodge-theory-iii-8.1.23.1]{Equation~\ref*{hodge-theory-iii-8.1.23.1}} can be written as

    \begin{eqn}{8.1.23.2}\setnumber{8.1.23.2}\label{hodge-theory-iii-8.1.23.2}\skipsinglepar\[         E_1^{pq}         = \operatorname {H}^q(U_{\bullet  p},\mathbb {Z})         \Rightarrow  \operatorname {H}^{p+q}(U_{\bullet \bullet },\mathbb {Z}).       \tag{8.1.23.2}       \]\end{eqn}\label{hodge-theory-iii-8.1.24}\par{}If \(K'\) and \(K''\) are mixed Hodge \(A\)-complexes, then then their product \(K'\otimes  K''\), defined by
    \[       \begin {aligned}         (K'\otimes  K'')_A         &= K'_A\overset {\mathrm {L}}{\otimes } K''_A       \\((K'\otimes  K'')_{A\otimes \mathbb {Q}},W)         &= (K'_{A\otimes \mathbb {Q}},W')\otimes (K''_{A\otimes \mathbb {Q}},W'')       \\((K'\otimes  K'')_\mathbb {C},W,F)         &= (K'_\mathbb {C},W',F')\otimes (K''_\mathbb {C},W'',F'')       \end {aligned}     \]
    is again a mixed Hodge \(A\)-complex.\par{}Let \(U'\) (resp. \(U''\)) be the complement of a normal crossing divisor \(Y'\) (resp. \(Y''\)) in a smooth and proper scheme \(X'\) (resp. \(X''\)).
    Then \(U=U'\times  U''\) is the complement of a normal crossing divisor \(Y\) in \(X=X'\times  X''\).
    Let \(j\) (resp. \(j'\), \(j''\)) be the inclusion of \(U\) (resp. \(U'\), \(U''\)) into \(X\) (resp. \(X'\), \(X''\)).\par{}We have a quasi-isomorphism
    \[       \mathrm {R} j'_*\mathbb {Z}\overset {\mathrm {L}}{\boxtimes } \mathrm {R} j''_*\mathbb {Z}       \coloneqq  \operatorname {pr}_1^*\mathrm {R} j'_*\mathbb {Z}\overset {\mathrm {L}}{\otimes }\operatorname {pr}_2^*\mathrm {R} j''_*\mathbb {Z}       \xrightarrow {\sim } \mathrm {R} j_*\mathbb {Z},     \]
    a filtered quasi-isomorphism
    \[       (\mathrm {R} j'_*\mathbb {Q},\tau _\leqslant )\overset {\mathrm {L}}{\boxtimes }(\mathrm {R} j''_*\mathbb {Q},\tau _\leqslant )       \xrightarrow {\sim } (\mathrm {R} j_*\mathbb {Q},\tau _\leqslant ),     \]
    and a bifiltered morphism
    \[       (\Omega _{X'}^\bullet (\log  Y'),W,F)\boxtimes (\Omega _{X''}(\log  Y''),W,F)       \to  (\Omega _X^\bullet (\log  Y),W,F).     \]\par{}Applying the functor \(\mathrm {R}\Gamma \) gives us, by the Künneth formula, and by its analogue for cohomology of coherent analytic sheaves, an isomorphism of mixed Hodge complexes
    \[       \mathrm {R}\Gamma (U',\mathbb {Z})\overset {\mathrm {L}}{\otimes }\mathrm {R}\Gamma (U'',\mathbb {Z})       \xrightarrow {\sim } \mathrm {R}\Gamma (U,\mathbb {Z})     \]
    whence an isomorphism of mixed Hodge \(\mathbb {Q}\)-structures
    \[       \operatorname {H}^\bullet (U',\mathbb {Q})\otimes \operatorname {H}^\bullet (U'',\mathbb {Q})       \xrightarrow {\sim } \operatorname {H}^\bullet (U,\mathbb {Q}).     \]\label{hodge-theory-iii-8.1.25}\par{}Standard arguments, based on the Cartier–Eilenberg–Zilber theorem, give us the following simplicial variant of \cref{hodge-theory-iii-8.1.24}.\par{}Let \(U'_\bullet \) (resp. \(U''_\bullet \)) be the complement of a normal crossing divisor in a smooth and proper simplicial scheme \(X'_\bullet \) (resp. \(X''_\bullet \)).
    Let \(U_\bullet =U'_\bullet \times  U''_\bullet \) and \(X_\bullet =X'_\bullet \times  X''_\bullet \).\par{}We have an isomorphism of mixed Hodge complexes
    \[       \mathrm {R}\Gamma (U'_\bullet ,\mathbb {Z})\overset {\mathrm {L}}{\otimes }\mathrm {R}\Gamma (U''_\bullet ,\mathbb {Z})       \xrightarrow {\sim } \mathrm {R}\Gamma (U_\bullet ,\mathbb {Z})     \]
    and, in particular, an isomorphism of mixed Hodge \(\mathbb {Q}\)-structures
    \[       \operatorname {H}^\bullet (U'_\bullet ,\mathbb {Q})\otimes \operatorname {H}^\bullet (U''_\bullet ,\mathbb {Q})       \xrightarrow {\sim } \operatorname {H}^\bullet (U_\bullet ,\mathbb {Q}).     \]\label{hodge-theory-iii-8.1.26}\par{}Let \(U'_{\bullet \bullet }\) (resp. \(U''_{\bullet \bullet }\)) be the complement of a normal crossing divisor in a smooth and proper bisimplicial scheme \(X'_{\bullet \bullet }\) (resp. \(X''_{\bullet \bullet }\)).
    Let \(U_{\bullet \bullet }=U'_{\bullet \bullet }\times  U''_{\bullet \bullet }\) and \(X_{\bullet \bullet }=X'_{\bullet \bullet }\times  X''_{\bullet \bullet }\).
    
    Standard arguments show that, in rational cohomology, the spectral sequence \hyperref[hodge-theory-iii-8.1.23.2]{Equation~\ref*{hodge-theory-iii-8.1.23.2}} for \(U_{\bullet \bullet }\) is the tensor product of the spectral sequences \hyperref[hodge-theory-iii-8.1.23.2]{Equation~\ref*{hodge-theory-iii-8.1.23.2}} for \(U'_{\bullet \bullet }\) and \(U''_{\bullet \bullet }\).\section{Separated algebraic spaces}\label{hodge-theory-iii-8.2}\label{hodge-theory-iii-8.2.1}\par{}Let \(X\) be a scheme (or an algebraic space) of finite type over \(\mathbb {C}\), and assume it to be separated.
    Then there exists a diagram

    \begin{eqn}{8.2.1.1}\setnumber{8.2.1.1}\label{hodge-theory-iii-8.2.1.1}\skipsinglepar\[         X         \xleftarrow {\alpha } Y_\bullet          \hookrightarrow  \overline {Y}_\bullet        \tag{8.2.1.1}       \]\end{eqn}

    in which \(Y_\bullet \) is the complement of a normal crossing divisor in a smooth and proper simplicial scheme \(\overline {Y}_\bullet \), and is a proper hypercover of \(X\).
    We have ((II) of \cref{hodge-theory-iii-5.3.5})

    \begin{eqn}{8.2.1.2}\setnumber{8.2.1.2}\label{hodge-theory-iii-8.2.1.2}\skipsinglepar\[         \operatorname {H}^\bullet (X)         \xrightarrow {\sim } \operatorname {H}^\bullet (Y_\bullet ).       \tag{8.2.1.2}       \]\end{eqn}\par{}In \cref{hodge-theory-iii-8.1.19}, we endow \(\operatorname {H}^\bullet (Y_\bullet )\) with a mixed Hodge structure.
    Let \(\operatorname {\mathscr {H}}(\alpha )\) be the mixed Hodge structure on \(\operatorname {H}^\bullet (X)\) induced by \hyperref[hodge-theory-iii-8.2.1.2]{Equation~\ref*{hodge-theory-iii-8.2.1.2}}.\begin{proposition}{8.2.2}\setnumber{8.2.2}\label{hodge-theory-iii-8.2.2}\skipsinglepar\par{}With the above notation, the mixed Hodge structure \(\operatorname {\mathscr {H}}(\alpha )\) on \(\operatorname {H}^\bullet (X)\) is independent of the choice of \(\alpha \).
    We call it the \emph{mixed Hodge structure of \(\operatorname {H}^\bullet (X,\mathbb {Z})\)}.
    For every morphism \(f\colon  X_1\to  X_2\) of schemes, \(f^*\colon \operatorname {H}^\bullet (X_2)\to \operatorname {H}^\bullet (X_1)\) is a morphism of mixed Hodge structures.\begin{proof}\par{}The proof uses \cref{hodge-theory-iii-6.2.8}, and is parallel to that of (C) of \cref{hodge-theory-ii-3.2.11} (cf. also \cref{hodge-theory-iii-8.3.3}).\end{proof}\end{proposition}\label{hodge-theory-iii-8.2.3}\par{}For smooth \(X\), we can show, by taking \(\overline {Y}_\bullet \) to be a constant simplicial scheme, that the mixed Hodge structure in \hyperref[hodge-theory-iii-8.2.2]{Proposition~\ref*{hodge-theory-iii-8.2.2}} coincides with that defined in \hyperref[hodge-theory-ii-3.2.12]{Definition~\ref*{hodge-theory-ii-3.2.12}}.\begin{theorem}{8.2.4}\setnumber{8.2.4}\label{hodge-theory-iii-8.2.4}\skipsinglepar\par{}The pairs \((p,q)\) such that the Hodge number \(h^{pq}\) of \(\operatorname {H}^n(X)\) are non-zero satisfy the following conditions:

    
 \begin{enumerate}
      
 \item[(i)]
        \((p,q)\in [0,n]\times [0,n]\).
      


      
 \item[(ii)]
        If \(\dim  X=N\), and \(n\geqslant  N\), then
        \[           (p,q)\in [n-N,N]\times [n-N,N].         \]
      


      
 \item[(iii)]
        If \(X\) is proper, then \(p+q\leqslant  n\).
      


      
 \item[(iv)]
        If \(X\) is smooth, then \(p+q\geqslant  n\).
        The same conclusion holds if \(X\) is an equidimensional "rational homology manifold" of dimension \(N\), i.e. if, for all \(x\in  X\),
        \[           \operatorname {H}^i_{(x)}(X,\mathbb {Q})           =           \begin {cases}             0 &\text {if }i\neq  2N           \\\mathbb {Q} &\text {if }i=2N.           \end {cases}         \]
      

    \end{enumerate}\par{}\begin {tikzpicture}[point/.style={circle, fill=black, inner sep=0pt, minimum size=2.5pt}]
      \begin {scope}
        \draw [->] (0,0) to (3,0) node[label={below:$p$}]{};
        \draw [->] (0,0) to (0,3) node[label={left:$q$}]{};
        \node [below] at (2,0) {\scriptsize $N$};
        \node [left] at (0,2) {\scriptsize $N$};
        \node [below] at (1.3,0) {\scriptsize $n$};
        \node [left] at (0,1.3) {\scriptsize $n$};
        \draw [dashed] (2,0) to (2,2.7);
        \draw [dashed] (0,2) to (2.7,2);
        \draw  (0,1.3) to (1.3,1.3) to (1.3,0) to cycle;
        \node  at (0.5,0.2) {\scriptsize  proper};
        \node  at (0.8,1.1) {\scriptsize  smooth};
      \end {scope}
      \begin {scope}[shift={(4,0)}]
        \draw [->] (0,0) to (4,0) node[label={below:$p$}]{};
        \draw [->] (0,0) to (0,3) node[label={left:$q$}]{};
        \node [below] at (2.35,0) {\scriptsize $N\qquad  n$};
        \node [left] at (0,2) {\scriptsize $N$};
        \node [left] at (0,2.7) {\scriptsize $n$};
        \draw [dashed] (2,0) to (2,2.7);
        \draw [dashed] (0,2) to (2.7,2);
        \draw [dashed] (2.7,0) to (0,2.7);
        \draw  (0.7,2) to (2,2) to (2,0.7) to (0.7,0.7) to cycle;
        \draw  (0.7,2) to (2,0.7);
        \node  at (1.2,0.9) {\scriptsize  proper};
        \node  at (1.5,1.8) {\scriptsize  smooth};
      \end {scope}
    \end {tikzpicture}\begin{proof}\par{}Claim (ii) will be proven in \cref{hodge-theory-iii-8.3.10}.
      Claim (i) follows from (iii) of \hyperref[hodge-theory-iii-8.1.20]{Proposition~\ref*{hodge-theory-iii-8.1.20}}.
      For proper \(X\), we can find a diagram as in \hyperref[hodge-theory-iii-8.2.1.1]{Equation~\ref*{hodge-theory-iii-8.2.1.1}} in which \(Y_\bullet =\overline {Y}_\bullet \), and (iii) then follows from (iv) of \hyperref[hodge-theory-iii-8.1.20]{Proposition~\ref*{hodge-theory-iii-8.1.20}}.
      For smooth \(X\), (iv) follows from \cref{hodge-theory-iii-8.2.3} and \hyperref[hodge-theory-ii-3.2.15]{Corollary~\ref*{hodge-theory-ii-3.2.15}}.
      The general case then follows: if \(p\colon \widetilde {X}\to  X\) is a resolution of singularities of \(X\), then \(p^*\colon \operatorname {H}^\bullet (X,\mathbb {Q})\to \operatorname {H}^\bullet (\widetilde {X},\mathbb {Q})\) is injective, since the Poincaré dual transpose \(p_!\) of \(p^*\colon \operatorname {H}^\bullet _\mathrm {c}(X,\mathbb {Q})\to \operatorname {H}^\bullet _\mathrm {c}(\widetilde {X},\mathbb {Q})\) is a retraction for it.\end{proof}\end{theorem}\begin{proposition}{8.2.5}\setnumber{8.2.5}\label{hodge-theory-iii-8.2.5}\skipsinglepar\par{}Suppose that \(X\) is proper.
    If \(u\colon  Y\to  X\) is a surjective proper morphism, with \(Y\) smooth, then the weight-\(n\) quotient of \(\operatorname {H}^n(X,\mathbb {Q})\) is the image of \(\operatorname {H}^n(X,\mathbb {Q})\) in \(\operatorname {H}^n(Y,\mathbb {Q})\).
    The sequence

    \begin{eqn}{8.2.5.1}\setnumber{8.2.5.1}\label{hodge-theory-iii-8.2.5.1}\skipsinglepar\[         \operatorname {H}^n(X,\mathbb {Q})         \to  \operatorname {H}^n(Y,\mathbb {Q})         \xrightarrow {\operatorname {pr}_1^*-\operatorname {pr}_2^*} \operatorname {H}^n(Y\times _X Y,\mathbb {Q})       \tag{8.2.5.1}       \]\end{eqn}

    is exact.\begin{proof}\par{}Let \(Y_\bullet \) be as in \hyperref[hodge-theory-iii-8.2.1.1]{Equation~\ref*{hodge-theory-iii-8.2.1.1}}, with \(Y_0=Y\) and \(Y_\bullet =\overline {Y}_\bullet \).
      By (ii) of \hyperref[hodge-theory-iii-8.1.20]{Proposition~\ref*{hodge-theory-iii-8.1.20}}, we have an exact sequence

      \begin{eqn}{8.2.5.2}\setnumber{8.2.5.2}\label{hodge-theory-iii-8.2.5.2}\skipsinglepar\[           0           \to  \operatorname {H}^n(X,\mathbb {Q})/W^{n-1}\operatorname {H}^n(X,\mathbb {Q})           \to  \operatorname {H}^n(Y_0,\mathbb {Q})           \underset {\delta _1}{\overset {\delta _0}{\rightrightarrows }} \operatorname {H}^n(Y_1,\mathbb {Q})         \tag{8.2.5.2}         \]\end{eqn}

      which implies \hyperref[hodge-theory-iii-8.2.5]{Proposition~\ref*{hodge-theory-iii-8.2.5}} (conversely, this exact sequence can easily be deduced from \hyperref[hodge-theory-iii-8.2.5]{Proposition~\ref*{hodge-theory-iii-8.2.5}} and \hyperref[hodge-theory-iii-8.2.7]{Proposition~\ref*{hodge-theory-iii-8.2.7}} below).\end{proof}\end{proposition}\begin{proposition}{8.2.6}\setnumber{8.2.6}\label{hodge-theory-iii-8.2.6}\skipsinglepar\par{}Let
    \[       Y       \xrightarrow {f} X       \overset {j}{\hookrightarrow } \overline {X}     \]
    be morphisms of schemes.
    Suppose that \(Y\) is proper and \(X\) is smooth, and that \(\overline {X}\) is a smooth compactification of \(X\).
    Then \(\operatorname {H}^\bullet (\widetilde {X},\mathbb {Q})\) and \(\operatorname {H}^\bullet (X,\mathbb {Q})\) have the same image in \(\operatorname {H}^\bullet (Y,\mathbb {Q})\).\begin{proof}\par{}The proof is parallel to that of \hyperref[hodge-theory-ii-3.2.18]{Corollary~\ref*{hodge-theory-ii-3.2.18}}.
      It suffices to prove, for any \(i\) and any \(n\), that \(\operatorname {Gr}_i^W(f^*)\) and \(\operatorname {Gr}_i^W((jf)^*)\) have the same image in \(\operatorname {Gr}_i^W\operatorname {H}^n(Y,\mathbb {Q})\).

      
 \begin{enumerate}
        
 \item[(a)]
          For \(i>n\), these images are both zero since \(\operatorname {Gr}_i^W\operatorname {H}^n(Y,\mathbb {Q})=0\) ((iii) of \hyperref[hodge-theory-iii-8.2.4]{Theorem~\ref*{hodge-theory-iii-8.2.4}}).
        

        
 \item[(b)]
          For \(i<n\), these images are both zero since \(\operatorname {Gr}_i^W\operatorname {H}^n(X,\mathbb {Q})=0\) ((iv) of \hyperref[hodge-theory-iii-8.2.4]{Theorem~\ref*{hodge-theory-iii-8.2.4}}).
        

        
 \item[(c)]
          For \(i=n\), these images are equal since \(\operatorname {Gr}_n^W(j^*)\) is surjective (\hyperref[hodge-theory-ii-3.2.17]{Corollary~\ref*{hodge-theory-ii-3.2.17}}).
        

      \end{enumerate}\end{proof}\end{proposition}\begin{proposition}{8.2.7}\setnumber{8.2.7}\label{hodge-theory-iii-8.2.7}\skipsinglepar\par{}Let
    \[       \widetilde {X}       \xrightarrow {\pi } X       \xrightarrow {f} Y     \]
    be morphisms of schemes.
    Suppose that \(Y\) is smooth, \(X\) is proper, \(\widetilde {X}\) is smooth and proper, and that \(\pi \) is surjective.
    Then the kernels of \(f^*\) and \((f\pi )^*\) in \(\operatorname {H}^\bullet (Y,\mathbb {Q})\) are equal.\begin{proof}\par{}We will prove that \(\operatorname {Ker}(\operatorname {Gr}_i^W(f^*))=\operatorname {Ker}(\operatorname {Gr}_i^W((f\pi )^*))\) in \(\operatorname {Gr}_i^W\operatorname {H}^n(Y,\mathbb {Q})\).

      
 \begin{enumerate}
        
 \item[(a)]
          For \(i<n\), these kernels are both zero since \(\operatorname {Gr}_i^W\operatorname {H}^n(Y,\mathbb {Q})=0\) ((iv) of \hyperref[hodge-theory-iii-8.2.4]{Theorem~\ref*{hodge-theory-iii-8.2.4}}).
        

        
 \item[(b)]
          For \(i>n\), these kernels are both all of \(\operatorname {Gr}_i^W\operatorname {H}^n(Y,\mathbb {Q})\) since \(\operatorname {Gr}_i^W\operatorname {H}^n(X,\mathbb {Q})=0\) ((iii) of \hyperref[hodge-theory-iii-8.2.4]{Theorem~\ref*{hodge-theory-iii-8.2.4}}).
        

        
 \item[(a)]
          For \(i=n\), these kernels are equal since \(\operatorname {Gr}_n^W(\pi ^*)\) is injective (\hyperref[hodge-theory-iii-8.2.5.2]{Equation~\ref*{hodge-theory-iii-8.2.5.2}}).
        

      \end{enumerate}\end{proof}\end{proposition}\begin{corollary}{8.2.8}\setnumber{8.2.8}\label{hodge-theory-iii-8.2.8}\skipsinglepar\par{}Let
    \[       \widetilde {X}       \xrightarrow {\pi } X       \overset {f}{\hookrightarrow } Y     \]
    be morphisms of schemes.
    Suppose that \(Y\) is smooth and proper, that \(X\) is a closed subscheme of \(Y\), and that \(\widetilde {X}\) is a resolution of singularities of \(X\).
    Let \(q=f\pi \) and \(U=Y\setminus  X\).
    Then the sequence

    \begin{eqn}{8.2.8.1}\setnumber{8.2.8.1}\label{hodge-theory-iii-8.2.8.1}\skipsinglepar\[         \operatorname {H}^\bullet (\widetilde {X},\mathbb {Q})         \xrightarrow {q_!} \operatorname {H}^\bullet (Y,\mathbb {Q})         \to  \operatorname {H}^\bullet (U,\mathbb {Q})       \tag{8.2.8.1}       \]\end{eqn}

    is exact.\begin{proof}\par{}By \hyperref[hodge-theory-iii-8.2.7]{Proposition~\ref*{hodge-theory-iii-8.2.7}}, we have that
      \[         \operatorname {Ker}(f^*\colon \operatorname {H}^\bullet (Y,\mathbb {Q})\to \operatorname {H}^\bullet (X,\mathbb {Q}))         = \operatorname {Ker}(q^*\colon \operatorname {H}^\bullet (Y,\mathbb {Q})\to \operatorname {H}^\bullet (\widetilde {X},\mathbb {Q})).       \]
      The long exact sequence in cohomology of the pair \(Y,X\) thus gives an exact sequence

      \begin{eqn}{8.2.8.2}\setnumber{8.2.8.2}\label{hodge-theory-iii-8.2.8.2}\skipsinglepar\[           \operatorname {H}_c^\bullet (U,\mathbb {Q})           \to  \operatorname {H}^\bullet (Y,\mathbb {Q})           \to  \operatorname {H}^\bullet (\widetilde {X},\mathbb {Q})         \tag{8.2.8.2}         \]\end{eqn}

      and \hyperref[hodge-theory-iii-8.2.8.1]{Equation~\ref*{hodge-theory-iii-8.2.8.1}} is the transpose sequence of \hyperref[hodge-theory-iii-8.2.8.2]{Equation~\ref*{hodge-theory-iii-8.2.8.2}} under Poincaré duality.\end{proof}\end{corollary}\begin{remark}{8.2.9}\setnumber{8.2.9}\label{hodge-theory-iii-8.2.9}\skipsinglepar\par{}\hyperref[hodge-theory-iii-8.2.8]{Corollary~\ref*{hodge-theory-iii-8.2.8}} is used under more general hypotheses in [\hyperref[G1968]{G1968}, (9.7), p. 161].
    As Lieberman remarked, the justification given loc. cit. is insufficient.\end{remark}\par{}From \cref{hodge-theory-iii-8.1.25} we deduce the following:\begin{proposition}{8.2.10}\setnumber{8.2.10}\label{hodge-theory-iii-8.2.10}\skipsinglepar\par{}The Künneth isomorphisms
    \[       \operatorname {H}^\bullet (X\times  Y,\mathbb {Q})       \simeq  \operatorname {H}^\bullet (X,\mathbb {Q})\otimes \operatorname {H}^\bullet (Y,\mathbb {Q})     \]
    are isomorphisms of mixed Hodge structures.\end{proposition}\begin{corollary}{8.2.11}\setnumber{8.2.11}\label{hodge-theory-iii-8.2.11}\skipsinglepar\par{}The cup product
    \[       \operatorname {H}^\bullet (X)\otimes \operatorname {H}^\bullet (X)       \to  \operatorname {H}^\bullet (X)     \]
    is a morphism of mixed Hodge structures.\begin{proof}\par{}This follows from \hyperref[hodge-theory-iii-8.2.10]{Proposition~\ref*{hodge-theory-iii-8.2.10}} and from \hyperref[hodge-theory-iii-8.2.2]{Proposition~\ref*{hodge-theory-iii-8.2.2}} applied to the diagonal morphism \(\Delta \colon  X\to  X\times  X\).\end{proof}\end{corollary}\section{Hodge theory of simplicial schemes}\label{hodge-theory-iii-8.3}\label{hodge-theory-iii-8.3.1}\par{}For a separated simplicial scheme (or separated simplicial algebraic space) \(X^\bullet \), we will have to consider diagrams

    \begin{eqn}{8.3.1.1}\setnumber{8.3.1.1}\label{hodge-theory-iii-8.3.1.1}\skipsinglepar\end{eqn}

    of simplicial schemes satisfying:
    
 \begin{enumerate}
      
 \item[(a)]
        The morphism \(j\) is an open immersion with dense image; \(\overline {Z}_\bullet \) is proper.
      

      
 \item[(b)]
        The evident morphisms \(Z_n\to (\operatorname {cosq}_{n-1}^{X_\bullet }\operatorname {sq}_{n-1}Z_\bullet )_n\) are proper and surjective.
      

    \end{enumerate}\par{}For reduced \(\overline {Z}_\bullet \), we can consider \(j\) as a simplicial object in the category \(\mathscr {C}\) of pairs \((Y,\overline {Y})\) consisting of a proper reduced scheme \(\overline {Y}\) and a dense (Zariski) open \(Y\) of \(\overline {Y}\).
    The category \(\mathscr {C}\) satisfies \cref{hodge-theory-iii-6.2.4.1}, which allows us to apply the analogue of \hyperref[hodge-theory-iii-6.2.4]{Proposition~\ref*{hodge-theory-iii-6.2.4}}.\begin{proposition}{8.3.2}\setnumber{8.3.2}\label{hodge-theory-iii-8.3.2}\skipsinglepar\begin{enumerate}
    
 \item[(i)]
      For every separated simplicial scheme \(X_\bullet \), there exists a diagram of the form \hyperref[hodge-theory-iii-8.3.1.1]{Equation~\ref*{hodge-theory-iii-8.3.1.1}} that satisfies (a) and (b), and such that \(\overline {Z}_\bullet \) is smooth and proper, and \(Z_\bullet \) is the complement of a normal crossing divisor.
    


    
 \item[(ii)]
      Let \(F\colon  {X_\bullet }_1\to {X_\bullet }_2\) be a morphism of separated simplicial schemes.
      Let \(({X_\bullet }_i,{Z_\bullet }_i,{\overline {Z}_\bullet }_i)_{i=1,2}\) be diagrams of the form \hyperref[hodge-theory-iii-8.3.1.1]{Equation~\ref*{hodge-theory-iii-8.3.1.1}} that satisfy (a) and (b).
      Then there exists \(({X_\bullet }_1,Z_\bullet ,\overline {Z}_\bullet )\) as in (i), and a commutative diagram

      
  \begin {tikzcd}
        Z_\bullet 
          \ar [r,hook]
          \ar [d]
          \ar [drr]
        & \overline {Z}_\bullet 
          \ar [d]
          \ar [drr]
      \\{Z_\bullet }_1
          \ar [r,hook,swap,"j_1"]
          \ar [d,swap,"p_1"]
        & {\overline {Z}_\bullet }_1
        & {Z_\bullet }_2
          \ar [r,hook,swap,"j_2"]
          \ar [d,swap,"p_2"]
        & {\overline {Z}_\bullet }_2
      \\{X_\bullet }_1
          \ar [rr,"f"]
        && {X_\bullet }_2
      \end {tikzcd}

    

  \end{enumerate}\begin{proof}\par{}The proof, which is a more complicated analogue of \cref{hodge-theory-iii-6.2.5}, is left to the reader.\end{proof}\end{proposition}\label{hodge-theory-iii-8.3.3}\par{}Consider a diagram of the form \hyperref[hodge-theory-iii-8.3.1.1]{Equation~\ref*{hodge-theory-iii-8.3.1.1}} as in (ii) of \hyperref[hodge-theory-iii-8.3.2]{Proposition~\ref*{hodge-theory-iii-8.3.2}}.
    By (V) of \cref{hodge-theory-iii-5.3.5}, the map \(\operatorname {H}^\bullet {(X_\bullet ,\mathbb {Z})\to \operatorname {H}^\bullet (Z_\bullet ,\mathbb {Z})}\) is an isomorphism.
    In \cref{hodge-theory-iii-8.1.19} and \hyperref[hodge-theory-iii-8.1.20]{Proposition~\ref*{hodge-theory-iii-8.1.20}}, we endowed \(\operatorname {H}^\bullet (Z_\bullet ,\mathbb {Z})\) with a mixed Hodge structure, and obtained a mixed Hodge structure on \(\operatorname {H}^\bullet (X_\bullet ,\mathbb {Z})\).
    
    This is independent of the choice of diagram of the form \hyperref[hodge-theory-iii-8.3.1.1]{Equation~\ref*{hodge-theory-iii-8.3.1.1}}: by (ii) of \hyperref[hodge-theory-iii-8.3.2]{Proposition~\ref*{hodge-theory-iii-8.3.2}} for \(f=\mathrm {Id}\). any two systems \(({Z_\bullet }_i,{\overline {Z}_\bullet }_i,p_i)_{i=1,2}\) are always covered by a third \((Z_\bullet ,\overline {Z}_\bullet ,p)\), whence a commutative diagram of isomorphisms

    
  \begin {tikzcd}
      & \operatorname {H}^\bullet (X_\bullet ,\mathbb {Z})
        \ar [dl] \ar [d] \ar [dr]
    \\\operatorname {H}^\bullet ({Z_\bullet }_1,\mathbb {Z})
        \ar [r]
      & \operatorname {H}^\bullet (Z_\bullet ,\mathbb {Z})
      & \operatorname {H}^\bullet ({Z_\bullet }_2,\mathbb {Z})
        \ar [l]
    \end {tikzcd}\begin{definition}{8.3.4}\setnumber{8.3.4}\label{hodge-theory-iii-8.3.4}\skipsinglepar\par{}The \emph{mixed Hodge structure} of \(\operatorname {H}^\bullet (X_\bullet ,\mathbb {Z})\) is the one defined above.\end{definition}\par{}It follows from (ii) of \hyperref[hodge-theory-iii-8.3.2]{Proposition~\ref*{hodge-theory-iii-8.3.2}} that this mixed Hodge structure is functorial in \(X_\bullet \).\begin{proposition}{8.3.5}\setnumber{8.3.5}\label{hodge-theory-iii-8.3.5}\skipsinglepar\par{}Let \(X_\bullet \) be a separated simplicial scheme.
    Then the spectral sequence
    \[       E_1^{pq}       = \operatorname {H}^q(X_p,\mathbb {Z})       \Rightarrow  \operatorname {H}^{p+q}(X_\bullet ,\mathbb {Z})     \]
    is a spectral sequence of mixed Hodge structures.\end{proposition}\par{}For a diagram as in \hyperref[hodge-theory-iii-8.3.1.1]{Equation~\ref*{hodge-theory-iii-8.3.1.1}}, the maps \(p_n\colon  Z_n\to  X_n\) are proper and surjective.
  We can thus apply \hyperref[hodge-theory-iii-6.2.4]{Proposition~\ref*{hodge-theory-iii-6.2.4}} to the category of simplicial objects in the category \(\mathscr {C}\) (\cref{hodge-theory-iii-8.3.1}) to obtain, as in \cref{hodge-theory-iii-6.2.5}, the following lemma:\begin{lemma}{8.3.6}\setnumber{8.3.6}\label{hodge-theory-iii-8.3.6}\skipsinglepar\par{}Let \(X_\bullet \) be a separated simplicial scheme.
    Then there exists a bisimplicial scheme \(Z_{\bullet \bullet }\) augmented over \(X_\bullet \)
    \[       \varepsilon _{n,m}       \colon  Z_{n,m}       \to  X_n     \]
    and \(i\colon  Z_{\bullet \bullet }\hookrightarrow \overline {Z}_{\bullet \bullet }\) such that:

    
 \begin{enumerate}
      
 \item[(a)]
        \(\overline {Z}_{\bullet \bullet }\) is a smooth and proper bisimplicial scheme, and \(Z_{\bullet \bullet }\) is the complement of a normal crossing divisor in \(\overline {Z}_{\bullet \bullet }\);
      

      
 \item[(b)]
        for all \(n\), the augmented simplicial scheme \(Z_{n\bullet }\to  X_n\) is a proper hypercover of \(X_n\).
      

    \end{enumerate}\end{lemma}\par{}In \cref{hodge-theory-iii-8.1.21}, we endowed \(\operatorname {H}^\bullet (Z_{\bullet \bullet },\mathbb {Z})\) with a mixed Hodge structure.
  If \(\delta \mathbb {Z}_{\bullet \bullet }\) is the diagonal simplicial scheme of \(Z_{\bullet \bullet }\), then we have a commutative diagram of morphisms
  
  \begin {tikzcd}
    \operatorname {H}^\bullet (X_\bullet ,\mathbb {Z})
      \ar [r,"\gamma "]
      \ar [dr,swap,"\alpha "]
    & \operatorname {H}^\bullet (Z_{\bullet \bullet },\mathbb {Z})
      \ar [d,equal,"\beta "]
  \\& \operatorname {H}^\bullet (\delta  Z_{\bullet \bullet },\mathbb {Z})
  \end {tikzcd}

  
  in which \(\beta \) and \(\gamma \) are isomorphisms, and \(\alpha \) and \(\beta \) are morphisms of mixed Hodge structures;
  \(\gamma \) is thus an isomorphism of mixed Hodge structures.\par{}We have a morphism of spectral sequences
  
  \begin {tikzcd}
    \operatorname {H}^q(X_p,\mathbb {Z})
      \ar [r,Rightarrow]
      \ar [d]
    & \operatorname {H}^{p+q}(X_\bullet ,\mathbb {Z})
      \ar [d]
    && (5.2.3)
  \\\operatorname {H}^q(Z_{p\bullet },\mathbb {Z})
      \ar [r,Rightarrow]
    & \operatorname {H}^{p+q}(Z_{\bullet \bullet },\mathbb {Z})
    && (8.1.23.1)
  \end {tikzcd}

  and this morphism is an isomorphism that is compatible with the mixed Hodge structures, so that \hyperref[hodge-theory-iii-8.3.5]{Proposition~\ref*{hodge-theory-iii-8.3.5}} follows from \cref{hodge-theory-iii-8.1.23}.\begin{example}{8.3.7}\setnumber{8.3.7}\label{hodge-theory-iii-8.3.7}\skipsinglepar\par{}Let \(X_\bullet \) be a simplicial scheme, not necessarily separated.
    Then there exists a separated simplicial scheme \(Y_\bullet \) and \(u\colon  Y_\bullet \to  X_\bullet \) such that, for \(n\geq 0\), the map \(u'_n\colon  Y_n\to (\operatorname {cosq}_{n-1}^{X_\bullet }\operatorname {sq}_{n-1}Y_\bullet )_n\) is étale and surjective (for example, the sum of opens of an open cover).\par{}The map \(u^\bullet \colon \operatorname {H}^\bullet (X_\bullet ,\mathbb {Z})\to \operatorname {H}^\bullet (Y_\bullet ,\mathbb {Z})\) is an isomorphism ((V) and (III) of \cref{hodge-theory-iii-5.3.5}).
    By definition, the \emph{mixed Hodge structure} of \(\operatorname {H}^\bullet (X_\bullet ,\mathbb {Z})\) is induced from that of \(\operatorname {H}^\bullet (Y_\bullet ,\mathbb {Z})\) \emph{via \(u^\bullet \)}.
    It does not depend on the choice of \(Y_\bullet \).\par{}For \(X\) a non-separated scheme, the \emph{mixed Hodge structure} of \(\operatorname {H}^\bullet (X,\mathbb {Z})\) is that of \(\operatorname {H}^\bullet (X_\bullet ,\mathbb {Z})\), where \(X_\bullet \) is the constant simplicial scheme defined by \(X\).
    It is plausible, but not proven, that (iv) of \hyperref[hodge-theory-iii-8.2.4]{Theorem~\ref*{hodge-theory-iii-8.2.4}} remains true for \(X\) smooth and non-separated.\end{example}\begin{example}{8.3.8}\setnumber{8.3.8}\label{hodge-theory-iii-8.3.8}\skipsinglepar\par{}Let \(u\colon  X_\bullet \to  Y_\bullet \) be a morphism of simplicial schemes.
    The relative cohomology groups \(\operatorname {H}^n(Y_\bullet \operatorname {mod} X_\bullet ,\mathbb {Z})\) are, by definition (\cref{hodge-theory-iii-6.3.2}), the cohomology groups of the simplicial scheme \(C(u)\).
    They are thus endowed with a mixed Hodge structure.\end{example}\label{hodge-theory-iii-8.3.9}\par{}The long exact sequence of relative cohomology
    \[       \operatorname {H}^n(Y_\bullet ,\mathbb {Z})       \to  \operatorname {H}^n(X_\bullet ,\mathbb {Z})       \xrightarrow {j} \operatorname {H}^{n+1}(Y_\bullet \operatorname {mod} X_\bullet ,\mathbb {Z})       \to  \operatorname {H}^{n+1}(Y_\bullet ,\mathbb {Z})     \]
    is a long exact sequence of mixed Hodge structures.\begin{proof}\par{}We can reduce to the case where \(X_\bullet \) and \(Y_\bullet \) are smooth, and where, furthermore, there exists a commutative diagram
      
  \begin {tikzcd}
        X_\bullet  \ar [r,"u"] \ar [d,hook]
        & Y_\bullet  \ar [d,hook]
      \\\overline {X}_\bullet  \ar [r,swap,"\overline {u}"]
        & \overline {Y}_\bullet 
      \end {tikzcd}

      with \(\overline {X}_\bullet \) and \(\overline {Y}_\bullet \) both smooth and proper, and \(X_\bullet \) and \(Y_\bullet \) both complements of normal crossing divisors \(D_\bullet \) and \(E_\bullet \) in \(\overline {X}_\bullet \) and \(\overline {Y}_\bullet \) (respectively).\par{}There thus exist bifiltered resolutions (graded with acyclic components) \(a\colon \Omega _{\overline {X}_\bullet }^\bullet (\log  D_\bullet )\to  K\) and \(b\colon \Omega _{\overline {Y}_\bullet }^\bullet (\log  E_\bullet )\to  L\) that fit into a commutative diagram of bifiltered complexes
      
  \begin {tikzcd}
        \Omega _{\overline {X}_\bullet }^\bullet (\log  D_\bullet )
          \ar [d,swap,"a"]
        & \Omega _{\overline {Y}_\bullet }^\bullet (\log  E_\bullet )
          \ar [l,swap,"\overline {u}^*"]
          \ar [d,"b"]
      \\K
        & L
          \ar [l,"v"]
      \end {tikzcd}

      (where \(\overline {u}^*\) and \(v\) are \(\overline {u}\)-morphisms).
      The long exact sequence in \cref{hodge-theory-iii-8.3.9} is then identified with the exact sequence given by the distinguished triangle
      
  \begin {tikzcd}[column sep=tiny]
        & C(\Gamma  v)
          \ar [dl]
      \\\Gamma  L
          \ar [rr]
        && \Gamma  K
          \ar [ul]
      \end {tikzcd}

      The complex \(C(v)\) over \(C(\overline {u})\) is a bifiltered resolution of \(C(u^*\colon \Omega _{\overline {Y}_\bullet }^\bullet (\log  E_\bullet )\to \Omega _{\overline {X}_\bullet }^\bullet (\log  D_\bullet ))\).
      From this it results that the Hodge and weight filtrations of \(\operatorname {H}^\bullet (Y_\bullet \operatorname {mod} X_\bullet ,\mathbb {C})\) are induced by the filtrations of the cone \(C(\Gamma  v)\), and \cref{hodge-theory-iii-8.3.9} then follows formally.\end{proof}Proof of (ii) of (8.2.4)\label{hodge-theory-iii-8.3.10}\begin{proof}\par{}The proof proceeds by induction on \(N=\dim (X)\).
      We can assume \(X\) to be reduced.
      Let \(U\) be a smooth dense Zariski open of \(X\), and let \(Y=X\setminus  U\).
      By the resolution of singularities, there exists a proper birational morphism \(p\colon  X'\to  X\), with smooth source \(X'\), that induces an isomorphism from \(p^{-1}(U)\) to \(U\).
      Let \(Y'=p^{-1}(Y)\).
      Consider, for \(n\geqslant  N\), the diagram of mixed Hodge structures
      
  \begin {tikzcd}
        & \operatorname {H}^n(X\operatorname {mod} Y,\mathbb {Z})
          \ar [d,swap,"p^*"]
          \ar [r]
        & \operatorname {H}^n(X,\mathbb {Z})
          \ar [d]
          \ar [r]
        & \operatorname {H}^n(Y,\mathbb {Z})
          \ar [d]
      \\\operatorname {H}^{n-1}(Y',\mathbb {Z})
          \ar [r]
        & \operatorname {H}^n(X'\operatorname {mod} Y',\mathbb {Z})
          \ar [r]
        & \operatorname {H}^n(X',\mathbb {Z})
          \ar [r]
        & \operatorname {H}^n(Y',\mathbb {Z})
      \end {tikzcd}

      Let \(\Phi \) be the family of supports of \(U\) consisting of those \(A\subset  U\) such that, in \(X\) (or, equivalently, in \(X'\)), \(\overline {A}\subset  U\).
      We have that
      \[         \operatorname {H}^n(X\operatorname {mod} Y,\mathbb {Z})         = \operatorname {H}_\Phi ^n(U,\mathbb {Z})         = \operatorname {H}^n(X'\operatorname {mod} Y',\mathbb {Z})       \]
      and so the arrow denoted \(p^*\) is an isomorphism.\par{}For \(h^{pq}\) a non-zero Hodge number of one of the mixed Hodge structures in play, we have:
      
      
 \begin{enumerate}
        
 \item[(a)]
          for \(\operatorname {H}^n(X',\mathbb {Z})\), \((p,q)\in [n-N,N]\times [n-N,N]\)
        


        
 \item[(b)]
          for \(\operatorname {H}^{n-1}(Y',\mathbb {Z})\), \((p,q)\in [n-N,N-1]\times [n-N,N-1]\)
        


        
 \item[(c)]
          for \(\operatorname {H}^n(X'\operatorname {mod} Y',\mathbb {Z})=\operatorname {H}^n(X\operatorname {mod} Y,\mathbb {Z})\), \((p,q)\in [n-N,N]\times [n-N,N]\)
        


        
 \item[(d)]
          for \(\operatorname {H}^n(Y,\mathbb {Z})\), \((p,q)\in [n-N+1,N-1]\times [n-N+1,N-1]\)
        


        
 \item[(e)]
          for \(\operatorname {H}^n(X,\mathbb {Z})\), \((p,q)\in [n-N,N]\times [n-N,N]\)
        

      \end{enumerate}


      (a) follows from the fact that \(X'\) is smooth, and from the fact that the claim is true for the \({}_WE_1\) pages of the spectral sequence \hyperref[hodge-theory-ii-3.2.4.1]{Equation~\ref*{hodge-theory-ii-3.2.4.1}}.
      (b) and (d) follow from the inductive hypothesis: we have that \(\dim  Y,\dim  Y'\leqslant  N-1\).
      (c) follows from (a) and (b), and (e) follows from (c) and (d).
      This finishes the proof.\end{proof}\chapter{Examples and applications}\thispagestyle{fancy}\label{hodge-theory-iii-9}\section{Cohomology of groups and of classifying spaces}\label{hodge-theory-iii-9.1}\par{}Let \(G\) be a complex connected linear algebraic group.
  We intend to calculated the mixed Hodge structure of \(\operatorname {H}^\bullet (G)\).
  We will start by calculating that of \(\operatorname {H}^\bullet (\mathrm {B}_{\bullet  G})\), by reduction to the case of tori;
  we will then calculate that of \(\operatorname {H}^\bullet (G)\) via the Eilenberg–Moore spectral sequence \cref{hodge-theory-iii-6.1.5}.
  As a bonus, we obtain the strange Theorem \hyperref[hodge-theory-iii-9.1.7]{Theorem~\ref*{hodge-theory-iii-9.1.7}}.\begin{theorem}{9.1.1}\setnumber{9.1.1}\label{hodge-theory-iii-9.1.1}\skipsinglepar\par{}Let \(G\) be a linear algebraic group.
    We have that \(\operatorname {H}^{2n+1}(\mathrm {B}_{\bullet  G},\mathbb {Q})=0\), and that \(\operatorname {H}^{2n}(\mathrm {B}_{\bullet  G},\mathbb {Q})\otimes \mathbb {C}\) is pure of type \((n,n)\) (\cref{hodge-theory-iii-0.5})\begin{proof}\par{}Let \(G^0\) be the identity component of \(G\), and let \(T\subset  G^0\) be a maximal torus.
      Since \(\operatorname {H}^\bullet (\mathrm {B}_{\bullet  G},\mathbb {Q})=\operatorname {H}^\bullet (\mathrm {B}_{\bullet  G^0},\mathbb {Q})^{G/G^0}\), we have, by the \emph{splitting principle} (see \cref{hodge-theory-iii-6.1.6})
      \[         \operatorname {H}^\bullet (\mathrm {B}_{\bullet  G},\mathbb {Q})         \hookrightarrow  \operatorname {H}^\bullet (\mathrm {B}_{\bullet  T},\mathbb {Q}).       \]
      This inclusion is a morphism of mixed Hodge structures: it suffices to consider the case of a torus, or even, by the Künneth formula, the case where \(G=\mathbb {G}_\mathrm {m}\).
      Here are two proofs of \hyperref[hodge-theory-iii-9.1.1]{Theorem~\ref*{hodge-theory-iii-9.1.1}} in this case.\par{}\emph{1st proof.} —
      The cohomology of \(\mathbb {G}_\mathrm {m}\) (i.e. of the projective line minus two points) is \(\operatorname {H}^0=\operatorname {H}^1=\mathbb {Z}\).
      The only Hodge numbers \(h^{pq}\) of \(\operatorname {H}^1(\mathbb {G}_\mathrm {m})\) that can be non-zero are \(h^{01}=h^{10}\) and \(h^{11}\) ((i) and (iv) of \hyperref[hodge-theory-iii-8.2.4]{Theorem~\ref*{hodge-theory-iii-8.2.4}}).
      Since \(h^{01}+h^{10}+h^{11}=2h^{01}+h^{11}=1\), we have that \(h^{01}=0\), and so \(\operatorname {H}^1(\mathbb {G}_\mathrm {m})\) is pure of type \((1,1)\) (for another proof, see \hyperref[hodge-theory-iii-10.3.8]{Construction~\ref*{hodge-theory-iii-10.3.8}}).
      The \(E_2^{p,q}\) page of the spectral sequence \cref{hodge-theory-iii-6.1.5}, which is zero for odd \(p+q\), is thus pure of type \(((p+q)/2,(p+q)/2)\).
      Its abutment has the same homogeneity.\par{}\emph{2nd proof.} —
      We know that \(\mathbb {P}^r(\mathbb {C})\), endowed with \(\mathcal {O}(1)\), is an "approximation" of \(\mathrm {B}_{\mathbb {G}_\mathrm {m}}\): the homomorphism \([\mathcal {O}(1)]\colon \operatorname {H}^i(\mathrm {B}_{\bullet \mathbb {G}_\mathrm {m}})\to \operatorname {H}^i(\mathbb {P}^r(\mathbb {C}))\) is bijective for \(i\leqslant  2r\).
      The cohomology group \(\operatorname {H}^{2i}(\mathbb {P}^r(\mathbb {C}))\) is generated by the cohomology class of an algebraic cycle (that is, a linear subspace), and is thus of type \((i,i)\).
      The following proposition allows us to deduce that \(\operatorname {H}^{2i}(\mathrm {B}_{\mathbb {G}_\mathrm {m}})\) is of type \((i,i)\).\end{proof}\end{theorem}\begin{proposition}{9.1.2}\setnumber{9.1.2}\label{hodge-theory-iii-9.1.2}\skipsinglepar\par{}Let \(G\) be an algebraic group (over \(\mathbb {C}\)), and \(P\) a (algebraic) principal homogeneous space of group \(G\) over a scheme \(X\).
    Then the morphism
    \[       [P]\colon \operatorname {H}^\bullet (\mathrm {B}_G)=\operatorname {H}^\bullet (\mathrm {B}_{\bullet  G})       \to  \operatorname {H}^\bullet (X)     \]
    is a morphism of mixed Hodge structures.\begin{proof}\par{}This follows from the definition (\hyperref[hodge-theory-iii-6.1.3.1]{Equation~\ref*{hodge-theory-iii-6.1.3.1}}) of \([P]\) and from functoriality (\hyperref[hodge-theory-iii-8.3.4]{Definition~\ref*{hodge-theory-iii-8.3.4}}).\end{proof}\end{proposition}\begin{corollary}{9.1.3}\setnumber{9.1.3}\label{hodge-theory-iii-9.1.3}\skipsinglepar\par{}The characteristic classes, in \(\operatorname {H}^{2n}(S)\), of a (algebraic) principal homogeneous space \(P\) over \(S\), with structure group a linear group \(G\), are pure of type \((n,n)\): they come from morphisms \(\mathbb {Q}(-n)\to \operatorname {H}^{2n}(S,\mathbb {Q})\).\end{corollary}\par{}In the particular case where \(S\) is smooth and quasi-projective, we know that the characteristic classes of \(P\) are even algebraic, i.e. in the image of the Chow ring \((\otimes \mathbb {Q})\).\label{hodge-theory-iii-9.1.4}\par{}Let \(G\) be a linear algebraic group over a finite field \(\mathbb {F}_q\), and let \(\overline {G}/\overline {\mathbb {F}}_q\) be the group induced by extension of scalars to an algebraic closure of \(\mathbb {F}_q\).
    A proof parallel to that of \hyperref[hodge-theory-iii-9.1.1]{Theorem~\ref*{hodge-theory-iii-9.1.1}} shows that, in \(\ell \)-adic cohomology, the proper values of the (geometric) Frobenius acting on \(\operatorname {H}^{2n}(\mathrm {B}_{\bullet \overline {G}},\mathbb {Q}_\ell )\) are of the form \(q^n\zeta \) (with \(\zeta \) a root of unity).
    This agrees with the philosophy of \cref{hodge-theory-i-3}.
    Similarly, the statement of \hyperref[hodge-theory-iii-9.1.5]{Theorem~\ref*{hodge-theory-iii-9.1.5}} below is inspired by the formulas that give the number of points of reductive groups over finite fields.\begin{theorem}{9.1.5}\setnumber{9.1.5}\label{hodge-theory-iii-9.1.5}\skipsinglepar\par{}Let \(G\) be a connected linear algebraic group.
    Let \(P^\bullet \subset \operatorname {H}^\bullet (G,\mathbb {Q})\) be the primitive part of the Hopf algebra \(\operatorname {H}^\bullet (G,\mathbb {Q})\).
    It is a mixed Hodge substructure.
    We have that \(P^{2i}(G)=0\), and \(P^{2i-1}(G)\) is pure of type \((i,i)\).
    Finally, the isomorphism
    
    \begin{eqn}{9.1.5.1}\setnumber{9.1.5.1}\label{hodge-theory-iii-9.1.5.1}\skipsinglepar\[         \bigwedge  P^\bullet (G)         \to  \operatorname {H}^\bullet (G,\mathbb {Q})       \tag{9.1.5.1}       \]\end{eqn}

    is an isomorphism of mixed Hodge structures.\begin{proof}\par{}In the spectral sequence \cref{hodge-theory-iii-6.1.5}, \(E^{0,q}=0\) for \(q>0\).
      For \(n>0\), the \emph{edge homomorphism}
      \[         \operatorname {H}^{2n}(\mathrm {B}_G,\mathbb {Q})         \to  E_1^{1,2n}         = \operatorname {H}^{2n-1}(G,\mathbb {Q})       \]
      is a morphism of mixed Hodge structures (\hyperref[hodge-theory-iii-8.3.5]{Proposition~\ref*{hodge-theory-iii-8.3.5}});
      its image \(P^{2n-1}(G)\) is thus a mixed Hodge substructure, and is pure of type \((n,n)\) by \hyperref[hodge-theory-iii-9.1.1]{Theorem~\ref*{hodge-theory-iii-9.1.1}}.
      The isomorphism \hyperref[hodge-theory-iii-9.1.5.1]{Equation~\ref*{hodge-theory-iii-9.1.5.1}} is given by the cup product, and we conclude by \hyperref[hodge-theory-iii-8.2.11]{Corollary~\ref*{hodge-theory-iii-8.2.11}}.\end{proof}\end{theorem}\begin{corollary}{9.1.6}\setnumber{9.1.6}\label{hodge-theory-iii-9.1.6}\skipsinglepar\par{}Let \(G\) be a linear algebraic group.
    For \(n>0\), we have that
    \[       W_n(\operatorname {H}^n(G,\mathbb {Q}))       = 0.     \]\begin{proof}\par{}We can reduce to the case where \(G\) is connected, and then apply \hyperref[hodge-theory-iii-9.1.5]{Theorem~\ref*{hodge-theory-iii-9.1.5}}.\end{proof}\end{corollary}\begin{theorem}{9.1.7}\setnumber{9.1.7}\label{hodge-theory-iii-9.1.7}\skipsinglepar\par{}Let \(X\) be a smooth and proper algebraic space on which a connected linear algebraic group \(G\) acts via \(\rho \colon  G\times  X\to  X\).
    Then the diagram
    
  \begin {tikzcd}
      \operatorname {H}^\bullet (X,\mathbb {Q})
        \ar [r,"\rho ^*"]
        \ar [d,equals]
      & \operatorname {H}^\bullet (X,\mathbb {Q})\otimes \operatorname {H}^\bullet (G,\mathbb {Q})
    \\\operatorname {H}^\bullet (X,\mathbb {Q})\otimes _\mathbb {Q}\mathbb {Q}
        \ar [r,equals]
      & \operatorname {H}^\bullet (X,\mathbb {Q})\otimes \operatorname {H}^0(G,\mathbb {Q})
        \ar [u]
    \end {tikzcd}

    commutes.\begin{proof}\par{}Since the composite map
      \[         \operatorname {H}^\bullet (X,\mathbb {Q})         \xrightarrow {p^*} \operatorname {H}^\bullet (X,\mathbb {Q})\otimes \operatorname {H}^\bullet (G,\mathbb {Q})         \to  \operatorname {H}^\bullet (X,\mathbb {Q})\otimes \operatorname {H}^0(G,\mathbb {Q})         = \operatorname {H}^\bullet (X,\mathbb {Q})       \]
      is the identity (the composite \(X=\{e\}\times  X\to  G\times  X\to  X\) is the identity), it suffices to prove that, for \(i>0\), the \(i\)-th Künneth component of \(\rho ^*\)
      \[         \rho _i^*\colon  \operatorname {H}^n(X,\mathbb {Q})         \to  \operatorname {H}^{n-i}(X,\mathbb {Q})\otimes \operatorname {H}^i(G,\mathbb {Q})       \]
      is zero.
      Since \(\operatorname {H}^n(X,\mathbb {Q})\) is pure of weight \(n\), and \(\operatorname {H}^{n-i}(X,\mathbb {Q})\) is pure of weight \(n-i\), we have that
      \[         \operatorname {Im}(\operatorname {H}^n(X,\mathbb {Q}))         \subset  W_n(\operatorname {H}^{n-i}(X,\mathbb {Q})\otimes \operatorname {H}^i(G,\mathbb {Q}))         = \operatorname {H}^{n-i}(X,\mathbb {Q})\otimes  W_i\operatorname {H}_i(G,\mathbb {Q}).       \]
      By \hyperref[hodge-theory-iii-9.1.6]{Corollary~\ref*{hodge-theory-iii-9.1.6}}, this image is thus null.\end{proof}\end{theorem}\begin{variant}{9.1.8}\setnumber{9.1.8}\label{hodge-theory-iii-9.1.8}\skipsinglepar\par{}Let \(X\) be a separated algebraic space on which a connected linear group \(G\) acts.

    
 \begin{enumerate}
      
 \item[(i)]
        If \(X\) is smooth, and \(\overline {X}\) is a compactification of \(X\), then the diagram

        
  \begin {tikzcd}
          \operatorname {H}^\bullet (\overline {X},\mathbb {Q})
            \ar [r]
            \ar [d]
          & \operatorname {H}^\bullet (X,\mathbb {Q})
            \ar [r,"\rho ^*"]
          & \operatorname {H}^\bullet (X,\mathbb {Q})\otimes \operatorname {H}^\bullet (G,\mathbb {Q})
        \\\operatorname {H}^\bullet (X,\mathbb {Q})
            \ar [rr,equals]
          && \operatorname {H}^\bullet (X,\mathbb {Q})\otimes \operatorname {H}^0(G,\mathbb {Q})
            \ar [u,hook]
        \end {tikzcd}

        
        commutes.
      


      
 \item[(ii)]
        If \(X\) is proper, and \(\widetilde {X}\) is a resolution of singularities of \(X\), then the diagram

        
  \begin {tikzcd}
          \operatorname {H}^\bullet (X,\mathbb {Q})
            \ar [r,"\rho ^*"]
            \ar [d]
          & \operatorname {H}^\bullet (X,\mathbb {Q})\otimes \operatorname {H}^\bullet (G,\mathbb {Q})
            \ar [r]
          & \operatorname {H}^\bullet (\widetilde {X},\mathbb {Q})\otimes \operatorname {H}^\bullet (G,\mathbb {Q})
        \\\operatorname {H}^\bullet (\widetilde {X},\mathbb {Q})
            \ar [rr,equals]
          && \operatorname {H}^\bullet (\widetilde {X},\mathbb {Q})\otimes \operatorname {H}^0(G,\mathbb {Q})
            \ar [u,hook]
        \end {tikzcd}

        
        commutes.
      

    \end{enumerate}\begin{proof}\par{}It suffices to prove that, for \(i>0\), the Künneth components
      \[         \begin {gathered}           \operatorname {H}^n(\overline {X},\mathbb {Q})           \to  \operatorname {H}^n(X,\mathbb {Q})           \to  \operatorname {H}^{n-i}(X,\mathbb {Q})\otimes \operatorname {H}^i(G,\mathbb {Q})         \\(\text {resp.}\quad            \operatorname {H}^n(X,\mathbb {Q})           \to  \operatorname {H}^{n-i}(X,\mathbb {Q})\otimes \operatorname {H}^i(G,\mathbb {Q})           \to  \operatorname {H}^{n-i}(\widetilde {X},\mathbb {Q})\otimes \operatorname {H}^i(G,\mathbb {Q})           )         \end {gathered}       \]
      are zero.
      In case (i) (resp. case (ii)), we have (\hyperref[hodge-theory-iii-8.2.4]{Theorem~\ref*{hodge-theory-iii-8.2.4}})
      \[         \begin {gathered}           W_n\operatorname {H}^n(\overline {X},\mathbb {Q})           = \operatorname {H}^n(\overline {X},\mathbb {Q})         \\(\text {resp.}\quad            W_n\operatorname {H}^n(X,\mathbb {Q})           = \operatorname {H}^n(X,\mathbb {Q})           )         \end {gathered}       \]
      and so, by \hyperref[hodge-theory-iii-8.2.4]{Theorem~\ref*{hodge-theory-iii-8.2.4}} and \hyperref[hodge-theory-iii-9.1.6]{Corollary~\ref*{hodge-theory-iii-9.1.6}},
      \[         \begin {gathered}           W_n(\operatorname {H}^{n-i}(X,\mathbb {Q})\otimes \operatorname {H}^i(G,\mathbb {Q}))           = 0         \\(\text {resp.}\quad            W_n(\operatorname {H}^{n-i}(\widetilde {X},\mathbb {Q})\otimes \operatorname {H}^i(G,\mathbb {Q}))           = 0           )         \end {gathered}       \]
      whence the claim.\end{proof}\end{variant}\section{Hodge theory of smooth hypersurfaces, following Griffiths}\label{hodge-theory-iii-9.2}\par{}This section depends only on \cref{hodge-theory-ii-1} to \cref{hodge-theory-ii-3}.
  We will explain Theorem (8.6) of \hyperref[G1969]{[G1969]}.\label{hodge-theory-iii-9.2.1}\par{}Let \(X\) be a smooth and proper algebraic variety, and \(Y\) a smooth divisor in \(X\).
    Let \(j\) be the inclusion of \(U=X\setminus  Y\) into \(X\).
    We have that \(\mathrm {R}^i j_*\mathbb {Z}=0\) for \(i\neq 0,1\), and the Leray spectral sequence for \(j\) can be identified with the long exact sequence of cohomology

    \begin{eqn}{9.2.1.1}\setnumber{9.2.1.1}\label{hodge-theory-iii-9.2.1.1}\skipsinglepar\[         \ldots          \to  \operatorname {H}_Y^n(X)         \to  \operatorname {H}^n(X)         \to  \operatorname {H}^n(U)         \xrightarrow {\partial } \ldots        \tag{9.2.1.1}       \]\end{eqn}

    with
    \[       \operatorname {H}_Y^n(X)       = \operatorname {H}^{n-2}(Y,\operatorname {\mathscr {H}}_Y^2).     \]\par{}For a reason that was made clear in \hyperref[hodge-theory-ii-3.1.9]{Proposition~\ref*{hodge-theory-ii-3.1.9}}, we identify \(\operatorname {\mathscr {H}}_Y^2\) with the constant sheaf \(\mathbb {Z}(-1)_\mathbb {Z}\).
    Since the Leray spectral sequence for \(j\) converges, by definition, to the mixed Hodge structure of \(\operatorname {H}^n(U)\), we find that the exact sequence

    \begin{eqn}{9.2.1.2}\setnumber{9.2.1.2}\label{hodge-theory-iii-9.2.1.2}\skipsinglepar\[         \ldots          \to  \operatorname {H}^{n-2}(Y,\mathbb {Z})(-1)         \to  \operatorname {H}^n(X,\mathbb {Z})         \to \operatorname {H}^n(U,\mathbb {Z})         \xrightarrow {\partial } \ldots        \tag{9.2.1.2}       \]\end{eqn}

    is an exact sequence of mixed Hodge structures.\par{}The Leray spectral sequence for \(j\) is also the spectral sequence of the filtered complex \((\Omega _X^\bullet (\log  Y),W)\);
    the exact sequence above can thus also be identified with that induced by the short exact sequence

    \begin{eqn}{9.2.1.3}\setnumber{9.2.1.3}\label{hodge-theory-iii-9.2.1.3}\skipsinglepar\[         0         \to  \Omega _X^\bullet          \to  \Omega _X^\bullet (\log  Y)         \xrightarrow {\operatorname {Res}} \Omega _Y^\bullet          \to  0.       \tag{9.2.1.3}       \]\end{eqn}

    For this reason, the homomorphism \(\partial \) in \hyperref[hodge-theory-iii-9.2.1.2]{Equation~\ref*{hodge-theory-iii-9.2.1.2}} is also called the \emph{residue operator}.\label{hodge-theory-iii-9.2.2}\par{}The residue operator

    \begin{eqn}{9.2.2.1}\setnumber{9.2.2.1}\label{hodge-theory-iii-9.2.2.1}\skipsinglepar\[         \operatorname {Res}\colon          \operatorname {H}^{n+1}(U,\mathbb {Z})         \to  \operatorname {H}^n(Y,\mathbb {Z})(-1)       \tag{9.2.2.1}       \]\end{eqn}

    is strictly compatible with the Hodge filtration.
    The Hodge filtration of \(\operatorname {H}^n(U,\mathbb {C})\) thus determines that of \(\operatorname {Ker}(\operatorname {H}^n(Y,\mathbb {C}))\to \operatorname {H}^{n+2}(X,\mathbb {C})\).\par{}Suppose that \(Y\) is ample.
    If \(X\) is pure of dimension \(N+1\), then \(\operatorname {H}^i(X,\mathbb {Z})\xrightarrow {\sim }\operatorname {H}^i(Y,\mathbb {Z})\) for \(i<N\) (Lefschetz).
    By Poincaré duality, for \(i\neq  N\), \(\operatorname {H}^i(Y,\mathbb {C})\) can thus be calculated from \(\operatorname {H}^\bullet (X,\mathbb {C})\).
    Set
    \[       \begin {aligned}         \operatorname {H}_\mathrm {ev}^N(Y,\mathbb {C})         &= \operatorname {Ker}(\operatorname {H}^N(Y,\mathbb {C})\to \operatorname {H}^{N+2}(X,\mathbb {C}))       \\&= \text {orthogonal complement of }\operatorname {H}^N(X,\mathbb {C})\hookrightarrow \operatorname {H}^N(Y,\mathbb {C})       \end {aligned}     \]
    (the evanescent part of the cohomology).
    We have that
    \[       \operatorname {H}^N(Y,\mathbb {C})       = \operatorname {H}_\mathrm {ev}^N(Y,\mathbb {C})\oplus \operatorname {H}^N(X,\mathbb {C})     \]
    and we thus obtain a lot of information about the Hodge structure of \(\operatorname {H}^\bullet (Y,\mathbb {C})\) when we know about that of \(\operatorname {H}^\bullet (X,\mathbb {C})\) and the Hodge filtration of \(\operatorname {H}^{N+1}(U,\mathbb {C})\).\label{hodge-theory-iii-9.2.3}\par{}By definition, the Hodge filtration of \(\operatorname {H}^\bullet (U,\mathbb {C})\) is the abutment of the spectral sequence that \emph{degenerates} at \(E_1\)
    \[       E_1^{pq}       = \operatorname {H}^q(X,\Omega _X^p(\log  Y))       \Rightarrow  \operatorname {H}^{p+q}(U,\mathbb {C}).     \]
    By \hyperref[hodge-theory-ii-3.1.11]{Proposition~\ref*{hodge-theory-ii-3.1.11}}, this spectral sequence agrees with the spectral sequence of the complex \(j_*\Omega _U^\bullet \) endowed with the filtration given by the order of the pole \(P\).\par{}In particular, we have that
    \[       F^p(\operatorname {H}^\bullet (U,\mathbb {C}))       = \operatorname {\mathbb {H}}^\bullet (X,P^p(j_*\Omega _U^\bullet )).     \]
    Recall that \(P^p(j_*\Omega _U^\bullet )\) is the following complex
    \[       \Omega _X^p(Y)       \to  \Omega _X^{p+1}(2Y)       \to  \ldots        \to  \Omega _X^i((i-p+1)Y)       \to  \ldots      \]
    with the first non-zero component placed in degree \(p\).\par{}If \(\operatorname {H}^i(X,\Omega _X^p(nY))=0\) for \(i>0\), \(n>0\), and \(p\geqslant 0\), then we simply have
    \[       \operatorname {\mathbb {H}}^\bullet (X,P^p(j_*\Omega _U^\bullet ))       = \operatorname {H}^\bullet (\Gamma (X,P^p(j_*\Omega _U^\bullet )))     \]
    whence the following result.\begin{proposition}{9.2.4}\setnumber{9.2.4}\label{hodge-theory-iii-9.2.4}\skipsinglepar\par{}Suppose that \(\operatorname {H}^i(X,\Omega _X^p(nY))=0\) for \(i>0\), \(n>0\), and \(p\geqslant 0\) (which is the case if \(Y\) is sufficiently ample).
    Then a cohomology class \(c\in \operatorname {H}^d(U,\mathbb {C})\) is of Hodge filtration \(\geqslant  p\) if and only if it can be represented by a closed \(d\)-form \(\alpha \) on \(U\), with a pole of order \(\leqslant  d-p+1\) along \(Y\).
    If \(\alpha \) is a closed form on \(U\), with a pole of order \(k\) along \(Y\), and if the class of \(\alpha \) in \(\operatorname {H}^d(U,\mathbb {C})\) is zero, then \(\alpha =d\beta \), with \(\beta \) presenting a pole of order \(\leqslant  k-1\) along \(Y\);
    for \(k\leqslant 1\), we have that \(\alpha =0\).\end{proposition}\label{hodge-theory-iii-9.2.5}\par{}It is well known (by a theorem of Bott) that the hypotheses of \hyperref[hodge-theory-iii-9.2.4]{Proposition~\ref*{hodge-theory-iii-9.2.4}} are satisfied for \(Y\) a hypersurface in the projective space \(\mathbb {P}^n(\mathbb {C})\) (for \(n\geqslant 1\)), i.e. that
    \[       \operatorname {H}^i(\mathbb {P}^n(\mathbb {C}),\Omega ^j(m))       = 0       \quad        \text {for }i>0\text { and }m>0     \]
    (a demonstration appears in [SGA 7, XI], for example).
    Furthermore, for \(n>1\), the residue map identifies \(\operatorname {H}^n(U,\mathbb {C})\) with the primitive part of the cohomology of \(\operatorname {H}^{n-1}(Y,\mathbb {C})\).
    Thus ([\hyperref[G1969]{G1969}, (8.6)]):\begin{proposition}{9.2.6}\setnumber{9.2.6}\label{hodge-theory-iii-9.2.6}\skipsinglepar\par{}Let \(Y\) be a smooth hypersurface in \(\mathbb {P}^n(\mathbb {C})\) (with \(n>1\)) and \(U=\mathbb {P}^n(\mathbb {C})\setminus  Y\).

    
 \begin{enumerate}
      
 \item[(i)]
        For a primitive cohomology class \(c\in \operatorname {H}^{n-1}(Y,\mathbb {C})\) to be of Hodge filtration \(\leq  p\), it is necessary and sufficient that it be the residue of a regular differential \(n\)-form on \(U\), with a pole of order \(\leqslant  n-p\) along \(Y\).
      


      
 \item[(ii)]
        Let \(\alpha \) be a regular differential \(n\)-form on \(U\), with a pole of order \(k\geqslant 1\) along \(Y\).
        For \(\operatorname {Res}(\alpha )\in \operatorname {H}^{n-1}(Y,\mathbb {C})\) to be zero, it is necessary and sufficient that \(\alpha =d\beta \) with \(\beta \) an \((n-1)\)-form on \(U\) with a pole of order \(\leqslant  k-1\) along \(Y\).
      

    \end{enumerate}\end{proposition}\section{Construction of complexes of first-order differential operators}\label{hodge-theory-iii-9.3}\par{}This section depends only on \cref{hodge-theory-iii-5} and \cref{hodge-theory-iii-6}.
  We will prove the following theorem.\begin{proposition}{9.3.1}\setnumber{9.3.1}\label{hodge-theory-iii-9.3.1}\skipsinglepar\par{}Let \(X\) be a quasi-projective scheme (over \(\mathbb {C}\)).
    Then there exists on \(X\) a complex \(K\) of sheaves that satisfies the following properties.

    
 \begin{enumerate}
      
 \item[(a)]
        The \(K^n\) are coherent algebraic sheaves, and \(K^n=0\) for \(n<0\).
      

      
 \item[(b)]
        The differentials \(\mathrm {d}\colon  K^n\to  K^{n+1}\) are algebraic differential operators of order 1.
      

      
 \item[(c)]
        The complex \(K^\mathrm {an}\) of coherent analytic sheaves on \(X^\mathrm {an}\) is a resolution of the constant sheaf \(\underline {\mathbb {C}}\).
        More precisely, \(\operatorname {\mathscr {H}}^i(K)=0\) for \(i>0\), and there exists
        \[           f\colon  \Omega _X^\bullet            \to  K         \]
        with \(f^n\) linear algebraic, such that the composite map \(\underline {\mathbb {C}}\to \mathcal {O}_X\xrightarrow {f^0}K^0\) identifies \(\underline {\mathbb {C}}\) with \(\operatorname {\mathscr {H}}^0(K)\).
      

      
 \item[(d)]
        The natural map
        \[           \operatorname {\mathbb {H}}^i(X,K)           \to  \operatorname {\mathbb {H}}^i(X^\mathrm {an},K^\mathrm {an})           \xleftarrow [\text {(c)}]{\sim } \operatorname {\mathbb {H}}^i(X^\mathrm {an},\mathbb {C})         \]
        is an isomorphism.
      

    \end{enumerate}\end{proposition}\par{}As a corollary, we recover a particular case of a theorem of Bloom and Herrera [\hyperref[BH1969]{BH1969}].\par{}We will first prove the corollary.
  The proposition \hyperref[hodge-theory-iii-9.3.1]{Proposition~\ref*{hodge-theory-iii-9.3.1}} will be obtained by redoing the proof at the level of complexes.\begin{corollary}{9.3.2}\setnumber{9.3.2}\label{hodge-theory-iii-9.3.2}\skipsinglepar\par{}The map
    \[       \operatorname {H}^\bullet (X^\mathrm {an},\mathbb {C})       \to  \operatorname {\mathbb {H}}^\bullet (X^\mathrm {an},\Omega _{X^\mathrm {an}}^\bullet )     \]
    induced by \(\underline {\mathbb {C}}\to \Omega _{X^\mathrm {an}}^\bullet \) identifies \(\operatorname {H}^\bullet (X^\mathrm {an},\mathbb {C})\) with a direct factor of \(\operatorname {H}^\bullet (X^\mathrm {an},\Omega _{X^\mathrm {an}}^\bullet )\).\begin{proof}\par{}Let \(\varepsilon \colon  X_\bullet \to  X\) be an augmented simplicial scheme over \(X\).
      Suppose that \(\varepsilon \) is a proper hypercover, and that the \(X_n\) are smooth.
      We have the commutative diagram

      \begin{eqn}{9.3.2.1}\setnumber{9.3.2.1}\label{hodge-theory-iii-9.3.2.1}\skipsinglepar\end{eqn}

      
      The arrow \((1)\) is bijective by (II) and (IV) of \cref{hodge-theory-iii-5.3.5}.
      The arrow \((2)\) is bijective since \(\Omega _{X_\bullet }^\bullet \) is a resolution of \(\underline {\mathbb {C}}\) (the Poincaré lemma on the \(X_n\)).
      The diagram \hyperref[hodge-theory-iii-9.3.2.1]{Equation~\ref*{hodge-theory-iii-9.3.2.1}} gives a retract, compatible with the cup product, of \(\operatorname {\mathbb {H}}^\bullet (X^\mathrm {an},\Omega _X^\bullet )\) onto \(\operatorname {H}^\bullet (X^\mathrm {an},\mathbb {C})\), and this proves \hyperref[hodge-theory-iii-9.3.2]{Corollary~\ref*{hodge-theory-iii-9.3.2}}.\end{proof}\end{corollary}\label{hodge-theory-iii-9.3.3}\par{}The complex \(K\) that we wish to construct is, in the derived category, exactly \(\mathrm {s}\mathrm {R}{\varepsilon _\bullet }_*\Omega _{X_\bullet }^\bullet \).
    We will only deal with the case where the augmented simplicial space \(X_\bullet \) is \(s\)-split, defined as in \cref{hodge-theory-iii-6.2.5} by proper and surjective maps, with smooth quasi-projective sources
    \[       f'_k\colon        N_k       \to  (\operatorname {cosq}\operatorname {sq}_{k-1}(X_\bullet ))_k     \]
    for \(k\geqslant 0\).\par{}We will construct a "compatible system" of affine open covers of the \(X_n\).
    It will be relevant, for this construction, to consider an open cover of a space \(Y\) as a diagram
    \[       Y \xleftarrow {u} U \xrightarrow {\varphi } I     \]
    where \(I\) is discrete, \(u\) is surjective, and \(u_i\colon  U_i\coloneqq \varphi ^{-1}(i)\hookrightarrow  Y\) is an open embedding.
    Proceeding as in \cref{hodge-theory-iii-6.2.5}, we construct a diagram

    \begin{eqn}{9.3.3.1}\setnumber{9.3.3.1}\label{hodge-theory-iii-9.3.3.1}\skipsinglepar\[         X_\bullet  \xleftarrow {u} U_\bullet  \xrightarrow {\varphi } I_\bullet        \tag{9.3.3.1}       \]\end{eqn}

    of the following type.

    
 \begin{enumerate}
      
 \item[(a)]
        \(I_\bullet \) is a discrete simplicial set. The \(I_n\) are finite.
      


      
 \item[(b)]
        For all \(k\geqslant  0\), \hyperref[hodge-theory-iii-9.3.3.1]{Equation~\ref*{hodge-theory-iii-9.3.3.1}} induces a diagram
        
  \begin {tikzcd}
          (\operatorname {cosq}\operatorname {sq}_{k-1}(X_\bullet ))_k
          & (\operatorname {cosq}\operatorname {sq}_{k-1}(U_\bullet ))_k
            \ar [l]
            \ar [r,"\varphi "]
          & (\operatorname {cosq}\operatorname {sq}_{k-1}(I_\bullet ))_k
        \\N_k
            \ar [u,"f'_k"]
          & N_k(U_\bullet )
            \ar [u]
            \ar [l]
            \ar [r,swap,"\varphi "]
          & N_k(I_\bullet )
            \ar [u,swap,"a'_k"]
        \end {tikzcd}

      


      
 \item[(c)]
        In the diagram in (b), for \(i\in \operatorname {cosq}\operatorname {sq}_{k-1}(I_\bullet )_k\), the \(\varphi ^{-1}(j)\) for \(a'_k(j)=i\) form an affine open cover of
        \[           N_k\times _{(\operatorname {cosq}\operatorname {sq}_{k-1}(X_\bullet ))_k} \varphi ^{-1}(i).         \]
      


      
 \item[(d)]
        (Follows by induction on (c)).
        The first (and thus the second) row of the diagram in (b) defines an open cover.
        The \(X_k\leftarrow  U_k\to  I_k\) thus also form open covers.
      


      
 \item[(e)]
        (Follows from a suitable choice of affine covers in (c)).
        \(U_{n_i}=\varphi _n^{-1}(i)\) is the complement of a subscheme \(D_{n,i}\) of \(X_n\), given by the union of a divisor and some connected components of \(X_n\).
        For every coherent sheaf \(\mathscr {F}\) on \(X_n\) and integers \((a_i)_{i\in  I_n}\) (with \((a_i\geqslant 0)\)), we set
        \[           \mathscr {F}(\textstyle \sum _i a_i D_{n,i})           =           \begin {cases}             0             &\text {on the components of }X_n\text { contained inside some }D_{n,i}\text { with }a_i>0           \\\mathscr {F}\otimes \mathcal {O}(\sum _{a_i>0}a_i D_{n,i})             &\text { on the others}.           \end {cases}         \]
      

    \end{enumerate}\par{}In the inductive construction of \hyperref[hodge-theory-iii-9.3.3.1]{Equation~\ref*{hodge-theory-iii-9.3.3.1}}, the open covers in (c) are chosen arbitrarily.
    This allows us to suppose the following.

    
 \begin{enumerate}
      
 \item[(f)]
        For any \(a_i>0\) (for \(i\in  I_k\)),
        \[           \mathrm {R}^\ell {\varepsilon _k}_*           \big ( \Omega _{X_k}^m(\textstyle \sum _i a_i D_{k,i}) \big )           = 0         \]
        for \(\ell >0\).
      


      
 \item[(g)]
        For any \(\alpha \colon \Delta _m\to \Delta _n\) and \(i\in  I_n\), we have that
        \[           X_\bullet (\alpha )^{-1}(D_{m,\alpha (i)})           \subset  D_{n,i}         \]
        (as schemes),
      

    \end{enumerate}\par{}For each \(k\), the complex of sheaves
    \[       \mathrm {R}{\varepsilon _k}_*(\Omega _{X_k}^\bullet )     \]
    on \(X\) can be calculated in a Čech manner.
    If, for \(P\subset  I_k\), \(j_P\) is the inclusion of \(U_P=\bigcap _{i\in  P}U_{k_i}\) in \(X_k\), then this is the simple complex associated to the double complex of the
    \[       \check {K}^{k,p,q}       = \sum _{\# P=p+1} {\varepsilon _k}_*\Omega _{X_k}^q\big ( (q+1)\textstyle \sum _{i\in  P} D_{k,i} \big ).     \]\par{}For fixed \(k\) and fixed \(q\), it follows from (f) that the inclusion of \(K^{k,\bullet ,q}\) into \(\check {K}^{k,\bullet ,q}\) is a quasi-isomorphism.
    For fixed \(k\), the inclusion of \(K^{k,\bullet ,\bullet }\) into \(\check {K}^{k,\bullet ,\bullet }\) is also a quasi-isomorphism.
    Finally, for varying \(k\), the \(K^{k,\bullet ,\bullet }\) form, by (g), a simplicial double complex.
    The associated simple complex \(K\) agrees with \(\mathrm {R}{\varepsilon _\bullet }_*\Omega _{X_\bullet }^\bullet \) in the derived category.\par{}We now prove that \(K\) satisfies \hyperref[hodge-theory-iii-9.3.1]{Proposition~\ref*{hodge-theory-iii-9.3.1}}.
    Properties (a) and (b) of \(K\) are clear;
    we define \(f\colon \Omega _X^\bullet \to  K\) by the evident maps from the \(\Omega _X^q\) to the \(K^{0,0,q}\).\par{}In the analytic category, the complex \(K^\mathrm {an}\) of coherent analytic sheaves is given by the same formulas as for \(K\).
    We have a commutative diagram
    
  \begin {tikzcd}
      \operatorname {\mathbb {H}}^i(X,K)
        \ar [r,"\sim "]
        \ar [d]
      & \operatorname {\mathbb {H}}^i(X_\bullet ,\Omega _{X_\bullet }^\bullet )
        \ar [d,"(1)"]
    \\\operatorname {\mathbb {H}}^i(X^\mathrm {an},K^\mathrm {an})
        \ar [r,"\sim "]
      & \operatorname {\mathbb {H}}^i(X_\bullet ^\mathrm {an},\Omega _{X_\bullet ^\mathrm {an}}^\bullet )
    \end {tikzcd}

    and \((1)\) is an abutment of a morphism of spectral sequences
    
  \begin {tikzcd}
      \operatorname {\mathbb {H}}^a(X_b,\Omega _{X_\bullet }^\bullet )
        \ar [r,Rightarrow]
        \ar [d]
      & \operatorname {\mathbb {H}}^{a+b}(X_\bullet ,\Omega _{X_\bullet }^\bullet )
        \ar [d]
    \\\operatorname {\mathbb {H}}^a(X_b^\mathrm {an},\Omega _{X_\bullet ^\mathrm {an}}^\bullet )
        \ar [r,Rightarrow]
      & \operatorname {\mathbb {H}}^{a+b}(X_\bullet ^\mathrm {an},\Omega _{X_\bullet ^\mathrm {an}}^\bullet ).
    \end {tikzcd}

    This morphism is an isomorphism by [\hyperref[G1966]{G1966}], whence (d).\par{}To prove (c), we consider the commutative diagram
    
  \begin {tikzcd}
      \underline {\mathbb {C}}
        \ar [rr,"\sim "]
        \ar [rr,swap,"(1)"]
        \ar [d]
      && \mathrm {R}{\varepsilon _\bullet }_*\underline {\mathbb {C}}
        \ar [d,"\sim "]
        \ar [d,swap,"(2)"]
    \\{\Omega _{X_\bullet }^\bullet }^\mathrm {an}
        \ar [r,swap,"f"]
      & K^\mathrm {an}
        \ar [r,"\sim "]
      & {\mathrm {R}{\varepsilon _\bullet }_*\Omega _{X_\bullet }^\bullet }^\mathrm {an}
    \end {tikzcd}\par{}The fact that \((1)\) is a quasi-isomorphism follows from \cref{hodge-theory-iii-5.3.5}.
    The fact that \((2)\) is a quasi-isomorphism follows from the Poincaré lemma on the \(X_n\).
    Claim (c) thus follows.\label{hodge-theory-iii-9.3.4}\par{}The same construction gives resolutions that satisfy (a) and (b) for sheaves \(\mathscr {F}\) on \(X^\mathrm {an}\) other than the constant sheaf \(\underline {\mathbb {C}}\) (we can probably take \(\mathscr {F}\) to be any "algebraically constructible" sheaf).
    Unfortunately, I can not prove that the resolutions thus obtained are essentially unique (in a suitable derived category).
    In the absence of such uniqueness, \hyperref[hodge-theory-iii-9.3.1]{Proposition~\ref*{hodge-theory-iii-9.3.1}} is not of great interest.\chapter{Hodge theory in level \(\leqslant  1\)}\thispagestyle{fancy}\label{hodge-theory-iii-10}\section{1-Motives}\label{hodge-theory-iii-10.1}\begin{definition}{10.1.1}\setnumber{10.1.1}\label{hodge-theory-iii-10.1.1}\skipsinglepar\par{}\begin{enumerate}
      
 \item[(i)]
        A mixed Hodge structure \(H\) is \emph{torsion free} if \(H_\mathbb {Z}\) is torsion free.
      


      
 \item[(ii)]
        If \(H\) is torsion free, then we again denote by \(W\) the filtration of \(H_\mathbb {Z}\subset \operatorname {H}_\mathbb {Q}\) induced by \(W\).
        We denote by \(\operatorname {Gr}_n^W(H)\) the torsion free \(\mathbb {Z}\)-module \(\operatorname {Gr}_n^W(H_\mathbb {Z})\) endowed with its Hodge structure of weight \(n\).
      

    \end{enumerate}\end{definition}\begin{definition}{10.1.2}\setnumber{10.1.2}\label{hodge-theory-iii-10.1.2}\skipsinglepar\par{}A \emph{1-motive} over an algebraically closed field \(k\) consists of the following:

    
 \begin{enumerate}
      
 \item[(a)]
        a free \(\mathbb {Z}\)-module \(X\) of finite type, an abelian variety \(A\), and a torus \(T\);
      

      
 \item[(b)]
        an extension \(G\) of \(A\) by \(T\);
      

      
 \item[(c)]
        a homomorphism \(u\colon  X\to  G(k)\).
      

    \end{enumerate}\end{definition}\par{}It is often useful to consider a 1-motive \(M\) as being a complex of group schemes, concentrated in degrees 0 and 1:
  \[     M = [X\xrightarrow {u}G].   \]\begin{construction}{10.1.3}\setnumber{10.1.3}\label{hodge-theory-iii-10.1.3}\skipsinglepar\par{}We will construct an equivalence \(M\mapsto  T(M)\) from the category of 1-motives over \(\mathbb {C}\) to the category of torsion free mixed Hodge structures \(H\) of type
    \[       \big \{ (0,0), (0,-1), (-1,0), (-1,-1) \big \}     \]
    
    such that \(\operatorname {Gr}_{-1}^W(H)\) is polarisable.
    Furthermore, we will construct isomorphisms
    \[       \begin {aligned}         \operatorname {H}_1(T,\mathbb {Z})         &\simeq  \operatorname {Gr}_{-2}^WT(M)_\mathbb {Z}       \\\operatorname {H}_1(A,\mathbb {Z})         &\simeq  \operatorname {Gr}_{-1}^WT(M)         \qquad \text {(isomorphism of Hodge structures)}       \\X         &\simeq  \operatorname {Gr}_0^WT(M)_\mathbb {Z}       \end {aligned}     \]
    that are natural in \(M=(X,A,T,G,u)\).\par{}Let \(M\) be a 1-motive.
    The exponential map \(\exp \colon \operatorname {Lie}(G)\to  G\) has \(\operatorname {H}_1(G)\) as its kernel.
    We define \(T_\mathbb {Z}(M)\) as the fibre product of \(\operatorname {Lie}(G)\) and \(X\) over \(G\), whence

    \begin{eqn}{10.1.3.1}\setnumber{10.1.3.1}\label{hodge-theory-iii-10.1.3.1}\skipsinglepar\[         \begin {CD}           0 @>>> \operatorname {H}_1(G) @>>> \operatorname {Lie}(G) @>{\exp }>> G @>>> 0         \\@. @| @A{\alpha }AA @AA{u}A @.         \\0 @>>> \operatorname {H}_1(G) @>>> T_\mathbb {Z}(M) @>>{\beta }> X @>>> 0         \end {CD}       \tag{10.1.3.1}       \]\end{eqn}\par{}We set
    \[       \begin {aligned}         W_{-1}(T_\mathbb {Z}(M))         &= \operatorname {Ker}(\beta )         = \operatorname {H}_1(G)       \\W_{-2}(T_\mathbb {Z}(M))         &= \operatorname {H}_1(T)         = \operatorname {Ker}\big (\operatorname {H}_1(G)\to \operatorname {H}_1(A)\big ).       \end {aligned}     \]
    This defines the \emph{weight filtration}.
    The morphism \(\alpha \) extends to
    \[       \alpha _\mathbb {C}\colon  T_\mathbb {Z}(M)\otimes \mathbb {C}       \to  \operatorname {Lie}(G).     \]
    We set \(F^0(T_\mathbb {Z}(M)\otimes \mathbb {C})=\operatorname {Ker}(\alpha _\mathbb {C})\), and this defines the \emph{Hodge filtration} on
    \[       T_\mathbb {C}(M)       = T_\mathbb {Z}(M)\otimes \mathbb {C}.     \]\begin{lemma}{10.1.3.2}\setnumber{10.1.3.2}\label{hodge-theory-iii-10.1.3.2}\skipsinglepar\par{}The triple \(T(M)=(T_\mathbb {Z}(M),W,F)\) is a torsion-free mixed Hodge structure of type \(\{(0,0),(0,-1),(-1,0),(-1,-1)\}\), and \(\operatorname {Gr}_{-1}^W(T(M))\) is polarisable.\begin{proof}\par{}\begin{enumerate}
          
 \item[(a)]
            We have that \(W_{-2}(T_\mathbb {C}(M))\cap  F^0=0\), since \((1)\) is injective;
            thus \((W_{-2}(T_\mathbb {C}(M)),F)\) is a Hodge structure of type \((-1,-1)\).
          


          
 \item[(b)]
            Since \((2)\) is surjective, \(F\) induces on \(\operatorname {Gr}_{-1}^W(T_\mathbb {Z}(M))\otimes \mathbb {C}\simeq \operatorname {H}_1(A,\mathbb {Z})\otimes \mathbb {C}\) the Hodge filtration of \(\operatorname {H}_1(A,\mathbb {Z})\otimes \mathbb {C}\);
            this is a polarisable Hodge structure of type \((-1,0)+(0,1)\) (cf. \hyperref[hodge-theory-ii-4.4.3]{Reminder~\ref*{hodge-theory-ii-4.4.3}}).
          


          
 \item[(c)]
            Consider the diagram
            
  \begin {tikzcd}[row sep=huge]
              & 0
                \ar [d]
              & 0
                \ar [d]
            \\0
                \ar [r]
            & W_{-1}(T_\mathbb {C}(M))\cap  F^0
                \ar [r]
                \ar [d]
            & F^0
                \ar [r]
                \ar [d]
            & V
                \ar [r]
                \ar [d,"\sim "]
            & 0
            \\0
                \ar [r]
            & \operatorname {H}_1(G,\mathbb {Z})\otimes \mathbb {C}
                \ar [r]
                \ar [d]
            & T_\mathbb {Z}(M)\otimes \mathbb {C}
                \ar [r]
                \ar [d]
            & X\otimes \mathbb {C}
            & 0
            \\
            & \operatorname {Lie}(G)
                \ar [r,equals]
                \ar [d]
            & \operatorname {Lie}(G)
                \ar [d]
            \\
            & 0
            & 0
            \end {tikzcd}

            Since \(F^0\) is sent to \(\operatorname {Gr}_0^W(T_\mathbb {C}(M))\), it is thus of type \((0,0)\).
          

        \end{enumerate}\par{}This finishes the construction of \(T(M)\).
        It is clear that \(T(M)\) is functorial in \(M\).\end{proof}\end{lemma}\par{}Suppose that \(\operatorname {Gr}_{-1}^W(H)\) is polarisable;
    the complex torus
    \[       A       = H_\mathbb {Z} \backslash  \operatorname {Gr}_{-1}^W(H_\mathbb {C}) / F^0\operatorname {Gr}_{-1}^W(H_\mathbb {C})     \]
    is then an abelian variety (cf. \hyperref[hodge-theory-ii-4.4.3]{Reminder~\ref*{hodge-theory-ii-4.4.3}}).
    Let \(T\) be the torus with character group the dual of \(\operatorname {Gr}_{-2}^W(H_\mathbb {Z})\), i.e.
    \[       \operatorname {H}_1(T)       = \operatorname {Gr}_{-2}^W(H_\mathbb {Z}).     \]
    The complex analytic group
    \[       G       = W_{-1}(H_\mathbb {Z}) \backslash  W_{-1}(H_\mathbb {C}) / F^0\cap  W_{-1}(H_\mathbb {C})     \]
    is an extension of \(A\) by \(T\).\begin{lemma}{10.1.3.3}\setnumber{10.1.3.3}\label{hodge-theory-iii-10.1.3.3}\skipsinglepar\par{}The functor \(E\mapsto  E^\mathrm {an}\) is an equivalence from the category of algebraic group extensions of \(A\) by \(T\) to the category of analytic group extensions of \(A\) by \(T\).\begin{proof}\par{}We reduce to the case where \(T=\mathbb {G}_\mathrm {m}\).
        By GAGA, \(E\mapsto  E^\mathrm {an}\) is more general an equivalence from the category of \(\mathbb {G}_\mathrm {m}\)-torsors over \(A\) to the category of \(\mathbb {G}_\mathrm {m}\)-torsors over \(A^\mathrm {an}\).\par{}We thus have \(G\), an extension of \(A\) by \(T\).
        Let \(X=\operatorname {Gr}_0^W(H_\mathbb {Z})\).
        We similarly define \(u\colon  X\to  G\) (cf. \cref{hodge-theory-ii-2.2.1}):
        
  \begin {tikzcd}
          0
            \ar [r]
        & W_{-1}(H_\mathbb {Z})
            \ar [r]
            \ar [dd,equals]
        & W_{-1}(H_\mathbb {C})/(W_{-1}(H_\mathbb {C})\cap  F^0)
            \ar [r]
            \ar [d,"\sim "]
        & G
            \ar [r]
        & 0
        \\
        &
        & H_\mathbb {C}{/}F^0
        &
        &
        \\0
            \ar [r]
        & W_{-1}(H_\mathbb {Z})
            \ar [r]
        & H_\mathbb {Z}
            \ar [r]
            \ar [u]
        & X
            \ar [r]
            \ar [uu]
        & 0
        \end {tikzcd}\par{}It is clear that the above construction, which associates to \(H\) a torsion-free 1-motive, is the inverse of \(T\), and that it is functorial.\end{proof}\end{lemma}\par{}If \(G_i\) (for \(i=1,2\)) is an extension of an abelian variety \(A_i\) by a torus \(A_i\), then, for every morphism \(u\colon  G_1\to  G_2\), we have that \(u(T_1)\subset  T_2\) and \(\operatorname {Ker}(u\colon  A_1\to  A_2)^0=\operatorname {Im}(\operatorname {Ker}(u\colon  G_1\to  G_2)\to  A_1)^0\).
    This corresponds to the fact that a morphism of mixed Hodge \(\mathbb {Q}\)-structures is strictly compatible with the weight filtration.\end{construction}\label{hodge-theory-iii-10.1.4}\par{}Let \(H\) be a torsion-free mixed Hodge structure of type
    \[       \big \{(0,0),(0,-1),(-1,0),(-1,-1)\big \}.     \]
    The filtration \(W\) of \(H_\mathbb {Z}\) thus defines a filtration on \(H\) by mixed Hodge substructures.
    At the same time, if \(M=(X,A,T,G,u)\) is a 1-motive, then we denote by \(W\) the following increasing filtration on \(M\):
    \[       \begin {cases}         W_i(M) = 0         & \text {for }i< -2;       \\W_i(M) = M         & \text {for }i\geqslant 0;       \\W_{-1}(M) = G         \quad \text {i.e. } = (\{0\},A,T,G,0);       \\W_{-2}(M) = T         \quad \text {i.e. } = (\{0\},0,T,T,0).       \end {cases}     \]
    In an evident way, the successive \(\operatorname {Gr}_i^W(M)\) for \(i=0,-1,-2\) are \(X\), \(A\), and \(T\) (respectively).\label{hodge-theory-iii-10.1.5}\par{}The construction in \hyperref[hodge-theory-iii-10.1.3]{Construction~\ref*{hodge-theory-iii-10.1.3}} that associates to a 1-motive \(M\) over \(\mathbb {C}\) the object \(T_\mathbb {Z}(M)\) is transcendental.
    We will show that
    \[       \widehat {T}(M)       = T_\mathbb {Z}(M)\otimes \widehat {\mathbb {Z}}       = \prod _\ell  T_\mathbb {Z}(M)\otimes \mathbb {Z}_\ell      \]
    can be defined in a purely algebraic way.\par{}Let \(M\) be a 1-motive over an algebraically closed field \(k\) of characteristic \(0\).
    We identify \(M\) with a complex \([X\xrightarrow {u}G]\) (concentrated in degrees \(0\) and \(1\)).
    For every integer \(n>0\), consider the complex \([\mathbb {Z}\xrightarrow {\cdot  n}\mathbb {Z}]\) (concentrated in degrees \(-1\) and \(0\)), and let \(T_{\mathbb {Z}{/}n\mathbb {Z}}(M)\) be the \(\operatorname {H}^0\) of the complex \([X\xrightarrow {u}G]\otimes [\mathbb {Z}\xrightarrow {\cdot  n}\mathbb {Z}]\), i.e.
    
    
  \begin {tikzcd}
      X
        \ar [r,"u"]
    & G
    \\X
        \ar [u,"n"]
        \ar [r,swap,"u"]
    & G
        \ar [u,swap,"-n"]
    \end {tikzcd}

    We have that
    \[       T_{\mathbb {Z}{/}n\mathbb {Z}}(M)       = \big \{ (x,g) \mid  u(x)=ng \big \} / \big \{ (nx,u(x)) \mid  x\in  X \big \}.     \]
    In the derived category, we have that
    \[       T_{\mathbb {Z}{/}n\mathbb {Z}}(M)       = \operatorname {H}^0(M\overset {\mathrm {L}}{\otimes }\mathbb {Z}{/}n\mathbb {Z}).     \]
    For \(n=md\), we define transition morphisms
    \[       \begin {aligned}         \varphi _{m,n}\colon  T_{\mathbb {Z}{/}n\mathbb {Z}}(M)         &\to  T_{\mathbb {Z}{/}m\mathbb {Z}}(M)       \\(x,g)         &\mapsto  (x,dg).       \end {aligned}     \]
    We set
    \[       \widehat {T}(M)       = \varprojlim  T_{\mathbb {Z}{/}n\mathbb {Z}}(M).     \]
    The filtration \(W\) of \(M\) induces a filtration \(W\) of the \(T_{\mathbb {Z}{/}n\mathbb {Z}}(M)\) and of \(\widehat {T}(M)\).
    We have that
    \[       \begin {aligned}         \operatorname {Gr}_0^W(\widehat {T}(M))         &= X\otimes \widehat {\mathbb {Z}}       \\\operatorname {Gr}_{-1}^W(\widehat {T}(M))         &= \varprojlim  A_n         = \widehat {T}(A)       \\\operatorname {Gr}_{-2}^W(\widehat {T}(M))         &= \varprojlim  T_n         = Y\otimes \widehat {\mathbb {Z}}(1),       \end {aligned}     \]
    where \(Y\) is the dual of the character group of \(T=W_{-2}(M)\).\begin{construction}{10.1.6}\setnumber{10.1.6}\label{hodge-theory-iii-10.1.6}\skipsinglepar\par{}Let \(M=[X\xrightarrow {u}G]\) be a 1-motive on \(\mathbb {C}\).
    We have that

    \begin{eqn}{10.1.6.1}\setnumber{10.1.6.1}\label{hodge-theory-iii-10.1.6.1}\skipsinglepar\[         \widehat {T}(M)         \simeq  T_\mathbb {Z}(M)\otimes \widehat {\mathbb {Z}}.       \tag{10.1.6.1}       \]\end{eqn}\par{}The natural morphism \([T_\mathbb {Z}(M)\to \operatorname {Lie}(G)]\to [X\to  G]\) is a quasi-isomorphism.
    The quasi-isomorphisms
    \[       \left [         \,\,         \begin {CD}           X @>{u}>> G         \\@A{n}AA @AA{-n}A         \\X @>>{u}> G         \end {CD}         \,\,\,\,\,\,       \right ]       \longleftarrow        \left [         \begin {CD}           T_\mathbb {Z} @>{u}>> \operatorname {Lie}(G)         \\@A{n}AA @A{\sim }A{-n}A         \\T_\mathbb {Z} @>>{u}> \operatorname {Lie}(G)         \end {CD}       \right ]       \longrightarrow        \left [         \begin {CD}           T_\mathbb {Z} @>>> 0         \\@A{n}AA @AAA         \\T_\mathbb {Z} @>>{c}> 0         \end {CD}       \right ]     \]
    give isomorphisms

    \begin{eqn}{10.1.6.2}\setnumber{10.1.6.2}\label{hodge-theory-iii-10.1.6.2}\skipsinglepar\[         T_{\mathbb {Z}{/}n\mathbb {Z}}(M)         \simeq  T_\mathbb {Z}{/}nT_\mathbb {Z}.       \tag{10.1.6.2}       \]\end{eqn}\par{}Under this isomorphism, to \(t\in  W_{-1}(T_\mathbb {Z})\) we associate \(\exp (t/n)\in  G_n\).
    By passing to the limit, we deduce \hyperref[hodge-theory-iii-10.1.6.1]{Equation~\ref*{hodge-theory-iii-10.1.6.1}} from \hyperref[hodge-theory-iii-10.1.6.2]{Equation~\ref*{hodge-theory-iii-10.1.6.2}}.\end{construction}\par{}Techniques of Grothendieck allow us to transpose to de Rham cohomology what we have just done in \(\ell \)-adic cohomology.
  Let \(M\) be a 1-motive over an algebraically closed field \(k\).\begin{construction}{10.1.7}\setnumber{10.1.7}\label{hodge-theory-iii-10.1.7}\skipsinglepar\par{}We will construct a vector space \(T_\mathrm {DR}(M)\) over \(k\), endowed with filtrations \(W\) and \(F\).\par{}Let \(M=(X,A,T,G,u)\), and consider \(M\) as a complex concentrated in degrees \(-1\) and \(0\).
    An \emph{extension} of \(M\) by \(\mathbb {G}_\mathrm {a}\) is an extension of \(M\) by the complex consisting of \(\mathbb {G}_\mathrm {a}\) concentrated in degree \(0\).
    We have that \(\operatorname {Hom}(G,\mathbb {G}_\mathrm {a})=\operatorname {Ext}^1(X,\mathbb {G}_\mathrm {a})=0\).
    Thus:

    
 \begin{enumerate}
      
 \item[(a)]
        extensions of \(M\) by \(\mathbb {G}_\mathrm {a}\) have no automorphisms;
      

      
 \item[(b)]
        the sequence
        \[           0           \to  \operatorname {Hom}(X,\mathbb {G}_\mathrm {a})           \to  \operatorname {Ext}^1(M,\mathbb {G}_\mathrm {a})           \to  \operatorname {Ext}^1(G,\mathbb {G}_\mathrm {a})           \to  0         \]
        is exact.
      

    \end{enumerate}


    We have that \(\operatorname {Hom}(T,\mathbb {G}_\mathrm {a})=\operatorname {Ext}^1(T,\mathbb {G}_\mathrm {a})=0\), and so

    
 \begin{enumerate}
      
 \item[(c)]
        \(\operatorname {Ext}^1(A,\mathbb {G}_\mathrm {a})\xrightarrow {\sim }\operatorname {Ext}^1(G,\mathbb {G}_\mathrm {a})\).
      

    \end{enumerate}\par{}The vector space \(\operatorname {Ext}^1(M,\mathbb {G}_\mathrm {a})\) is thus of finite dimension, and, by (a), there thus exists a universal extension of \(M\) by a vector group;
    we denote it by \(M^\natural \colon  X\to  G^\natural \).

    
  \begin {tikzcd}
    & 0
        \ar [r]
        \ar [d]
    & X
        \ar [r,equals]
        \ar [d]
    & X
        \ar [d]
    \\0
        \ar [r]
    & \operatorname {Ext}^1(M,\mathbb {G}_\mathrm {a})^*
        \ar [r]
    & G^\natural 
        \ar [r]
    & G
        \ar [r]
    & 0
    \end {tikzcd}\par{}We set
    \[       \begin {aligned}         T_\mathrm {DR}(M)         &= \operatorname {Lie}(G^\natural )       \\F^0T_\mathrm {DR}(M)         &= \operatorname {Ker}(\operatorname {Lie} G^\natural \to \operatorname {Lie} G)         \simeq  (\operatorname {Ext}^1(M,\mathbb {G}_\mathrm {a}))^*.       \end {aligned}     \]
    Then \(T_\mathrm {DR}\) is functorial in \(M\) (and thus also in \(M^\natural \)), and we define the filtration \(W\) of \(T_\mathrm {DR}(M)\) as coming from the filtration \(W\) of \(M\).\end{construction}\begin{construction}{10.1.8}\setnumber{10.1.8}\label{hodge-theory-iii-10.1.8}\skipsinglepar\par{}Let \(M\) be a 1-motive over \(\mathbb {C}\).
    We have that
    \[       (T_\mathrm {DR}(M),F,W)       \simeq  (T_\mathbb {C}(M),F,W).     \]\begin{proof}\par{}The \(\operatorname {Ext}^i(X,\mathbb {G}_\mathrm {a})\), \(\operatorname {Ext}^i(A,\mathbb {G}_\mathrm {a})\), and \(\operatorname {Ext}^i(T,\mathbb {G}_\mathrm {a})\) (for \(i=0,1\)) are the same in the algebraic and analytic categories.
      Consider the map
      \[         T_\mathbb {C}(M)         = T_\mathbb {Z}(M)\otimes \mathbb {C}         \to  \operatorname {Lie}(G)         \to  G.       \]
      The diagram
      \[         \begin {CD}           X @>>> T_\mathbb {C}(M)/\operatorname {H}_1(G)         \\@| @VVV         \\X @>>> G         \end {CD}       \]
      
      defines an extension of \(M^\mathrm {an}\) by the vector group \(F^0T_\mathbb {C} M\), and thus an extension of \(M\) by this group.
      We have to prove that this is the universal extension.
      For this, it suffices to verify that the category of extensions of \([X\to  T_\mathbb {C}(M)/\operatorname {H}_1(G)]\) by \(\mathbb {G}_\mathrm {a}\) is trivial (i.e. that it has a single object, and no automorphisms).
      This is also the category of extensions of \(T_\mathbb {C}(M)/T_\mathbb {Z}(M)\) by \(\mathbb {G}_\mathrm {a}\), and, indeed,
      \[         \operatorname {Ext}^i(T_\mathbb {C}(M)/T_\mathbb {Z}(M),\mathbb {G}_\mathrm {a})         = 0         \qquad \text {for }i=0,1.       \]\end{proof}\end{construction}\begin{remark}{10.1.9}\setnumber{10.1.9}\label{hodge-theory-iii-10.1.9}\skipsinglepar\par{}The universal extension \(M^\natural \) is characterised by the commutative diagram
    \[       \begin {CD}         X @>>> G^\natural  @>>> G       \\@AAA @AAA @AAA       \\T_\mathbb {Z}(M) @>>> \operatorname {Lie}(G^\natural ) @>>> \operatorname {Lie}(G)       \end {CD}     \]
    whose exterior square is \hyperref[hodge-theory-iii-10.1.3.1]{Equation~\ref*{hodge-theory-iii-10.1.3.1}}, and that induces an isomorphism \(T_\mathbb {C}(M)\xrightarrow {\sim }\operatorname {Lie}(G^\natural )\).\end{remark}\begin{variant}{10.1.10}\setnumber{10.1.10}\label{hodge-theory-iii-10.1.10}\skipsinglepar\par{}A \emph{smooth 1-motive} \(M\) over a scheme \(S\) consists of the following:

    
 \begin{enumerate}
      
 \item[(a)]
        a group scheme \(X\) over \(S\) that is locally (for the étale topology) a constant group scheme defined by a free \(\mathbb {Z}\)-module of finite type;
        an abelian scheme \(A\) over \(S\), and a torus \(T\) over \(S\);
      

      
 \item[(b)]
        an extension \(G\) of \(A\) by \(T\);
      

      
 \item[(c)]
        a morphism \(u\colon  X\to  G\).
      

    \end{enumerate}\par{}For every integer \(n\), a construction analogous to \cref{hodge-theory-iii-10.1.5} associates to a smooth 1-motive \(M\) over \(S\) a group scheme \(T_{\mathbb {Z}{/}n\mathbb {Z}}(M)\) that is finite and flat over \(S\).
    For \(\ell \) invertible over \(S\), the projective system
    \[       T_\ell (M)       = \varprojlim  T_{\mathbb {Z}{/}\ell ^n\mathbb {Z}}(M)     \]
    is a \(\mathbb {Z}_\ell \)-sheaf on \(S\).
    In general, \(T_\ell (M)\) corresponds to an \(\ell \)-divisible group (i.e. a Barsotti–Tate group) over \(S\).\par{}The filtration \(W\) of \(M\), defined as in \cref{hodge-theory-iii-10.1.4}, defines a filtration \(W\) on \(T_\ell (M)\).
    We again have that
    \[       \begin {aligned}         \operatorname {Gr}_0^W(T_\ell (M))         &= X\otimes \widehat {\mathbb {Z}}       \\\operatorname {Gr}_{-1}^W(T_\ell (M))         &= T_\ell (A)       \\\operatorname {Gr}_{-2}^W(T_\ell (M))         &= Y\otimes \widehat {\mathbb {Z}}_\ell (1),       \end {aligned}     \]
    where \(Y=\underline {\operatorname {Hom}}(T,\mathbb {G}_\mathrm {m})^\vee \).\par{}Similarly, a construction analogous to that of \hyperref[hodge-theory-iii-10.1.7]{Construction~\ref*{hodge-theory-iii-10.1.7}} associates to \(M\) a vector bundle \(T_\mathrm {DR}(M)\).\end{variant}\begin{terminology}{10.1.11}\setnumber{10.1.11}\label{hodge-theory-iii-10.1.11}\skipsinglepar\par{}Let \(M\) be a 1-motive.
    We call \(T(M)\), \(T_\ell (M)\), and \(T_\mathrm {DR}(M)\) the \emph{Hodge}, \emph{\(\ell \)-adic}, and \emph{de Rham} \emph{realisations} of \(M\) (respectively).\end{terminology}\section{1-Motives and bi-extensions}\label{hodge-theory-iii-10.2}\par{}The results of this section will not be used elsewhere.
  We will make use of the notion of biextensions [SGA 7, VII, (2.1)] and of the following generalisation.\label{hodge-theory-iii-10.2.1}\par{}Let \(K_i\colon [A_i\to  B_i]\) (for \(i=1,2\)) be complexes of sheaves, concentrated in degrees {0} and \(-1\) (in an arbitrary topos).
    A \emph{biextension} of \(K_1\) and \(K_2\) by an abelian sheaf \(H\) consists of the following:

    
 \begin{enumerate}
      
 \item[(a)]
        a biextension \(P\) of \(B_1\) and \(B_2\) by \(H\), i.e. an \(H\)-torsor \(P\) over \(B_1\times  B_2\), with \(P_{b_1b_2}\) depending biadditively on \(b_1\) and \(b_2\);
      


      
 \item[(b)]
        a trivialisation (i.e. a biadditive section) of the biextension of \(B_1\) and \(A_2\) by \(H\), given by the inverse image of \(P\);
      


      
 \item[(c)]
        a trivialisation of the biextension of \(A_1\) and \(B_2\) by \(H\), given by the inverse image of \(P\);
      

    \end{enumerate}


    where we demand that the trivialisations in (b) and (c) agree on \(A_1\times  A_2\).\par{}We denote by \(\operatorname {\textbf {Biext}}(K_1,K_2;H)\) the category of biextensions of \(K_1\) and \(K_2\) by \(H\).
    This is a Picard category (the objects add).
    The group \(\operatorname {Biext}^0(K_1,K_2;H)\) of automorphisms of an arbitrary biextension of \(K_1\) and \(K_2\) by \(H\) is
    \[       \operatorname {Biext}^0(K_1,K_2;H)       = \operatorname {Hom}(\operatorname {H}^0(K_1)\otimes \operatorname {H}^0(K_2),H).     \]
    We denote by \(\operatorname {Biext}^1(K_1,K_2;H)\) the group of isomorphism classes of biextensions.
    We can show, as in [SGA 7, VII, (3.6.5)], that
    \[       \operatorname {Biext}^i(K_1,K_2;H)       = \operatorname {Ext}^i(K_1\overset {\mathrm {L}}{\otimes } K_2;H)       \qquad \text {for }i=0,1.     \]
    From this formula, or by elementary means, we can show that, if \(K'_i\to  K_i\) is a quasi-isomorphism (for \(i=1,2\)), then giving a biextension of \(K_1\) and \(K_2\) by \(H\) is equivalent to giving a biextension of \(K'_1\) and \(K'_2\) by \(H\).\par{}This also applies to replacing "sheaves" by "algebraic groups" (resp. by "complex analytic groups"): we can interpret these as sheaves on a suitable site.\label{hodge-theory-iii-10.2.2}\par{}In this section, we identity a 1-motive \(M=(X,T,A,G,u)\) with a complex concentrated in \emph{degrees \(0\) and \(-1\)}
    \[       M\colon  [X\to  G].     \]
    Over \(\mathbb {C}\), we denote by \(M^\mathrm {an}\) the complex of complex analytic groups \([X\to  G^\mathrm {an}]\).
    We have the following rigidity result:\begin{lemma}{10.2.2.1}\setnumber{10.2.2.1}\label{hodge-theory-iii-10.2.2.1}\skipsinglepar\par{}Let \(M_i=(X_i,A_i,T_i,G_i,u_i)\) (for \(i=1,2\)) be 1-motives.
      Then
      \[         \operatorname {Biext}^0(M_1,M_2;\mathbb {G}_\mathrm {m})         = 0.       \]
      Over \(\mathbb {C}\), we similarly have that \(\operatorname {Biext}^0(M_1^\mathrm {an},M_2^\mathrm {an};\mathbb {G}_\mathrm {m})=0\).\begin{proof}\par{}Every biadditive morphism \(G_1\times  G_2\to \mathbb {G}_\mathrm {m}\) is trivial.\end{proof}\end{lemma}\begin{construction}{10.2.3}\setnumber{10.2.3}\label{hodge-theory-iii-10.2.3}\skipsinglepar\par{}Let \(M_i=(X_i,A_i,T_i,G_i,u_i)\) (for \(i=1,2\)) be 1-motives over \(\mathbb {C}\).
    We have that
    \[       \operatorname {Biext}^1(M_1,M_2;\mathbb {G}_\mathrm {m})       = \operatorname {Hom}(T(M_1)\otimes  T(M_2),\mathbb {Z}(1)).     \]\begin{proof}\par{}We will first calculate \(\operatorname {Biext}^1(M_1^\mathrm {an},M_2^\mathrm {an};\mathbb {G}_\mathrm {m})\).\par{}If \(V\) and \(W\) are vector spaces, then, \emph{in the analytic category},

      \begin{eqn}{10.2.3.1(an)}\setnumber{10.2.3.1(an)}\label{hodge-theory-iii-10.2.3.1-an}\skipsinglepar\[           \operatorname {Biext}^0(V,W;\mathbb {G}_\mathrm {m})           \underset {\exp }{\xleftarrow {\sim }} \operatorname {Hom}(V\otimes  W,\mathbb {C})         \tag{10.2.3.1(\text{an})}         \]\end{eqn}

      \begin{eqn}{10.2.3.2(an)}\setnumber{10.2.3.2(an)}\label{hodge-theory-iii-10.2.3.2-an}\skipsinglepar\[           \operatorname {Biext}^1(V,W;\mathbb {G}_\mathrm {m})           = 0.         \tag{10.2.3.2(\text{an})}         \]\end{eqn}

      From the fact that an extension of \(W\) by \(\mathbb {G}_\mathrm {m}\), over an arbitrary base, is always locally trivial, it follows, essentially formally, that
      \[         \begin {aligned}           \operatorname {Biext}^i(V,W;\mathbb {G}_\mathrm {m})           &= \operatorname {Ext}^i(V,\underline {\operatorname {Hom}}(W,\mathbb {G}_\mathrm {m}))         \\&= \operatorname {Ext}^i(V,W^*)           \qquad \text {for }i=1,2.         \end {aligned}       \]\par{}Let \(V_\mathbb {Z}\) and \(W_\mathbb {Z}\) be free \(\mathbb {Z}\)-modules of finite rank.
      Since
      \[         [V_\mathbb {Z}\to  V_\mathbb {C}]         \to  [0\to  V_\mathbb {C}{/}V_\mathbb {Z}]       \]
      is a quasi-isomorphism, the "inverse image" functor is an equivalence.

      \begin{eqn}{10.2.3.3(an)}\setnumber{10.2.3.3(an)}\label{hodge-theory-iii-10.2.3.3-an}\skipsinglepar\[           \operatorname {\textbf {Biext}}(V_\mathbb {C}{/}V_\mathbb {Z},W_\mathbb {C}{/}W_\mathbb {Z};\mathbb {G}_\mathrm {m})           \simeq  \operatorname {\textbf {Biext}}([V_\mathbb {Z}\to  V_\mathbb {C}],[W_\mathbb {Z}\to  W_\mathbb {C}];\mathbb {G}_\mathrm {m}).         \tag{10.2.3.3(\text{an})}         \]\end{eqn}

      Every biadditive morphism from \(V_\mathbb {C}{/}V_\mathbb {Z}\times  W_\mathbb {C}{/}W_\mathbb {Z}\) to \(\mathbb {G}_\mathrm {m}\) is trivial, whence (a) in the following:\begin{lemma}{10.2.3.4}\setnumber{10.2.3.4}\label{hodge-theory-iii-10.2.3.4}\skipsinglepar\par{}\begin{enumerate}
          
 \item[(a)]\(\operatorname {Biext}^0(V_\mathbb {C}{/}V_\mathbb {Z},W_\mathbb {C}{/}W_\mathbb {Z};\mathbb {G}_\mathrm {m})=0\)

          
 \item[(b)]\(\operatorname {Biext}^1(V_\mathbb {C}{/}V_\mathbb {Z},W_\mathbb {C}{/}W_\mathbb {Z};\mathbb {G}_\mathrm {m})=\operatorname {Hom}(V_\mathbb {Z}\otimes  W_\mathbb {Z},\mathbb {Z}(1))\).

        \end{enumerate}\begin{proof}\par{}Let \(\psi _i\colon  V_\mathbb {C}\otimes  W_\mathbb {C}\to \mathbb {C}\) be bilinear maps, with \(\psi =\psi _1-\psi _2\) taking values in \(\mathbb {Z}(1)=2\pi  i\mathbb {Z}\) on \(V_\mathbb {Z}\otimes  W_\mathbb {Z}\).
          Let \(P(\psi _1,\psi _2)\) be the following biextension of \([V_\mathbb {Z}\to  V_\mathbb {C}]\) and \([W_\mathbb {Z}\to  W_\mathbb {C}]\) by \(\mathbb {G}_\mathrm {m}\): the trivial biextension of \(V_\mathbb {C}\) and \(W_\mathbb {C}\) by {\(\mathbb {G}_\mathrm {m}\)}, endowed with the trivialisations \(\varphi _i=\exp \psi _i\colon  V_\mathbb {Z}\otimes  W_\mathbb {C}\to \mathbb {G}_\mathrm {m}\) for \(i=1,2\).
          By \hyperref[hodge-theory-iii-10.2.3.2-an]{Equation~\ref*{hodge-theory-iii-10.2.3.2-an}}, every biextension is of this form;
          by \hyperref[hodge-theory-iii-10.2.3.1-an]{Equation~\ref*{hodge-theory-iii-10.2.3.1-an}}, \(P(\psi _1,\psi _2)\simeq  P(\psi '_1,\psi '_2)\) if and only if \(\psi '_1-\psi '_2=\psi _1-\psi _2\).
          The \(P(\psi ,0)\) thus represent each biextension exactly once.\end{proof}\end{lemma}\label{hodge-theory-iii-10.2.3.5}\par{}Under the general hypotheses of \cref{hodge-theory-iii-10.2.1}, let \(K_i=[A_i+F_i\to  B_i]\) be complexes (for \(i=1,2\)).
        Let \(K'_i\) be the subcomplex \([A_i\to  B_i]\).
        We can show that giving a biextension of \(K_1\) and \(K_2\) by \(H\) is equivalent to giving (both of) the following:

        
 \begin{enumerate}
          
 \item[(a)]
            a biextension \(P\) of \(K'_1\) and \(K'_2\) by \(H\); and
          

          
 \item[(b)]
            trivialisations of the biextensions of \([0\to  F_1]\) and \(K'_2\) (resp. \(K'_1\) and \([0\to  F_2]\)) by \(H\), given by the inverse images of \(P\), that agree on \(F_1\times  F_2\).
          

        \end{enumerate}\par{}In the particular case where
        \[           \begin {aligned}             \operatorname {Biext}^i(F_1,K'_2;H)             &= 0           \\\operatorname {Biext}^i(K'_1,F_2;H)             &= 0           \end {aligned}         \]
        
        for \(i=1,2\), there exists exactly one trivialisation \(\varphi _1\), a section of \(P\) over \(F_1\times  B_2\) (resp. \(\varphi _2\), a section of \(P\) over \(B_1\times  F_2\)) as in (b).
        The category \(\operatorname {\textbf {Biext}}(K_1,K_2;H)\) can then be identified with the subcategory of \(\operatorname {\textbf {Biext}}(K'_1,K'_2;H)\) consisting of the \(P\) for which \(\varphi _1=\varphi _2\) on \(F_1\times  F_2\).\label{hodge-theory-iii-10.2.3.6}\par{}Let \(V_\mathbb {Z}\) and \(W_\mathbb {Z}\) be free \(\mathbb {Z}\)-modules of finite type, and \(F\) and \(G\) vector subspaces of \(V_\mathbb {C}\) and \(W_\mathbb {C}\) (respectively).
        We have that
        \[           \operatorname {Biext}^i(F,[W_\mathbb {Z}\to  W_\mathbb {C}];\mathbb {G}_\mathrm {m})           = 0.         \]
        For \(i=0\), this follows from \hyperref[hodge-theory-iii-10.2.3.1-an]{Equation~\ref*{hodge-theory-iii-10.2.3.1-an}} and from the fact that a bilinear map
        \[           B\colon  F\times  W_\mathbb {C}\to \mathbb {C}           \quad \text {such that}\quad            B(F,W_\mathbb {Z})\subset \mathbb {Z}(1)         \]
        is zero.
        For \(i=1\), this implies that every biadditive morphism \(F\times  W_\mathbb {Z}\to \mathbb {G}_\mathrm {m}\) comes from a biadditive morphism \(F\times  W_\mathbb {C}\to \mathbb {G}_\mathrm {m}\).
        In other words, we use the long exact sequence
        \[           \begin {aligned}             0             &\to  \operatorname {Biext}^0(F,[W_\mathbb {Z}\to  W_\mathbb {C}];\mathbb {G}_\mathrm {m})             \to  \operatorname {Biext}^0(F,W_\mathbb {C};\mathbb {G}_\mathrm {m})             \xrightarrow {\sim } \operatorname {Biext}^0(F,W_\mathbb {Z};\mathbb {G}_\mathrm {m})           \\&\to  \operatorname {Biext}^1(F,[W_\mathbb {Z}\to  W_\mathbb {C}];\mathbb {G}_\mathrm {m})             \to  \operatorname {Biext}^1(F,W_\mathbb {C};\mathbb {G}_\mathrm {m}).           \end {aligned}         \]\par{}We can thus apply \cref{hodge-theory-iii-10.2.3.5} to the complexes \(K_1=[V_\mathbb {Z}\oplus  F\to  V_\mathbb {C}]\) and \(K_2=[W_2\oplus  G\to  W_\mathbb {C}]\).
        We find that the biextensions of \(K_1\) and \(K_2\) by \(\mathbb {G}_\mathrm {m}\) can be identified with certain biextensions of \(V_\mathbb {C}{/}V_\mathbb {Z}\) and \(W_\mathbb {C}{/}W_\mathbb {Z}\) (respectively) by \(\mathbb {G}_\mathrm {m}\), i.e. with certain forms \(\psi \colon  V_\mathbb {Z}\otimes  W_\mathbb {Z}\to  \mathbb {Z}(1)\) (\hyperref[hodge-theory-iii-10.2.3.4]{Lemma~\ref*{hodge-theory-iii-10.2.3.4}}).
        For \(P(\psi _1,\psi _2)\) (\hyperref[hodge-theory-iii-10.2.3.4]{Lemma~\ref*{hodge-theory-iii-10.2.3.4}}) to come from a biextension of \(K_1\) and \(K_2\) by \(\mathbb {G}_\mathrm {m}\), it is necessary and sufficient that \(\exp \psi _2\colon  F\times  W_\mathbb {C}\to \mathbb {G}_\mathrm {m}\) and \(\exp \psi _1\colon  V_\mathbb {C}\times  G\to \mathbb {G}_\mathrm {m}\) agree on \(F\times  G\), i.e. that \(\psi (F,G)=0\).\begin{lemma}{10.2.3.7}\setnumber{10.2.3.7}\label{hodge-theory-iii-10.2.3.7}\skipsinglepar\par{}With the above notation, we have that \(\operatorname {Biext}^0(K_1,K_2;\mathbb {G}_\mathrm {m})=0\), and \(\operatorname {Biext}^1(K_1,K_2;\mathbb {G}_\mathrm {m})\) can be identified with the set of forms \(\psi \colon  V_\mathbb {Z}\otimes  W_\mathbb {Z}\to \mathbb {Z}(1)\) with complexifications satisfying \(\Psi (F,G)=0\).\begin{proof}\par{}For \(\psi \) as in \hyperref[hodge-theory-iii-10.2.3.7]{Lemma~\ref*{hodge-theory-iii-10.2.3.7}}, with \(\psi =\psi _1-\psi _2\), the biextension corresponding to \(\psi \) is given by the following trivialisations of the trivial biextension of \(V_\mathbb {C}\) and \(W_\mathbb {C}\) by \(\mathbb {G}_\mathrm {m}\):
          \begin{itemize}\item{}on \((V_\mathbb {Z}\oplus  F)\times  W_\mathbb {C}\) : \(\exp \psi _1(v_\mathbb {Z},w_\mathbb {C})\cdot \exp \psi _2(f,w_\mathbb {C})\)
            \item{}on \(V_\mathbb {C}\times (W_\mathbb {C}\oplus  G)\) : \(\exp \psi _2(v_\mathbb {C},w_\mathbb {Z})\cdot \exp \psi _2(v_\mathbb {C},g)\).\end{itemize}\end{proof}\end{lemma}\label{hodge-theory-iii-10.2.3.8}\par{}For the 1-motives \(M_i\), we have quasi-isomorphisms
        \[           \begin {CD}             T_\mathbb {Z}(M_i)\oplus  F^0T_\mathbb {C}(M_i) @>>> T_\mathbb {Z}(M_i) @>>> X_i           \\@VVV @VVV @VVV           \\T_\mathbb {C}(M_i) @>>> \operatorname {Lie}(G_i) @>>> G_i           \end {CD}         \]
        
        The group \(\operatorname {Biext}^1(M_1^\mathrm {an},M_2^\mathrm {an};\mathbb {G}_\mathrm {m})\) can thus be identified with the group of maps
        \[           \psi \colon  T_\mathbb {Z}(M_1)\otimes  T_\mathbb {Z}(M_2)           \to  \mathbb {Z}(1)         \]
        whose complexification is compatible with the Hodge filtration.\begin{lemma}{10.2.3.9}\setnumber{10.2.3.9}\label{hodge-theory-iii-10.2.3.9}\skipsinglepar\par{}The map \(\operatorname {Biext}^1(M_1,M_2;\mathbb {G}_\mathrm {m})\to \operatorname {Biext}^1(M_1^\mathrm {an},M_2^\mathrm {an};\mathbb {G}_\mathrm {m})\) is injective.
        Its image consists of those biextensions \(P\) whose restriction to \(G_1\times  G_2\) is the inverse image of a biextension of \(A_1\) and \(A_2\) by \(\mathbb {G}_\mathrm {m}\).\begin{proof}\par{}We have that
          \[             \begin {aligned}               \operatorname {Biext}^1(A_1,A_2;\mathbb {G}_\mathrm {m})               &\xrightarrow {\sim } \operatorname {Biext}^1(G_1,G_2;\mathbb {G}_\mathrm {m})             \\\operatorname {Biext}^1(A_1,A_2;\mathbb {G}_\mathrm {m})               &\xrightarrow {\sim } \operatorname {Biext}^1(A_1^\mathrm {an},A_2^\mathrm {an};\mathbb {G}_\mathrm {m})             \end {aligned}           \]
          where the former follows from [SGA 7, VIII, (3.5)] and the latter from GAGA.\par{}Let \(P\) be a biextension of \(M_1\) and \(M_2\) by \(\mathbb {G}_\mathrm {m}\).
          If \(P^\mathrm {an}\) is trivial, then:

          
 \begin{enumerate}
            
 \item[(a)]
              The restriction of \(P\) to \(G_1\times  G_2\) is the image of a biextension \(P_0\) of \(A_1\) and \(A_2\) by \(\mathbb {G}_\mathrm {m}\);
              the latter is trivial.
              If this were not the case, then it would be analytically non-trivial, and would thus give rise to a non-zero bilinear form \(\psi _0\), and the form \(\psi \) corresponding to \(P\) would be \emph{a fortiori} non-zero.
              The restriction of \(P\) to \(G_1\times  G_2\) is thus trivial (and uniquely trivialisable, even analytically).
            


            
 \item[(b)]
              \(P\) is thus defined by biadditive maps \(X_1\times  G_2\to \mathbb {G}_\mathrm {m}\) and \(G_1\times  X_2\to \mathbb {G}_\mathrm {m}\).
              These are zero in the analytic category, and thus thus zero, and so \(P\) is trivial.
            

          \end{enumerate}\par{}Let \(P\) be a biextension of \(M_1^\mathrm {an}\) and \(M_2^\mathrm {an}\) by \(\mathbb {G}_\mathrm {m}\) that satisfies the condition \hyperref[hodge-theory-iii-10.2.3.9]{Lemma~\ref*{hodge-theory-iii-10.2.3.9}}.
          It remains to prove that \(P\) is algebrisable.
          By hypothesis, and GAGA, its restriction \(P_1\) to \(G_1\times  G_2\) is algebrisable.
          We can conclude by showing that a trivialisation of \(P_1\) on \(X_1\times  G_2\) (resp. \(G_1\times  X_2\)) is automatically algebraic: \(P_1\) can be thought of as an extension of \(X_1\otimes  G_2\) by \(\mathbb {G}_\mathrm {m}\), which allows us to reduce to showing that
          \[             \operatorname {Ext}^i(G_1,\mathbb {G}_\mathrm {m})             \xrightarrow {\sim } \operatorname {Ext}^i(G_1^\mathrm {an},\mathbb {G}_\mathrm {m})             \qquad \text {for }i=1,2.           \]
          This follows from the same claim for an abelian variety (GAGA) and for \(\mathbb {G}_\mathrm {m}\).\par{}By \hyperref[hodge-theory-iii-10.2.3.9]{Lemma~\ref*{hodge-theory-iii-10.2.3.9}}, \(\operatorname {Biext}^1(M_1,M_2;\mathbb {G}_\mathrm {m})\) can be identified with the set of
          \[             \psi \colon              T_\mathbb {Z}(M_1)\otimes  T_\mathbb {Z}(M_2)             \to  \mathbb {Z}(1)           \]
          such that

          
 \begin{enumerate}
            
 \item[(a)]
              \(\psi \) is compatible with the Hodge filtration (\cref{hodge-theory-iii-10.2.3.8}); and
            

            
 \item[(b)]
              the restriction of \(\psi \) to \(W_{-1}T_\mathbb {Z}(M_1)\otimes  W_{-1}T_\mathbb {Z}(M_2)\) comes from a form on
              \[                 \operatorname {Gr}_{-1}^WT_\mathbb {Z}(M_1)\otimes \operatorname {Gr}_{-1}^WT_\mathbb {Z}(M_2).               \]
            

          \end{enumerate}\par{}Condition (b) implies that \(\psi \) is compatible with \(W\), and this concludes \hyperref[hodge-theory-iii-10.2.3]{Construction~\ref*{hodge-theory-iii-10.2.3}}.\end{proof}\end{lemma}\end{proof}\end{construction}\begin{remark}{10.2.4}\setnumber{10.2.4}\label{hodge-theory-iii-10.2.4}\skipsinglepar\par{}A biextension \(P\) of \(M_1\) and \(M_2\) by \(\mathbb {G}_\mathrm {m}\) defines, in an evident way, a biextension \(P^0\) of \(M_2\) and \(M_1\) by \(\mathbb {G}_\mathrm {m}\).
    
    If \(\psi \) (resp. \(\psi ^0\)) is the bilinear form defined by \(P\) (resp. by \(P^0\)), then \(\psi (x,y)=-\psi ^0(y,x)\).
    This is clear from the equations in \hyperref[hodge-theory-iii-10.2.3.7]{Lemma~\ref*{hodge-theory-iii-10.2.3.7}}.\par{}\hyperref[hodge-theory-iii-10.2.3]{Construction~\ref*{hodge-theory-iii-10.2.3}} has purely algebraic analogues in \(\ell \)-adic or de Rham cohomology.\end{remark}\begin{construction}{10.2.5}\setnumber{10.2.5}\label{hodge-theory-iii-10.2.5}\skipsinglepar\par{}Let \(M_i=(X_i,A_i,T_i,G_i,u_i)\) (for \(i=1,2\)) be 1-motives over an algebraically closed field \(k\) of characteristic 0.
    Let \(P\) be a biextension of \(M_1\) and \(M_2\) by \(\mathbb {G}_\mathrm {m}\).
    We associate to \(P\) a morphism
    \[       T_{\mathbb {Z}{/}n\mathbb {Z}}(M_1)\otimes  T_{\mathbb {Z}{/}n\mathbb {Z}}(M_2)       \to  \mathbb {Z}{/}n\mathbb {Z}(1).     \]\begin{proof}\par{}Let \(m_i\in  T_{\mathbb {Z}{/}n\mathbb {Z}}(M_1)\) be represented by \((x_i,g_i)\), with \(u_i(x_i)=ng_i\).
      We will construct isomorphisms \(a_1\) and \(a_2\) from \(P_{g_1,g_2}^{\otimes  n}\) to the trivial torsor \(\mathbb {G}_\mathrm {m}\):
      \[         \begin {aligned}           a_1\colon  P_{g_1,g_2}^{\otimes  n}           \xrightarrow {\sim } P_{ng_1,g_2}           &= P_{u(x_1),g_2}           \xrightarrow {\sim } \mathbb {G}_\mathrm {m}         \\a_2\colon  P_{g_1,g_2}^{\otimes  n}           \xrightarrow {\sim } P_{g_1,ng_2}           &= P_{g_1,u(x_2)}           \xrightarrow {\sim } \mathbb {G}_\mathrm {m}.         \end {aligned}       \]
      We set \(a_2=\Phi (m_1,m_2)a_1\).
      We can verify that \(\Phi (m_1,m_2)\) does not depend on the choice of the \((x_i,g_i)\), and is an \(n\)-th root of unity.
      This is the desired form.\end{proof}\end{construction}\begin{proposition}{10.2.6}\setnumber{10.2.6}\label{hodge-theory-iii-10.2.6}\skipsinglepar\par{}On \(\mathbb {C}\), the form in \hyperref[hodge-theory-iii-10.2.5]{Construction~\ref*{hodge-theory-iii-10.2.5}} is induced by the form in \hyperref[hodge-theory-iii-10.2.3]{Construction~\ref*{hodge-theory-iii-10.2.3}} by reduction modulo \(n\).\begin{proof}\par{}This can be proven by using the formulas given in \hyperref[hodge-theory-iii-10.2.3.7]{Lemma~\ref*{hodge-theory-iii-10.2.3.7}}.\end{proof}\end{proposition}\begin{construction}{10.2.7}\setnumber{10.2.7}\label{hodge-theory-iii-10.2.7}\skipsinglepar\par{}Let \(M_i=(X_i,A_i,T_i,G_i,u_i)\) (for \(i=1,2\)) be 1-motives over an algebraically closed field \(k\).
    Let \(P\) be a biextension of \(M_1\) and \(M_2\) by \(\mathbb {G}_\mathrm {m}\).
    We associate to \(P\) a morphism
    \[       T_\mathrm {DR}(M_1)\otimes _k T_\mathrm {DR}(M_2)       \to  k.     \]\begin{proof}\(\natural \)\(\natural \)\label{hodge-theory-iii-10.2.7.1}\par{}A \(\natural \)-extension of a smooth commutative group \(G\) by a smooth commutative group \(H\) consists of the following:

        
 \begin{enumerate}
          
 \item[(a)]
            an extension \(E\) of \(G\) by \(H\);
            if \(\mu \colon  G\times  G\to  G\) is the addition, then we consider \(E\) as an \(H\)-torsor over \(G\) endowed with a morphism
            \[               \nu \colon  \operatorname {pr}_1^*E + \operatorname {pr}_2^*E               \to  \mu ^*E             \]
            of \(H\)-torsors over \(G\times  G\) (addition);
          

          
 \item[(b)]
            a connection on the \(H\)-torsor \(E\), such that the map \(\nu \) is horizontal.
            This connection is automatically flat.
          

        \end{enumerate}\par{}If \(G\) is the extension of an abelian variety by a torus, then every extension of \(G\) by \(\mathbb {G}_\mathrm {m}\) admits a \(\natural \)-structure;
        two \(\natural \)-structures differ by an invariant form on \(G\) (with values in the Lie algebra of \(\mathbb {G}_\mathrm {m}\)).
        We define, in an evident way, the Baer sum of two \(\natural \)-extensions.\label{hodge-theory-iii-10.2.7.2}\par{}Let \(P\) be a biextension of \(G_1\) and \(G_2\) by \(H\).
        Consider \(P\) as an \(H\)-torsor \(P\) over \(G_1\times  G_2\), endowed with morphisms
        \[           \begin {aligned}             \nu _1\colon              \operatorname {pr}_{13}^*P + \operatorname {pr}_{23}^*P             &\to  (\mu _1\times \mathrm {Id})^*P             \quad \text {on}             \quad  G_1\times  G_1\times  G_2           \\\nu _2\colon              \operatorname {pr}_{12}^*P + \operatorname {pr}_{13}^*P             &\to  (\mathrm {Id}\times \mu _2)^*P             \quad \text {on}             \quad  G_1\times  G_2\times  G_2.           \end {aligned}         \]
        A \emph{\(\natural \)-1-structure} on \(P\) is a connection on the \(H\)-torsor \(P\) with respect to \(G_1\times  G_2\to  G_2\) (the covariant derivative is defined only for vector fields that are parallel to \(G_1\)) such that \(\nu _1\) and \(\nu _2\) are horizontal.
        We similarly define \emph{\(\natural \)-2-structures}.
        A \emph{\(\natural \)-structure} is the data of a \(\natural \)-1-structure and of a \(\natural \)-2-structure.
        This is also a connection on the \(H\) torsor \(P\) such that \(\nu _1\) and \(\nu _2\) are horizontal.\par{}The \emph{curvature} \(R\) of a \(\natural \)-biextension \(P\) is the curvature of the connection of \(P\).
        This is a 2-form on \(G_1\times  G_2\) that is invariant under translation, with values in \(\operatorname {Lie}(H)\), i.e. it is an alternating form on \(\operatorname {Lie}(G_1)\times \operatorname {Lie}(G_2)\), with values in \(\operatorname {Lie}(H)\).
        Its restriction to \(\operatorname {Lie}(G_1)\) and \(\operatorname {Lie}(G_2)\) is zero, and so \(R\) defines a pairing

        \begin{eqn}{10.2.7.3}\setnumber{10.2.7.3}\label{hodge-theory-iii-10.2.7.3}\skipsinglepar\[             \Phi \colon  \operatorname {Lie}(G_1)\otimes \operatorname {Lie}(G_2)             \to  \operatorname {Lie}(H)           \tag{10.2.7.3}           \]\end{eqn}

        (also called the \emph{curvature} of \(P\)), with
        \[           R(g_1+g_2, g'_1+g'_2)           = \Phi (g_1,g'_2) - \Phi (g'_1,g_2).         \]\par{}For complexes of smooth groups \(K_i\colon [A_i\to  B_i]\), a \emph{\(\natural \)-biextension} (resp. a \emph{\(\natural \)-\(i\)-extension}) of \(K_1\) and \(K_2\) by \(H\) is a \(\natural \)-biextension (resp. \(\natural \)-\(i\)-biextension) of \(B_1\) and \(B_2\) by \(H\), endowed with trivialisations (as \(\natural \)-biextensions, resp. \(\natural \)-\(i\)-biextensions) as in \cref{hodge-theory-iii-10.2.1}.\begin{proposition}{10.2.7.4}\setnumber{10.2.7.4}\label{hodge-theory-iii-10.2.7.4}\skipsinglepar\par{}Let \(P^\natural \) be the biextension of \(M_1^\natural \) and \(M_2^\natural \) by \(\mathbb {G}_\mathrm {m}\) given by the inverse image of \(P\).
        Then there exists exactly one \(\natural \)-structure on \(P^\natural \).\begin{proof}\par{}We first prove uniqueness: we have to show that, on the trivial biextension of \(M_1^\natural \) and \(M_2^\natural \) by \(\mathbb {G}_\mathrm {m}\), every \(\natural \)-structure is trivial.
          A \(\natural \)-structure is, first of all, a connection on the trivial \(\mathbb {G}_\mathrm {m}\)-torsor over \(G_1^\natural \times  G_2^\natural \), i.e. a field of forms on \(G_1^\natural \times  G_2^\natural \), say \(\Gamma _{g_1,g_2}(t_1+t_2)\).
          The axiom of \(\natural \)-structures implies that
          \[             \Gamma _{g_1,g_2}(t_1+t_2)             = \Phi _1(g_1,t_2) + \Phi _2(t_1,g_2)           \]
          where \(\Phi _1\) and \(\Phi _2\) are biadditive.
          Now, \(G_i^\natural \) is an extension of \(X_i\otimes  k\) by the universal extension \(G'_i\) of \(G_i\) by an additive group.
          Every morphism \(G'_i\to \mathbb {G}_\mathrm {a}\) is trivial.
          Furthermore, \(\Phi _1\) factors through \((X_1\otimes  k)\otimes \operatorname {Lie}(G_2^\natural )\).
          Finally, the restriction of \(\Gamma \) to \(X_1\times  G_2^\natural \) must be trivial, and thus so too must be the restriction of \(\Phi _1\) to \(X_1\times \operatorname {Lie}(G_2^\natural )\).
          We thus have that \(\Phi _1=0\).
          Similarly, \(\Phi _2=0\), and so \(\Gamma =0\).\par{}To prove existence, it suffices to construct a \(\natural \)-2-structure on the biextension of \(M_1^\natural \) and \(M_2\) by \(\mathbb {G}_\mathrm {m}\) given by the inverse image of \(P\): by the inverse image, we will obtain a \(\natural \)-2-structure on \(P^\natural \) and, symmetrically, a \(\natural \)-1-structure.\par{}For \(g\in  G_1\), \(P_g\) is an extension of \(G_2\) by \(\mathbb {G}_\mathrm {m}\).
          Let \(C_g\) be the set of \(\natural \)-structures on \(P_g\).
          This is a torsor under \(\operatorname {Lie}(G_2)^*\).
          For \(g\in  G_1\), the \(C_g\) are the fibres of a torsor \(C\) over \(G_1\);
          the Baer sum of \(\natural \)-extensions further makes \(C\) an extension of \(G_1\) by \(\operatorname {Lie}(G_2)^*\)
          
          We lift \(u_1\colon  X\to  G_1\) to \(u'_1\colon  X\to  C\) by sending \(x\in  X\) to the trivial connection on the trivial extension \(P_{u_1(x)}\).
          We thus obtain a \(\natural \)-2-structure on the biextension of \([X\to  C]\) and \(M_2\) by \(\mathbb {G}_\mathrm {m}\) given by the inverse image of \(P\).\par{}By the universal property of \(M_1^\natural \), there exists a unique commutative diagram

          
  \begin {tikzcd}
          & {[ X}
            \ar [r]
          & {G]}
          \\M_1^\natural 
            \ar [ur]
            \ar [rrr,swap,"v"]
          &&&{[ X}
            \ar [ul]
            \ar [r]
          & {C]}
          \end {tikzcd}


          by taking an inverse image under \(v\), we find the desired \(\natural \)-2-structure.\par{}By definition, the pairing in \hyperref[hodge-theory-iii-10.2.7]{Construction~\ref*{hodge-theory-iii-10.2.7}} is the negative of the curvature of the \(\natural \)-biextension \hyperref[hodge-theory-iii-10.2.7.4]{Proposition~\ref*{hodge-theory-iii-10.2.7.4}}.\end{proof}\end{proposition}\end{proof}\end{construction}\begin{proposition}{10.2.8}\setnumber{10.2.8}\label{hodge-theory-iii-10.2.8}\skipsinglepar\par{}Over \(\mathbb {C}\), the pairing in \hyperref[hodge-theory-iii-10.2.7]{Construction~\ref*{hodge-theory-iii-10.2.7}} is the complexification of the pairing in \hyperref[hodge-theory-iii-10.2.3]{Construction~\ref*{hodge-theory-iii-10.2.3}}.\begin{proof}\(\natural \)\(P^\natural \)\(\natural \)\([T_\mathbb {Z}(M_1)\to  T_\mathbb {C}(M_1)]\)\([T_\mathbb {Z}(M_2)\to  T_\mathbb {C}(M_2)]\)\(\mathbb {G}_\mathrm {m}\)\end{proof}\end{proposition}\begin{lemma}{10.2.9}\setnumber{10.2.9}\label{hodge-theory-iii-10.2.9}\skipsinglepar\par{}Let \(V_\mathbb {Z}\) and \(W_\mathbb {Z}\) be two free \(\mathbb {Z}\)-modules of finite rank, and \(P\) a biextension of \([V_\mathbb {Z}\to  V_\mathbb {C}]\) and \([W_\mathbb {Z}\to  W_\mathbb {C}]\) by \(\mathbb {G}_\mathrm {m}\) with underlying biextension \(P\).
    Then the curvature \(\psi \colon  V_\mathrm {c}\otimes  W_\mathbb {C}\to \mathbb {C}\) of \(P^\natural \) agrees with the negative of the complexification of the form corresponding to \(P\) under \hyperref[hodge-theory-iii-10.2.3.4]{Lemma~\ref*{hodge-theory-iii-10.2.3.4}}.\begin{proof}\(P=P(\psi _1,\psi _2)\)\(\natural \)\(P\)\(\mathbb {G}_\mathrm {m}\)\(V_\mathbb {C}\times  W_\mathbb {C}\)\(\Gamma _{v,w}(v'+w')\)\(V_\mathbb {C}\times  W_\mathbb {C}\)\(\natural \)\[       \Gamma _{v,w}(v'+w')       = \Phi _1(v,w') + \Phi _2(v',w)     \]\(\Phi _1\)\(\Phi _2\)\(\exp \psi _1(v_\mathbb {Z},w_\mathbb {C})\)\(\exp \psi _2(v_\mathbb {C},w_\mathbb {Z})\)\(\Phi _i=-\psi _i\)\(\mathrm {d}\Gamma \)\[       R_{v,w}(v'+w',v''+w'')       = \Phi _1(v',w'') + \Phi _2(v'',w') - \Phi _1(v'',w') - \Phi _2(v',w'').     \]\(P\)\[       \Phi _1(v',w'') - \Phi _2(v',w'')       = -(\psi _1(v',w'') - \psi _2(v',w''))       = -\psi (v',w'').     \qedhere\]\end{proof}\end{lemma}\label{hodge-theory-iii-10.2.10}\par{}Let \(M\) be a 1-motive over \(\mathbb {C}\).
    Then \(M\) has a corresponding mixed Hodge structure \(T(M)\) of type \(\{(0,0),(0,-1),(-1,0),(-1,-1)\}\).
    The mixed Hodge structure \(\operatorname {Hom}(T(M),\mathbb {Z}(1))\) is again of this type, and thus defines a 1-motive \(M^*\).
    The evident map \(T(M)\otimes  T(M^*)\to \mathbb {Z}(1)\) defines a biextension \(P\) of \(M\) and \(M^*\) by \(\mathbb {G}_\mathrm {m}\).
    The 1-motive \(M^*\), endowed with the biextension \(P\), is the \emph{Cartier dual} of \(M\).
    We call \(P\) the \emph{Poincaré biextension}.\label{hodge-theory-iii-10.2.11}\par{}The construction in \cref{hodge-theory-iii-10.2.10} can be made purely algebraic.
    To a 1-motive \(M=(X,A,T,G,u)\) over an algebraically closed field \(k\), we will associated a \emph{Cartier dual} \(M^*=(X',A',T',G',u')\):

    
 \begin{enumerate}
      
 \item[(a)]
        \(X'\) is the character group of \(T\), \(A'\) is the dual abelian variety of \(A\), and \(T'\) is the character group of \(X\).
      


      
 \item[(b)]
        We first consider the case where \(T=0\).
        An extension of \(M\) by \(\mathbb {G}_\mathrm {m}\) then consists of an extension \(E\) of \(A\) by \(\mathbb {G}_\mathrm {m}\), along with a trivialisation \(\widetilde {u}\) of the extension \(u^*E\) of \(X\) by \(\mathbb {G}_\mathrm {m}\).

        
  \begin {tikzcd}
        &&& X
            \ar [dl,swap,"\widetilde {u}"]
            \ar [d,"u"]
        \\0 \ar [r]
        & \mathbb {G}_\mathrm {m} \ar [r]
        & E \ar [r]
        & A \ar [r]
        & 0
        \end {tikzcd}


        We have an exact sequence
        \[           0           \to  \operatorname {Hom}(X,\mathbb {G}_\mathrm {m})           \to  \operatorname {Ext}^1(M,\mathbb {G}_\mathrm {m})           \to  \operatorname {Ext}^1(A,\mathbb {G}_\mathrm {m})           \to  0.         \]

        The extensions of \(M\) by \(\mathbb {G}_\mathrm {m}\) have no non-trivial automorphisms.
        The functor of isomorphism classes of extensions is representable;
        we define \(G'\) as the algebraic group that represents this functor;
        it is an extension of \(A'\) by \(T'\).
        We have a biextension \(P\) of \(M\) and \(G'\) by \(\mathbb {G}_\mathrm {m}\).
      


      
 \item[(c)]
        In the general case, we set \(G'=\operatorname {Ext}^1(M/W_{-2}M,\mathbb {G}_\mathrm {m})\).
        The natural biextension \(P''\) of \(M/W_{-2}M\) and \(G'\) by \(\mathbb {G}_\mathrm {m}\) induces a biextension \(P'\) of \(M\) and \(G'\) by \(\mathbb {G}_\mathrm {m}\).
        For each \(\chi \in  X'\), the extension \(M\) of \(M/W_{-2}M\) by \(T\) defines an extension of \(M/W_{-2}M\) by \(\mathbb {G}_\mathrm {m}\), whence \(u'(\chi )\in  G'\).

        
  \begin {tikzcd}
          0
            \ar [r]
        & T
            \ar [r]
            \ar [d]
        & M
            \ar [r]
            \ar [d,"v"]
        & M/W_{-2}M
            \ar [r]
            \ar [d,equal]
        & 0
        \\0
            \ar [r]
        & \mathbb {G}_\mathrm {m}
            \ar [r]
        & P''_{u(\chi )}
            \ar [r]
        & M/W_{-2}M
            \ar [r]
        & 0
        \end {tikzcd}


        The trivial map \(v\) trivialises the extension \(P'_{u(\chi )}\) of \(M\) by \(\mathbb {G}_\mathrm {m}\)

        
  \begin {tikzcd}
          0
            \ar [r]
        & \mathbb {G}_\mathrm {m}
            \ar [r]
            \ar [d,equal]
        & P_{u(\chi )}
            \ar [r]
            \ar [d]
        & M
            \ar [r]
            \ar [d]
            \ar [dl,swap,"v"]
        & 0
        \\{}
        & \mathbb {G}_\mathrm {m}
            \ar [r]
        & P'_{u(\chi )}
            \ar [r]
        & M/W_{-2}M
        \end {tikzcd}


        These trivialisations make \(P'\) into a biextension of \(M\) and \(M^*=[X'\xrightarrow {u'}G']\) by \(\mathbb {G}_\mathrm {m}\).
        The 1-motive \(M^*\), endowed with \(P'\), is the desired Cartier dual.
      

    \end{enumerate}\label{hodge-theory-iii-10.2.12}\par{}This construction allows us to give a more symmetric definition of 1-motives.\par{}Let \(M=(X,A,T,G,u)\) be a 1-motive, with Cartier dual \(M^*=(X',A',T',G',u')\), and let \(P_0\) be the Poincaré biextension.
    By passing to the quotient, the morphisms \(u\) and \(u'\) define morphisms

    \begin{eqn}{10.2.12.1}\setnumber{10.2.12.1}\label{hodge-theory-iii-10.2.12.1}\skipsinglepar\[         \overline {u}\colon  X\to  A         \quad \text {and}\quad          \overline {u}'\colon  X'\to  A'.       \tag{10.2.12.1}       \]\end{eqn}\par{}The biextension of \(G\) and \(G'\) by \(\mathbb {G}_\mathrm {m}\) induced by \(P_0\) is the inverse image of the Poincaré biextension \(P\) of \(A\) and \(A'\) by \(\mathbb {G}_\mathrm {m}\).
    By hypothesis, the trivialisations (of the inverse images) of \(P\) on \(X\times  G'\) and \(G\times  X'\) agree on \(X\times  X'\), whence

    \begin{eqn}{10.2.12.2}\setnumber{10.2.12.2}\label{hodge-theory-iii-10.2.12.2}\skipsinglepar\[         \text {a trivialisation }\psi \text { of }(\overline {u}\times \overline {u}')^*P.       \tag{10.2.12.2}       \]\end{eqn}\par{}A morphism \(f\colon  M_1\to  M_2\), with transpose (in an evident sense) \(M_2^*\to  M_1^*\), gives rise to commutative diagrams

    \begin{eqn}{10.2.12.3}\setnumber{10.2.12.3}\label{hodge-theory-iii-10.2.12.3}\skipsinglepar\[         \begin {CD}           X_1 @>>> A_1         \\@V{f_X}VV @VV{f_A}V         \\X_2 @>>> A_2         \end {CD}         \qquad \qquad          \begin {CD}           X'_1 @>>> A'_1         \\@A{f'_X}AA @AA{f'_A}A         \\X'_2 @>>> A'_2         \end {CD}       \tag{10.2.12.3}       \]\end{eqn}

    Furthermore, \(f_A\) and \(f'_A\) are transpose to one another, and, via the isomorphism \((1,f'_A)^*P_1\simeq (f_A,1)^*P_2\), we have

    \begin{eqn}{10.2.12.4}\setnumber{10.2.12.4}\label{hodge-theory-iii-10.2.12.4}\skipsinglepar\[         \psi _1(x_1,f'_X(x'_2))         = \psi _2(f_X(x_1),x'_2)       \tag{10.2.12.4}       \]\end{eqn}

    
  \begin {tikzcd}
      X_1\times  X'_1
        \ar [r,"\psi _1"]
    & (P_1\text { on }A_1\times  A'_1)
    \\X_1\times  X'_2
        \ar [u]
        \ar [r]
        \ar [d]
    & (1,f'_A)^*P_1 = (f_A,1)^*P_2\text { on }A_1\times  A'_2
        \ar [u]
        \ar [d]
    \\X_2\times  X'_2
        \ar [r,swap,"\psi _2"]
    & (P_2\text { on }A_2\times  A'_2)
    \end {tikzcd}\label{hodge-theory-iii-10.2.13}\par{}We thus obtain a functor form the category of 1-motives to the following category.
    Its objects consist of:

    
 \begin{enumerate}
      
 \item[(a)]
        two abelian varieties \(A\) and \(A'\), dual to one another;
      

      
 \item[(b)]
        two free \(\mathbb {Z}\)-modules \(X\) and \(X'\), of finite rank;
      

      
 \item[(c)]
        two morphisms \(v\colon  X\to  A\) and \(v'\colon  X'\to  A'\);
      

      
 \item[(d)]
        a trivialisation \(\psi \) of the inverse image under \((v,v')\) of the Poincaré biextension of \(A\) and \(A'\) by \(\mathbb {G}_\mathrm {m}\).
      

    \end{enumerate}\par{}Its arrows consist of diagrams of the form \hyperref[hodge-theory-iii-10.2.12.3]{Equation~\ref*{hodge-theory-iii-10.2.12.3}} that satisfy the conditions of \cref{hodge-theory-iii-10.2.12}\begin{proposition}{10.2.14}\setnumber{10.2.14}\label{hodge-theory-iii-10.2.14}\skipsinglepar\par{}The functor in \cref{hodge-theory-iii-10.2.13} is an equivalence of categories.\begin{proof}\end{proof}\end{proposition}\section{Algebraic interpretation of the mixed \(\operatorname {H}^1\): the case of curves}\label{hodge-theory-iii-10.3}\label{hodge-theory-iii-10.3.1}\par{}Let \(X_0\) be an algebraic curve over an algebraically closed field \(k\).
    Let \(p\colon  X'\to  X_0\) be the normalisation of \(X_0\), and let \(\overline {X}'\) be the non-singular projective curve in which \(X'\) is a dense open.
    We denote by \(X\) (resp. by \(\overline {X}\)) the curve induced from \(X'\) (resp. from \(\overline {X}'\)) by contracting each of the finite sets \(p^{-1}(s)\) (for \(s\in  X_0\)) to a point.

    
  \begin {tikzcd}
      X'
        \ar [r,hook,"j'"]
        \ar [d,swap,"r"]
    & \overline {X'}
        \ar [d,"\overline {r}"]
    \\X
        \ar [r,hook,swap,"j"]
        \ar [d,swap,"p_0"]
    & \overline {X}
    \\X_0
    \end {tikzcd}\par{}We can characterise \(X\) and \(\overline {X}\) by the following properties.

    
 \begin{enumerate}
      
 \item[(a)]
        \(p_0\) if finite and radicial, and thus induces isomorphisms
        \[           \operatorname {H}^\bullet (X_0,\mathbb {Z}_\ell )           \to  \operatorname {H}^\bullet (X,\mathbb {Z}_\ell )         \]
        (with \(\ell \) coprime to the characteristic).
      

      
 \item[(b)]
        The singularities of \(X\) are analytically isomorphic to those presented by a union of coordinate axes in affine space.
      

      
 \item[(c)]
        \(\overline {X}\) is projective and \(\overline {X}\setminus {X}\) is a finite set of points where \(\overline {X}\) is smooth.
      

    \end{enumerate}\label{hodge-theory-iii-10.3.2}\par{}The Picard group \(\operatorname {Pic}(\overline {X})\) \textbf{\emph{TO-DO: unfolds?}} as follows.

    
 \begin{enumerate}
      
 \item[(a)]
        Let \(I\) be the set of irreducible components of \(\overline {X}\) (which is equal to the set of irreducible components of \(X_0\)).
        For each irreducible component \(Y\) of \(\overline {X}\), the degree of the restriction to \(Y\) of an invertible sheaf is a homomorphism
        \[           \begin {aligned}             \operatorname {deg}_Y\colon  \operatorname {Pic}(\overline {X})           & \to  \mathbb {Z}           \\\mathscr {L}           & \mapsto  \operatorname {deg}_Y(\mathscr {L}).           \end {aligned}         \]
        We thus obtain
        \[           \operatorname {deg}\colon \operatorname {Pic}(\overline {X})\to \mathbb {Z}^I.         \]
        This homomorphism is surjective, with kernel \(\operatorname {Pic}^0(\overline {X})\).
      


      
 \item[(b)]
        The morphism from \(\operatorname {Pic}^0(\overline {X})\) to the abelian variety \(\operatorname {Pic}^0(\overline {X}')\) is surjective, and its kernel is a torus.
      

    \end{enumerate}\label{hodge-theory-iii-10.3.3}\par{}Let \(S\) be the finite set \(\overline {X}\setminus  X\), and \(\alpha _0\colon  S\to  I\) be the map that sends \(s\in  S\) to the unique irreducible component of \(\overline {X}\) such that \(s\in \alpha (s)\); let \(\alpha \colon \mathbb {Z}^S\to \mathbb {Z}^I\) be induced by \(\alpha _0\).\par{}Each \(s\in  S\) defines an invertible sheaf \(\mathcal {O}(s)\), which defines \(u_0\colon \mathbb {Z}^S\to \operatorname {Pic}(\overline {X})\).
    We have that \(\operatorname {deg} u_0=\alpha \).\begin{definition}{10.3.4}\setnumber{10.3.4}\label{hodge-theory-iii-10.3.4}\skipsinglepar\par{}The \emph{motivic \(\operatorname {H}^1\) of \(X_0\)}, denoted \(\operatorname {H}_\mathrm {m}^1(X_0)(1)\), is the 1-motive
    \[       \operatorname {Ker}(\alpha )       \to  \operatorname {Pic}^0(\overline {X}).     \]\par{}By definition, we thus have that \(\operatorname {H}_\mathrm {m}^1(X_0)(1)=\operatorname {H}_\mathrm {m}^1(X)(1)\).\end{definition}\begin{remarks}{10.3.5}\setnumber{10.3.5}\label{hodge-theory-iii-10.3.5}\skipsinglepar\begin{enumerate}
    
 \item[(i)]
      We have that
      \[         \begin {aligned}           \operatorname {Gr}_0^W(\operatorname {H}_\mathrm {m}^1(X)(1))           &= \operatorname {Ker}(\alpha \colon \mathbb {Z}^S\to \mathbb {Z}^I)         \\\operatorname {Gr}_{-1}^W(\operatorname {H}_\mathrm {m}^1(X)(1))           &= \operatorname {Pic}^0(\overline {X}').         \end {aligned}       \]
    


    
 \item[(ii)]
      The cokernel of \(\alpha \) is torsion free: it can be identified with the free \(\mathbb {Z}\)-module with basis \(I\setminus \alpha _0(S)\).
      We thus have a short exact sequence of complexes and a quasi-isomorphism
      
  \begin {tikzcd}
        0
          \ar [r]
      & {[\operatorname {Ker}(\alpha )\to \operatorname {Pic}^0(\overline {X}_0)]}
          \ar [r]
      & {[\mathbb {Z}^S\to \operatorname {Pic}(\overline {X}_0)]}
          \ar [r]
      & {[\operatorname {Im}(\alpha )\to \mathbb {Z}^I]}
          \ar [r]
          \ar [d,"\simeq "]
      & 0
      \\&&&{[0\to \operatorname {Coker}(\alpha )]}
      \end {tikzcd}

    

  \end{enumerate}\end{remarks}\begin{construction}{10.3.6}\setnumber{10.3.6}\label{hodge-theory-iii-10.3.6}\skipsinglepar\par{}For \(n>0\) integer and coprime to the characteristic of \(k\), we will construct an isomorphism
    \[       \operatorname {H}^1(X,\mathbb {Z}{/}n\mathbb {Z})(1)       = \operatorname {H}^1(X,\mu _n)       \cong  T_{\mathbb {Z}{/}n\mathbb {Z}}(\operatorname {H}_\mathrm {m}^1(X)(1)).     \]\par{}We know that \(\operatorname {H}^1(X,\mu _n)\) can be identified with the isomorphism classes of invertible sheaves \(\mathscr {L}\) on \(X\), endowed with an isomorphism \(\mathscr {L}^{\otimes  n}\cong \mathcal {O}\).
    Since \(p_0\colon  X_0\to  X\) is finite and radicial, we have that \(\operatorname {H}^1(X,\mathbb {Z}{/}n\mathbb {Z})\xrightarrow {\sim }\operatorname {H}^1(X_0,\mathbb {Z}{/}n\mathbb {Z})\), and \(\operatorname {H}^1(X,\mu _n)\) is again the set of isomorphism classes of pairs \((\mathscr {L},\alpha )\) with \(\mathscr {L}\) an invertible sheaf on \(X_0\) and \(\alpha \) an isomorphism \(\alpha \colon \mathscr {L}^{\otimes {n}}\xrightarrow {\sim }\mathcal {O}\).\par{}Now consider such a pair \((\mathscr {L},\alpha )\).
    The invertible sheaf \(\mathscr {L}\) extends to \(\overline {\mathscr {L}}\) on \(\overline {X}\), and there exists a divisor \(D\) of support \(S\) such that \(\alpha \) extends to \(\overline {\alpha }\colon \overline {\mathscr {L}}^{\otimes  n}\xrightarrow {\sim }\mathcal {O}(D)\).
    If there exists an isomorphism \(\beta \) from \(\overline {\mathscr {L}}^{\otimes  n}\) to \(\mathcal {O}(E)\), then it is uniquely determined up to multiplication by an element of the divisible group \(\operatorname {H}^0(\overline {X},\mathcal {O}^*)\).
    We thus deduce that \((\overline {\mathscr {L}},D)\) determines \((\mathscr {L},\alpha )\) up to isomorphism.
    For a pair \((\overline {\mathscr {L}},D)\) to come from a suitable pair \((\mathscr {L},\alpha )\), it is necessary and sufficient that \(n[\mathscr {L}]=[\mathcal {O}(D)]\) in \(\operatorname {Pic}(\overline {X})\); it comes from \((\mathcal {O}_{X_0},0)\) if and only if it is of the form \((\mathcal {O}_{X_0}(D),nD)\).\par{}This identifies \(\operatorname {H}^1(X,\mu _n)\) with the \(\operatorname {H}^0\) of the tensor product of \([\mathbb {Z}^S\to \operatorname {Pic}(\overline {X})]\) (concentrated in degrees \(0\) and \(1\)) with \([\mathbb {Z}\xrightarrow {n}\mathbb {Z}]\) (concentrated in degrees \(-1\) and \(0\)):
    \[       \begin {CD}         \mathbb {Z}^S @>{u_0}>> \operatorname {Pic}(\overline {X})       \\@A{n}AA @AA{-n}A       \\\mathbb {Z}^S @>>{u_0}> \operatorname {Pic}(\overline {X})       \end {CD}     \]
    \[       \operatorname {H}^1(X,\mu _n)       \cong  \operatorname {H}^0\big (         [\mathbb {Z}^S\to \operatorname {Pic}(\overline {X})] \overset {\mathrm {L}}{\otimes } \mathbb {Z}{/}n\mathbb {Z}.       \big )     \]
    Since \(\operatorname {Coker}(\alpha )\) is torsion free, we read from (ii) of \hyperref[hodge-theory-iii-10.3.5]{Remarks~\ref*{hodge-theory-iii-10.3.5}} that
    \[       \begin {aligned}         \operatorname {H}^1(X,\mu _n)         &\cong  \operatorname {H}^0\big (           [\mathbb {Z}^S\to \operatorname {Pic}(\overline {X})] \overset {\mathrm {L}}{\otimes } \mathbb {Z}{/}n\mathbb {Z}.         \big )       \\&\xleftarrow {\sim } \operatorname {H}^0\big (           \operatorname {H}_\mathrm {m}^1(X) \overset {\mathrm {L}}{\otimes } \mathbb {Z}{/}n\mathbb {Z}         \big )       \\&\cong  T_{\mathbb {Z}{/}n\mathbb {Z}}(\operatorname {H}_\mathrm {m}^1(X)(1)).       \end {aligned}     \]\end{construction}\begin{variant}{10.3.7}\setnumber{10.3.7}\label{hodge-theory-iii-10.3.7}\skipsinglepar\par{}For \(\ell \) coprime to the characteristic of \(k\), we have that
    \[       \operatorname {H}^1(X_0,\mathbb {Z}_\ell (1))       \simeq  T_\ell (\operatorname {H}_\mathrm {m}^1(X_0)(1)).     \]\end{variant}\begin{construction}{10.3.8}\setnumber{10.3.8}\label{hodge-theory-iii-10.3.8}\skipsinglepar\par{}Let \(X_0\) be a curve on \(\mathbb {C}\).
    We will construct an isomorphism of mixed Hodge structures
    \[       \operatorname {H}^1(X_0)(1)       \simeq  T(\operatorname {H}_\mathrm {m}^1(X_0)(1)).     \]\begin{proof}\(X=X_0\)\end{proof}\end{construction}\begin{construction}{10.3.9}\setnumber{10.3.9}\label{hodge-theory-iii-10.3.9}\skipsinglepar\begin{enumerate}
    
 \item[(i)]
      Let \(j_*^\mathrm {m}\mathcal {O}_X^*\) be the subsheaf of meromorphic functions in \(j_*\mathcal {O}_X^*\).
      We have that
      \[         \operatorname {H}^1(X,\mathbb {Z}(1))         = \operatorname {\mathbb {H}}^1(\overline {X},[\mathcal {O}_{\overline {X}}\xrightarrow {\exp }j_*^\mathrm {m}\mathcal {O}_X^*]).       \]
    


    
 \item[(ii)]
      We have that
      \[         \operatorname {H}^1(X,\mathbb {C})         = \operatorname {\mathbb {H}}^1(\overline {X},[\mathcal {O}_{\overline {X}}\xrightarrow {\mathrm {d}}r_*\Omega _{\overline {X}'}^1(\log  S)]).       \]
    


    
 \item[(iii)]
      The inclusion of \(\operatorname {H}^1(X,\mathbb {Z})\) into \(\operatorname {H}^1(X,\mathbb {C})\) is defined by the morphism of complexes
      \[         \begin {CD}           \mathcal {O}_{\overline {X}} @>{\exp }>> j_*^\mathrm {m}\mathcal {O}_X^*         \\@V{\simeq }VV @VV{\mathrm {d} f/f}V         \\\mathcal {O}_{\overline {X}} @>>> r_*\Omega _{\overline {X}'}^1(\log  S)         \end {CD}       \]
    

  \end{enumerate}\begin{proof}\par{}Consider the commutative diagram
      
  \begin {tikzcd}[row sep=huge]
        \operatorname {H}^1(X,\mathbb {Z}(1))
          \ar [r,"\sim "]
          \ar [d]
          \ar [dr,phantom,"\text {(1)}"]
      & \operatorname {\mathbb {H}}^1(X, [\mathcal {O}_X\xrightarrow {\exp }\mathcal {O}_X^*])
          \ar [d]
          \ar [dr,phantom,"\text {(2)}"]
      & \operatorname {\mathbb {H}}^1(\overline {X}, [j_*\mathcal {O}_X\xrightarrow {\exp }j_*\mathcal {O}_X^*])
          \ar [l,swap,"\sim "]
          \ar [d]
          \ar [dr,phantom,"\text {(3)}"]
      & \operatorname {\mathbb {H}}^1(\overline {X}, [\mathcal {O}_{\overline {X}}\xrightarrow {\exp }j_*^\mathrm {m}\mathcal {O}_X^*])
          \ar [l,swap,"\sim "]
          \ar [d]
      \\\operatorname {H}^1(X,\mathbb {C})
          \ar [r]
      & \operatorname {\mathbb {H}}^1(X, [\mathcal {O}_X\xrightarrow {\mathrm {d}}r_*\Omega _{\overline {X}'}^1])
      & \operatorname {\mathbb {H}}^1(\overline {X}, [j_*\mathcal {O}_X\to  j_*r_*\Omega _{X'}^1])
          \ar [l]
      & \operatorname {\mathbb {H}}^1(\overline {X}, [\mathcal {O}_{\overline {X}}\to  r_*\Omega _{\overline {X}'}^1(\log  S)])
          \ar [l,swap,"\sim "]
      \end {tikzcd}\par{}The horizontal arrows in the square labelled (1) are isomorphisms, since the sequences
      \[         \begin {gathered}           0 \to  \mathbb {Z}(1) \to  \mathcal {O}_X \xrightarrow {\exp } \mathcal {O}_X^* \to  0         \\0 \to  \mathbb {C} \to  \mathcal {O}_X \xrightarrow {\mathrm {d}} r_*\Omega _{X'}^1 \to  0         \end {gathered}       \]
      are exact.
      Those in (2) are isomorphisms since \(\mathrm {R}^1j_*\mathcal {O}_X=0\).
      Those in (3) are isomorphisms since the morphisms of the complexes that define them are quasi-isomorphisms.
      This diagram defines the desired construction.\end{proof}\end{construction}\label{hodge-theory-iii-10.3.10}\par{}Recall that if \(d\colon  F\to  G\) is a morphism of abelian sheaves on a space \(Z\), then \(\operatorname {\mathbb {H}}^1(Z,[F\to  G])\) can be identified with the set of isomorphism classes of pairs \((P,\alpha )\), where \(P\) is an \(F\)-torsor (i.e. a principal homogeneous space under \(F\)) and \(\alpha \) is a trivialisation of the \(G\)-torsor \(dP\).\par{}In particular:

    
 \begin{enumerate}
      
 \item[(a)]
        \par{}\(\operatorname {\mathbb {H}}^1(\overline {X}, [\mathcal {O}_{\overline {X}}^*\to  r_*\Omega _{\overline {X}'}^1(\log  S)])\) can be identified with the set of isomorphism classes of pairs \((\mathscr {L},\alpha )\), where \(\mathscr {L}\) is an invertible sheaf on \(\overline {X}\) and \(\alpha \) is a connection on \(\overline {r}^*\mathscr {L}\) that is holomorphic on \(X'\) and can present a simple pole at the points of \(S\).
          Let \(a\) be the projection from \(\overline {X}\) to \(\operatorname {Spec}(\mathbb {C})\) (i.e. a point).
          In the fppf topology, the sheaf \(\mathrm {R}^1a_*([\mathscr {O}_{\overline {X}}^*\to  r_*\Omega _{\overline {X}'}^1(\log  S)])\) is representable by a group scheme \(\operatorname {Pic}^\natural (X)\), of which \(\operatorname {\mathbb {H}}^1(\overline {X},[\mathscr {O}_{\overline {X}}^*\to  r_*\Omega _{\overline {X}'}^1(\log  S)])\) is the set of points.
          \(\operatorname {Pic}^\natural (X)\) is an extension of a subgroup of \(\operatorname {Pic}(\overline {X})\), containing \(\operatorname {Pic}^0(\overline {X})\), by \(\operatorname {H}^0(\overline {X}',\Omega _{\overline {X}'}^1(\log  S))\).

        \par{}We define a map

          \begin{eqn}{10.3.10.1}\setnumber{10.3.10.1}\label{hodge-theory-iii-10.3.10.1}\skipsinglepar\[               \mathbb {Z}^S               \to  \operatorname {Pic}^\natural (X)             \tag{10.3.10.1}             \]\end{eqn}

          by sending each divisor \(D\) of \(\overline {X}\), concentrated in \(S\), to the invertible sheaf \(\mathcal {O}(D)\).
      


      
 \item[(b)]
        \(\operatorname {H}^1(X,\mathbb {C})=\operatorname {\mathbb {H}}^1(\overline {X}, [\mathcal {O}_{\overline {X}}\to  r_*\Omega _{\overline {X}'}^1(\log  S)])\) is the Lie algebra of \(\operatorname {Pic}^\natural (X)\).
        It is the set of isomorphism classes of pairs \((P,\alpha )\), where \(P\) is an \(\mathcal {O}_{\overline {X}}\)-torsor and \(\alpha \) is a connection on \(P\) as in (a).
      


      
 \item[(c)]
        \par{}\(\operatorname {H}^1(X,\mathbb {Z}(1))=\operatorname {\mathbb {H}}^1(\overline {X}, [\mathcal {O}_{\overline {X}}\to  j_*^\mathrm {m}\mathcal {O}_{\overline {X}}^*])\) is the set of isomorphism classes of pairs \((P,\alpha )\), where \(P\) is an \(\mathcal {O}_{\overline {X}}\)-torsor and \(\alpha \) is an isomorphism from the invertible sheaf \(\exp (P)\) to an invertible sheaf \(\mathcal {O}_X(D)\) for some \(D\) concentrated in \(S\).
          The map
          \[             \operatorname {Aut}(P)             = \mathbb {C}             \to  \operatorname {Aut}(\exp (P))             = \mathbb {C}^*           \]
          is surjective.
          This allows us to also identify \(\operatorname {H}^1(X,\mathbb {Z}(1))\) with the set of pairs \((p,d)\), where \(p\) is an isomorphism class of \(\mathcal {O}_{\overline {X}}\)-torsors (i.e. an element of \(\operatorname {Lie}(\operatorname {Pic}\overline {X})\)) and \(d\in \operatorname {Ker}(\alpha )\), defining a divisor concentrated in \(S\) with whose image in \(\operatorname {Pic}^0(\overline {X})\) is \(\exp (p)\).
          In other words, we have defined an isomorphism

          \begin{eqn}{10.3.10.2}\setnumber{10.3.10.2}\label{hodge-theory-iii-10.3.10.2}\skipsinglepar\[               \operatorname {H}^1(X,\mathbb {Z}(1))               \xrightarrow {\sim } T_\mathbb {Z}(\operatorname {H}_\mathrm {m}^1(X)(1)).             \tag{10.3.10.2}             \]\end{eqn}

        \par{}From part (i) of \hyperref[hodge-theory-iii-10.3.11]{Lemma~\ref*{hodge-theory-iii-10.3.11}}, we deduce that this isomorphism is compatible with the weight filtration.
      

    \end{enumerate}\begin{lemma}{10.3.11}\setnumber{10.3.11}\label{hodge-theory-iii-10.3.11}\skipsinglepar\begin{enumerate}
    
 \item[(i)]
      
      We have that
      \[         \begin {aligned}           W^1(\operatorname {H}^1(X,\mathbb {Z}))           &= \operatorname {Im}\big ( \operatorname {H}^1(\overline {X},\mathbb {Z})\to \operatorname {H}^1(X,\mathbb {Z}) \big )         \\W^0(\operatorname {H}^1(X,\mathbb {Z}))           &= \operatorname {Ker}\big ( \operatorname {H}^1(\overline {X},\mathbb {Z})\to \operatorname {H}^1(\overline {X}',\mathbb {Z}) \big ).         \end {aligned}       \]
    


    
 \item[(ii)]
      The spectral sequence defined by the stupid filtration of \([\mathcal {O}_{\overline {X}}\to \overline {r}_*\Omega _{\overline {X}'}^1(\log  S)]\) degenerates and abuts to the Hodge filtration of \(\operatorname {H}^*(X,\mathbb {C})\).
    

  \end{enumerate}\begin{proof}\par{}We first prove (i).
      It suffices to prove the first claim: the second follows from \hyperref[hodge-theory-iii-8.2.5]{Proposition~\ref*{hodge-theory-iii-8.2.5}}.
      We have a comparison of exact sequences
      \[         \begin {CD}           \operatorname {H}^1(\overline {X},\,\mathbb {Z}) @>>> \operatorname {H}^1(X,\,\mathbb {Z}) @>{\partial }>> \operatorname {H}^2(\overline {X}\,\operatorname {mod}\, X,\,\mathbb {Z})         \\@VVV @VVV @VV{\simeq }V         \\\operatorname {H}^1(\overline {X}',\,\mathbb {Z}) @>>> \operatorname {H}^1(X',\,\mathbb {Z}) @>{\partial }>> \operatorname {H}^2(\overline {X}'\,\operatorname {mod}\, X',\,\mathbb {Z}) @>>> \operatorname {H}^2(\overline {X}',\,\mathbb {Z})         \end {CD}       \]
      It follows from \hyperref[hodge-theory-iii-8.3.7]{Example~\ref*{hodge-theory-iii-8.3.7}} and \hyperref[hodge-theory-ii-3.2.17]{Corollary~\ref*{hodge-theory-ii-3.2.17}} that \(\operatorname {H}^2(\overline {X}'\operatorname {mod} X,\mathbb {Z})\), and thus \(\operatorname {H}^2(\overline {X}\operatorname {mod} X,\mathbb {Z})\), is pure of weight 2.
      Taking \hyperref[hodge-theory-iii-8.2.4]{Theorem~\ref*{hodge-theory-iii-8.2.4}}, this proves (i).\par{}To prove (ii), we have to return to \cref{hodge-theory-iii-8.2.1} and use the fact that the simplicial scheme \(((X'/X)^{\Delta _n})_{n\geqslant 0}\) is smooth (the \((n+1)\)-tuple \((X'/X)^{\Delta _n}\) given by the iterated fibre product is the disjoint sum of \(X'\), the diagonal, and a finite number of points) and that it admits \(((\overline {X}'/\overline {X})^{\Delta _n})_{n\geqslant 0}\) as a smooth compactification.
      Let \(\varepsilon \colon ((\overline {X}'/\overline {X})^{\Delta _n})_{n\geqslant 0}\to \overline {X}\) the augmentation morphism.
      Let \((K,F)\) be the complex \(\mathcal {O}\to \Omega ^1(\log  S)\) on \(((\overline {X}'/\overline {X})^{\Delta _n})_{n\geqslant 0}\) endowed with the stupid filtration.
      We have an augmentation morphism
      \[         \alpha \colon          [\mathcal {O}_{\overline {X}}\to \overline {r}_*\Omega _{\overline {X}'}^1(\log  S)]         \to  \mathrm {s}{\varepsilon _\bullet }_*K.       \]
      The diagram
      
  \begin {tikzcd}
        \operatorname {H}^1(X,\mathbb {C})
          \ar [r,equals]
          \ar [d,equals]
      & \operatorname {\mathbb {H}}^1\big ( ((\overline {X}'/\overline {X})^{\Delta _n})_\bullet , K \big )
      & \operatorname {\mathbb {H}}^1(\overline {X},\mathrm {s}{\varepsilon _\bullet }_*K)
          \ar [l,swap,"\text {(1)}"]
      \\\operatorname {H}^1(X,\mathbb {C})
          \ar [rr,equals]
      &&\operatorname {\mathbb {H}}^1\big ( \overline {X}, [\mathcal {O}_{\overline {X}}\to \overline {r}_*\Omega _{\overline {X}'}^1*\log  S] \big )
          \ar [u]
      \end {tikzcd}

      commutes, and the arrow labelled (1) is an isomorphism since the \(\varepsilon _n\) are finite, so \(\mathrm {R}^i{\varepsilon _n}_*=0\) for \(i>0\).
      The claim then follows from the fact that \(\alpha \) is a filtered quasi-isomorphism.\end{proof}\end{lemma}\label{hodge-theory-iii-10.3.12}\par{}We now realise \hyperref[hodge-theory-iii-10.3.8]{Construction~\ref*{hodge-theory-iii-10.3.8}}.
    The diagram
    \[       \begin {CD}         \operatorname {H}^1(X,\mathbb {Z}(1)) @>{\simeq }>> T_\mathbb {Z}(\operatorname {H}_\mathrm {m}^1(X)(1)) @>>> \operatorname {Lie}\operatorname {Pic}(\overline {X}) = \operatorname {H}^1(\overline {X},\mathcal {O}_{\overline {X}})       \\ @| @. @|       \\\operatorname {H}^1(X,\mathbb {C}) @>{\simeq }>> \operatorname {Lie}\operatorname {Pic}^\natural (X) @>>> \operatorname {Lie}\operatorname {Pic}(\overline {X}) = \operatorname {H}^1(\overline {X},\mathcal {O}_{\overline {X}})       \end {CD}     \]
    commutes.
    By (ii) of \hyperref[hodge-theory-iii-10.3.11]{Lemma~\ref*{hodge-theory-iii-10.3.11}} and the definition of Hodge filtrations, \hyperref[hodge-theory-iii-10.3.10.2]{Equation~\ref*{hodge-theory-iii-10.3.10.2}} is compatible with the Hodge filtration.\begin{corollary}{10.3.13}\setnumber{10.3.13}\label{hodge-theory-iii-10.3.13}\skipsinglepar\par{}\(\operatorname {H}_\mathrm {m}^1(X)(1)^\natural \) is the extension
    \[       \operatorname {Ker}(\alpha )       \to  \operatorname {Pic}^\natural (X)^0     \]
    of \(\operatorname {H}_\mathrm {m}^1(X)(1)\) induced by \hyperref[hodge-theory-iii-10.3.10.1]{Equation~\ref*{hodge-theory-iii-10.3.10.1}}.\begin{proof}\[       \begin {CD}         \operatorname {Ker}(\alpha ) @>>> \operatorname {Pic}^\natural (X)^0 @>>> \operatorname {Pic}^0(\overline {X})       \\@AAA @AAA @AAA       \\\operatorname {H}^1(X,\mathbb {Z}) @>>> \operatorname {Lie}\operatorname {Pic}^\natural (X) @>>> \operatorname {Lie}\operatorname {Pic}(\overline {X})       \end {CD}     \qedhere\]\end{proof}\end{corollary}\begin{exercise}{10.3.14}\setnumber{10.3.14}\label{hodge-theory-iii-10.3.14}\skipsinglepar\par{}Show that the diagram
    \[       \begin {CD}         \operatorname {H}^1(X,\mathbb {Z}(1)) @= T_\mathbb {Z}(\operatorname {H}_\mathrm {m}^1(X)(1))       \\@VVV @VVV       \\\operatorname {H}^1(X,\mu _n) @= T_{\mathbb {Z}{/}n\mathbb {Z}}(\operatorname {H}_\mathrm {m}^1(X)(1))       \end {CD}     \]
    commutes.\end{exercise}\begin{remark}{10.3.15}\setnumber{10.3.15}\label{hodge-theory-iii-10.3.15}\skipsinglepar\par{}The isomorphism induced by \hyperref[hodge-theory-iii-10.3.8]{Construction~\ref*{hodge-theory-iii-10.3.8}} and (ii) of \hyperref[hodge-theory-iii-10.3.9]{Construction~\ref*{hodge-theory-iii-10.3.9}}, between \(\operatorname {\mathbb {H}}(\overline {X},\mathcal {O}_{\overline {X}}\to  r_*\Omega _{\overline {X}'}^1(\log  S))\) (which we may call \(\operatorname {H}_\mathrm {DR}^1(X)\)) and \(T_\mathrm {DR}(\operatorname {H}_\mathrm {m}^1(X)(1))\), equal to \(\operatorname {Lie}\operatorname {Pic}^\natural (X)\) by \hyperref[hodge-theory-iii-10.3.13]{Corollary~\ref*{hodge-theory-iii-10.3.13}}, has been defined in a purely algebraic manner.\end{remark}\section{Translation of a theorem of Picard}\label{hodge-theory-iii-10.4}\label{hodge-theory-iii-10.4.1}\par{}Let \(\mathscr {H}\) be a mixed Hodge structure such that \(h^{pq}=0\) for \(p<0\) or \(q<0\) (resp. for \(p>n\) or \(q>n\)).
    I provisionally denote by \(\operatorname {I}(\mathscr {H})\) (resp. by \(\operatorname {II}_n(\mathscr {H})\)) the 1-motive defined by the largest mixed Hodge sub-structure of \((\mathscr {H}_\mathbb {Z}{/}\mathrm {torsion})(1)\) (resp. by the largest quotient of \((\mathscr {H}_\mathbb {Z}{/}\mathrm {torsion})(n)\)) that is of pure of type
    \(       \big \{ (-1,-1), (-1,0), (0,-1), (0,0) \big \}.     \)
    If \(X\) in a complex algebraic variety of dimension \(\leqslant  N\), I conjecture that the 1-motives \(\operatorname {I}(\operatorname {H}^n(X,\mathbb {Z}))\), \(\operatorname {II}_n(\operatorname {H}^n(X,\mathbb {Z}))\) (for \(n\leqslant  N\)), and \(\operatorname {II}_n(\operatorname {H}^n(X,\mathbb {Z}))\) (for \(n\geqslant  N\)) admit a purely algebraic definition.
    The morphisms
    \[       \begin {aligned}         T_\ell (\operatorname {I}(\operatorname {H}^n(X,\mathbb {Z})))         & \to  (\operatorname {H}^n(X,\mathbb {Z}_\ell )/\mathrm {torsion})(1)       \\T_\ell (\operatorname {II}_n(\operatorname {H}^n(X,\mathbb {Z})))         & \leftarrow  \operatorname {H}^n(X,\mathbb {Z}_\ell )(n)           \quad \text {for }n\leqslant  N       \\T_\ell (\operatorname {II}_N(\operatorname {H}^n(X,\mathbb {Z})))         & \leftarrow  \operatorname {H}^n(X,\mathbb {Z}_\ell )(N)           \quad \text {for }n\geqslant  N       \end {aligned}     \]
    (and their analogues in de Rham cohomology) should also admit a purely algebraic definition.
    This is what we have shown for \(\dim (X)=1\).
    Here is another example.\label{hodge-theory-iii-10.4.2}\par{}Let \(S\) be a smooth projective surface, \(U\) a dense Zariski open of \(S\), and \(C=S\setminus  U\).
    Set
    \[       \operatorname {H}_2^{\log }(U,\mathbb {Z})       = \operatorname {Ker}\big ( \operatorname {H}_2(U,\mathbb {Z})\to \operatorname {H}_2(S,\mathbb {Z}) \big ).     \]
    By [\hyperref[G1968]{G1968}, §9], this group does not depend on the chosen smooth compactification \(S\) of \(U\).
    Modulo finite groups, this is nothing more (\hyperref[hodge-theory-ii-3.2.17]{Corollary~\ref*{hodge-theory-ii-3.2.17}}) than \(W_{-3}\operatorname {H}_2(U,\mathbb {Z})\) (\(\operatorname {H}_2(U,\mathbb {Q})\), the dual of \(\operatorname {H}^2(U,\mathbb {Q})\), is endowed with the dual mixed Hodge structure).
    We can show that the 1-motive defined by \((\operatorname {H}_2^{\log }(U,\mathbb {Z})/\mathrm {torsion})(-1)\) admits the following description:

    
 \begin{enumerate}
      
 \item[(a)]
        By Poincaré duality, we have that \(\operatorname {H}_2(U,\mathbb {Z})=\operatorname {H}_\mathrm {c}^2(U,\mathbb {Z})\).
        The long exact sequence of cohomology of \((S,C)\) gives an exact sequence
        \[           \operatorname {H}^1(S,\mathbb {Z})           \to  \operatorname {H}^1(C,\mathbb {Z})           \to  \operatorname {H}_2^{\log }(U,\mathbb {Z})           \to  0.         \]
      


      
 \item[(b)]
        Let \(\operatorname {H}_\mathrm {m}^1(C)\) be the quotient of \(\operatorname {Pic}^0(C)\) by its unipotent radical.
        We know that \(\operatorname {H}^1(C,\mathbb {Z})=T_\mathbb {Z}(\operatorname {H}_\mathrm {m}^1(C))\) and \(\operatorname {H}^1(S,\mathbb {Z})=T_\mathbb {Z}(\operatorname {Pic}^0(S))\).
        We thus have a map
        \[           \operatorname {H}_2^{\log }(U,\mathbb {Z})           \to  T_\mathbb {Z}(H_\mathrm {m}^1(C)/\operatorname {Pic}^0(S))         \]
        which is an isomorphism modulo finite groups.
      


      
 \item[(c)]
        The map

        \begin{eqn}{10.4.2.1}\setnumber{10.4.2.1}\label{hodge-theory-iii-10.4.2.1}\skipsinglepar\[             (\operatorname {H}_2^{\log }(U,\mathbb {Z})/\mathrm {torsion})(-1)             \to  T_\mathbb {Z}(\operatorname {H}_\mathrm {m}^1(C)/\operatorname {Pic}^0(S))           \tag{10.4.2.1}           \]\end{eqn}

        induced from (b) is a morphism of mixed Hodge structures (and an isomorphism modulo finite groups).
      

    \end{enumerate}\par{}We can define the isomorphisms (modulo finite groups) induced by \hyperref[hodge-theory-iii-10.4.2.1]{Equation~\ref*{hodge-theory-iii-10.4.2.1}} between the analogues of \(\operatorname {H}_2^{\log }(U,\mathbb {Z})(-1)\) in \(\ell \)-adic or de Rham cohomology, as well as the \(\ell \)-adic or de Rham realisations of \(\operatorname {H}_\mathrm {m}^1(C)/\operatorname {Pic}^0(S)\), in a purely algebraic way.\label{hodge-theory-iii-10.4.3}\par{}We will not prove the above constructions and compatibilities.
    The seed of the idea is in the following theorem of E. Picard ([\hyperref[PS1895]{PS1895}, p. 62]).
    Here \(P\) and \(A\) denote polynomials.\begin{theorem}{10.4.3.1}\setnumber{10.4.3.1}\label{hodge-theory-iii-10.4.3.1}\skipsinglepar\par{}Every residue of the double integral
      \[         \iint \frac {P(x,y)\mathrm {d} x\mathrm {d} y}{A(x,y)}       \]
      can be thought of as a logarithmic or cyclic period of an abelian integral.\end{theorem}\par{}Recall that:

    
 \begin{enumerate}
      
 \item[(a)]
        A \emph{\(k\)-fold integral} on a non-singular projective variety \(V\) of dimension \(d\) is a closed rational \(k\)-form on \(V\), i.e. it is an element \(\alpha \in \operatorname {H}^0(U,\Omega ^k)\), where \(U\) is a dense open of \(V\), such that \(\mathrm {d}\alpha =0\).
        Here, \(k=2\) and \(V=\mathbb {P}^2\), and so \(k=\dim  V\) and the condition \(\mathrm {d}\alpha =0\) is trivial.
      


      
 \item[(b)]
        A \(k\)-fold integral \(\alpha \) defines an element, again denoted by \(\alpha \), of \(\operatorname {H}^k(U,\mathbb {C})\).
        The \emph{periods} of \(\alpha \) are the numbers \(\langle \alpha ,c\rangle \) for \(c\in \operatorname {H}_k(U,\mathbb {Z})\).
        If \(c\) is defined by a singular chain \(C\), then
        \[           \langle \alpha ,c\rangle            = \int _C\alpha .         \]
      


      
 \item[(c)]
        The \emph{residues} of a \(k\)-fold integral \(\alpha \in \operatorname {H}^0(U,\Omega ^k)\) are the numbers \(\langle \alpha ,c\rangle \) for
        \[           c\in \operatorname {Ker}\big ( \operatorname {H}_k(U,\mathbb {Z})\to \operatorname {H}_k(V,\mathbb {Z}) \big ).         \]
        In the case considered by Picard, the map \(\operatorname {H}_2(U,\mathbb {Z})\to \operatorname {H}_2(\mathbb {P}^3,\mathbb {Z})\) is zero, and there is no need to distinguish between periods and residues.
        For the general case, see [\hyperref[L1924]{L1924}, note I].
      


      
 \item[(d)]
        To Picard, "logarithmic or cyclic period of an abelian integral" simply means "period of a simple integral \(\beta \) on a curve \(Y\)".
        It is implicit that \(\beta \) and \(Y\) are rationally determined by \(P\) and \(A\).
      

    \end{enumerate}\label{hodge-theory-iii-10.4.4}\par{}Let us translate this theorem.
    Let \(S\) and \(U\) be as in \cref{hodge-theory-iii-10.4.2} (for example, \(S=\mathbb {P}^2(\mathbb {C})\)), and let \(\beta \in \operatorname {H}^0(U,\Omega _U^2)\) be a double integral and \(c\in \operatorname {H}_2^{\log }(U,\mathbb {Z})\) a cycle.
    By definition,
    \[       \langle \beta ,c\rangle        = \int _c\beta      \]
    is a residue of \(\beta \).
    Let \(A=\operatorname {H}_\mathrm {m}^1(C)(1)/\operatorname {Pic}^0(S)\) (an extension of an abelian variety by a torus).
    The 2-form \(\beta \) defines a cohomology class \(\operatorname {cl}(\beta )\in \operatorname {H}_\mathrm {DR}^2(U)\) and induces a linear form \(\beta '\) on \(\operatorname {H}_{\mathrm {DR},2}^{\log }(U)=T_\mathrm {DR}(A)\).
    
    The map \(\beta \mapsto \beta '\) is rational.
    The class \(c\) defines \(c'\in  T_\mathbb {Z}(A)\otimes \mathbb {Q}\), and we have that

    \begin{eqn}{10.4.4.1}\setnumber{10.4.4.1}\label{hodge-theory-iii-10.4.4.1}\skipsinglepar\[         \langle \beta ,c\rangle          = 2\pi  i\langle \beta ',c'\rangle .       \tag{10.4.4.1}       \]\end{eqn}

    We can represent \(c'\) by a cycle on \(C\), and \(\beta '\) as a simple integral on \(C\).